% V : produire une lettre officielle ( lettre de motivation) — Français 1re Tech — Cours signé OnBuch+
% Programme officiel MINESEC. Compiler avec : tectonic N21-v-produire-lettre-officielle-lettre-motivation.tex
\newcommand{\DOCMATIERE}{Français}
\newcommand{\DOCNIVEAU}{1re Tech}
\newcommand{\DOCMODULE}{Français – classe de Première technique}
\newcommand{\DOCLECON}{N21}
\newcommand{\DOCTITRE}{V : produire une lettre officielle ( lettre de motivation)}
% OnBuch+ — préambule « cours signés » ENRICHI (compilé avec tectonic / XeLaTeX)
% Système visuel « Pop » d'OnBuch+ : fond crème (jamais blanc), bordures encre
% épaisses, ombres « dures » décalées, titres Archivo Black en capitales.
% Version MATHÉMATIQUES : graphes (pgfplots/tikz), boîtes pédagogiques
% (définition, propriété, méthode, attention, le savais-tu, exercice résolu,
% exercice, TP, bilan), page de garde et signature OnBuch+ ancrée en bas.
%
% Macros attendues dans main.tex AVANT \input{preamble} :
%   \DOCMATIERE  \DOCNIVEAU  \DOCMODULE  \DOCLECON (numéro)  \DOCTITRE
\documentclass[11pt]{article}

\usepackage{fontspec}
\usepackage[french]{babel}
\usepackage[a4paper,margin=2cm,top=3.4cm,bottom=2.8cm,headheight=1.4cm,footskip=1.3cm]{geometry}
\usepackage{amsmath,amssymb,amsfonts}
\usepackage{mathtools} % \xmapsto, \coloneqq... souvent utilisés par les modèles
\usepackage{cancel} % simplifications barrées (bilans de forces)
% \abs{z} (module, valeur absolue) et \norm{u} (norme), utilisés par les modèles
\providecommand{\abs}{}\let\abs\relax\DeclarePairedDelimiter{\abs}{\lvert}{\rvert}
\providecommand{\norm}{}\let\norm\relax\DeclarePairedDelimiter{\norm}{\lVert}{\rVert}
\usepackage[table]{xcolor} % [table] : fournit aussi \rowcolors (alternance de lignes)
\usepackage{pagecolor}
\usepackage{tikz}
\usetikzlibrary{arrows.meta,positioning,calc,shapes.geometric,shapes.misc,shapes.arrows,decorations.pathmorphing,decorations.pathreplacing,decorations.markings,patterns,shadows,mindmap,trees,fit,backgrounds,matrix,intersections,angles,quotes}
\usepackage{pgfplots}
\usepackage[version=4]{mhchem} % équations nucléaires \ce{^{235}_{92}U}
\usepackage[european,siunitx]{circuitikz} % schémas électriques
\pgfplotsset{compat=1.18}
\usepgfplotslibrary{fillbetween,groupplots} % aires entre courbes (calcul intégral)
\usepackage{enumitem}
\usepackage{fancyhdr}
% [most] charge toutes les bibliothèques tcolorbox (listings, minted, raster,
% documentation...) et gonfle inutilement la mémoire TeX sur un document long
% (~300 boîtes) : on ne charge que ce qui est réellement utilisé.
\usepackage[skins,breakable]{tcolorbox}
\usepackage{graphicx}
\usepackage{colortbl}
\usepackage{booktabs}
\usepackage{tabularx}
\usepackage{multirow}
\usepackage{diagbox} % case d'en-tête diagonale des tableaux à double entrée
% Colonnes extensibles centrée / gauche / droite (C, L, R), souvent utilisées par les modèles
\newcolumntype{C}{>{\centering\arraybackslash}X}
\newcolumntype{L}{>{\raggedright\arraybackslash}X}
\newcolumntype{R}{>{\raggedleft\arraybackslash}X}
\usepackage{array}
\usepackage{multicol}
\usepackage{siunitx}
% parse-numbers=false : accepte aussi \SI{2,5 \times 10^{-3}}{...}
\sisetup{locale=FR,output-decimal-marker={,},group-minimum-digits=4,parse-numbers=false}
\DeclareSIUnit\ohm{\ensuremath{\symup{\Omega}}} % U+2126 absent de la police
\DeclareSIUnit\division{div} % réglages d'oscilloscope (ms/div, V/div)
\DeclareSIUnit\tr{tr} % tours (vitesse de rotation, ex. \SI{1500}{\tr\per\minute})
\DeclareSIUnit\molar{M} % concentration molaire (ex. \SI{3}{\molar}, \SI{1}{\micro\molar})
\DeclareSIUnit\an{an} % année (le modèle confond parfois avec \year, un compteur TeX) — ex. \SI{5}{\kilo\meter\per\an}
\DeclareSIUnit\FCFA{FCFA} % franc CFA (ex. \SI{500}{\FCFA\per\kilo\gram})
\DeclareSIUnit\degreeBrix{\textdegree Brix} % teneur en sucre d'un sirop/jus (réfractométrie)
\DeclareSIUnit\month{mois} % le modèle utilise parfois \month, un compteur TeX, comme unité de durée
\DeclareSIUnit\microliter{\micro\liter} % le modèle écrit parfois \microliter en un seul token
\DeclareSIUnit\million{millions} % dénombrement (ex. \SI{15}{\million\per\mL} spermatozoïdes)
\usepackage{qrcode}
% \degree : commande fréquemment utilisée par les modèles pour un angle ou une
% température en dehors de \SI{}{} (non fournie par les paquets chargés).
\providecommand{\degree}{^{\circ}}
% \qed (fin de démonstration), utilisé par les modèles sans amsthm
\providecommand{\qed}{\hfill$\square$}
% \barre : double barre des valeurs interdites dans un tableau de variations (tkz-tab)
\providecommand{\barre}{\ensuremath{\|}}
% environnement proof (amsthm non chargé) utilisé par les modèles
\ifdefined\proof\else\newenvironment{proof}[1][Démonstration]{\par\noindent\textit{#1.}\ }{\qed\par}\fi
\usepackage{lastpage}
\usepackage{hyperref}
\hypersetup{hidelinks}

% ============================================================
% Palette « Pop » OnBuch+ — mêmes hex que la classe P (plus_ui.dart)
% ============================================================
\definecolor{poporange}{HTML}{F59321}
\definecolor{popdark}{HTML}{E07A0C}
\definecolor{popcream}{HTML}{FFF6E8}
\definecolor{popink}{HTML}{1C1714}
\definecolor{poppurple}{HTML}{7A5AE0}
\definecolor{popgreen}{HTML}{2E9E6A}
\definecolor{poppink}{HTML}{E0457B}
\definecolor{popblue}{HTML}{2F7DE1}
\definecolor{popgold}{HTML}{C98A12}
\definecolor{popmuted}{HTML}{8A7F76}
\definecolor{popline}{HTML}{EADFD2}
\definecolor{popgray}{HTML}{8A7F76}
\colorlet{popgrey}{popgray}
% Alias tolérants (le modèle nomme parfois les couleurs différemment)
\colorlet{popwhite}{white}
\colorlet{popblack}{popink}
\colorlet{popred}{poppink}
\colorlet{popyellow}{popgold}
% Teintes claires (fonds de tableaux/encadrés) — le modèle nomme parfois
% les couleurs avec un suffixe L (ex. popblueL) sans les avoir définies.
\colorlet{poporangeL}{poporange!20}
\colorlet{popdarkL}{popdark!20}
\colorlet{poppurpleL}{poppurple!20}
\colorlet{popgreenL}{popgreen!20}
\colorlet{poppinkL}{poppink!20}
\colorlet{popblueL}{popblue!20}
\colorlet{popgoldL}{popgold!20}
\colorlet{popredL}{popred!20}
\colorlet{popgrayL}{popgray!20}
% Teintes claires pour fonds de tableaux / zones
\colorlet{popblueL}{popblue!10!white}
\colorlet{poporangeL}{poporange!14!white}
\colorlet{popgreenL}{popgreen!12!white}
\colorlet{poppurpleL}{poppurple!12!white}
\colorlet{poppinkL}{poppink!12!white}
\colorlet{popgoldL}{popgold!14!white}

\pagecolor{popcream}

% ============================================================
% Typographie
% ============================================================
\defaultfontfeatures{Ligatures=TeX}
\setmainfont{PlusJakartaSans-Medium.ttf}[Path=../../../fonts/,BoldFont=PlusJakartaSans-Bold.ttf]
% Maths en sans-serif (Fira Math) avec chiffres et lettres droites de Plus
% Jakarta, pour que les formules se fondent dans le texte.
\usepackage[math-style=ISO,bold-style=ISO]{unicode-math}
\setmathfont{FiraMath-Regular.otf}[Path=../../../fonts/]
\setmathfont{PlusJakartaSans-Medium.ttf}[Path=../../../fonts/,range={up/num,up/{Latin,latin},\mathup}]
\usepackage{icomma} % virgule décimale sans espace en mode math
% Caractères mathématiques Unicode absents des polices de texte (Plus Jakarta,
% Archivo Black) : rendus par leur équivalent mathématique, en texte comme en
% formule — sinon ils s'affichent en carré vide (ex. « Arithmétique dans ℤ »).
\usepackage{newunicodechar}
\newunicodechar{ℝ}{\ensuremath{\mathbb{R}}}
\newunicodechar{ℤ}{\ensuremath{\mathbb{Z}}}
\newunicodechar{ℂ}{\ensuremath{\mathbb{C}}}
\newunicodechar{ℕ}{\ensuremath{\mathbb{N}}}
\newunicodechar{ℚ}{\ensuremath{\mathbb{Q}}}
\newunicodechar{𝔻}{\ensuremath{\mathbb{D}}}
\newunicodechar{α}{\ensuremath{\alpha}}
\newunicodechar{β}{\ensuremath{\beta}}
\newunicodechar{γ}{\ensuremath{\gamma}}
\newunicodechar{δ}{\ensuremath{\delta}}
\newunicodechar{ε}{\ensuremath{\varepsilon}}
\newunicodechar{θ}{\ensuremath{\theta}}
\newunicodechar{λ}{\ensuremath{\lambda}}
\newunicodechar{μ}{\ensuremath{\mu}}
\newunicodechar{σ}{\ensuremath{\sigma}}
\newunicodechar{τ}{\ensuremath{\tau}}
\newunicodechar{φ}{\ensuremath{\varphi}}
\newunicodechar{ψ}{\ensuremath{\psi}}
\newunicodechar{ω}{\ensuremath{\omega}}
\newunicodechar{Σ}{\ensuremath{\Sigma}}
% majuscules produites par \MakeUppercase dans les titres (α-aminés -> Α)
\newunicodechar{Α}{\ensuremath{\alpha}}
\newunicodechar{Β}{\ensuremath{\beta}}
\newunicodechar{Γ}{\ensuremath{\Gamma}}
\newunicodechar{Δ}{\ensuremath{\Delta}}
\newunicodechar{Ε}{\ensuremath{\varepsilon}}
\newunicodechar{Θ}{\ensuremath{\Theta}}
\newunicodechar{Λ}{\ensuremath{\Lambda}}
\newunicodechar{Μ}{\ensuremath{\mu}}
\newunicodechar{Τ}{\ensuremath{\tau}}
\newunicodechar{Φ}{\ensuremath{\Phi}}
\newunicodechar{Ψ}{\ensuremath{\Psi}}
\newunicodechar{Ω}{\ensuremath{\Omega}}
\newunicodechar{↦}{\ensuremath{\mapsto}}
\newunicodechar{∈}{\ensuremath{\in}}
\newunicodechar{∉}{\ensuremath{\notin}}
\newunicodechar{∧}{\ensuremath{\wedge}}
\newunicodechar{⇔}{\ensuremath{\Leftrightarrow}}
\newunicodechar{⟺}{\ensuremath{\Longleftrightarrow}}
\newunicodechar{⇒}{\ensuremath{\Rightarrow}}
\newunicodechar{⟹}{\ensuremath{\Longrightarrow}}
\newunicodechar{∘}{\ensuremath{\circ}}
\newunicodechar{∪}{\ensuremath{\cup}}
\newunicodechar{∩}{\ensuremath{\cap}}
\newunicodechar{≡}{\ensuremath{\equiv}}
\newunicodechar{⊂}{\ensuremath{\subset}}
\newunicodechar{⊥}{\ensuremath{\perp}}
\newunicodechar{∃}{\ensuremath{\exists}}
\newunicodechar{∀}{\ensuremath{\forall}}
\newunicodechar{∛}{\ensuremath{\sqrt[3]{\,}}}
\newunicodechar{∜}{\ensuremath{\sqrt[4]{\,}}}
\newunicodechar{⃗}{\ensuremath{{}^{\to}}}
\newunicodechar{ⁿ}{\ifmmode^{n}\else\textsuperscript{\ensuremath{n}}\fi}
\newunicodechar{ˣ}{\ifmmode^{x}\else\textsuperscript{\ensuremath{x}}\fi}
\newunicodechar{ᵇ}{\ifmmode^{b}\else\textsuperscript{\ensuremath{b}}\fi}
\newunicodechar{ʳ}{\ifmmode^{r}\else\textsuperscript{\ensuremath{r}}\fi}
\newunicodechar{ᵘ}{\ifmmode^{u}\else\textsuperscript{\ensuremath{u}}\fi}
\newunicodechar{ʸ}{\ifmmode^{y}\else\textsuperscript{\ensuremath{y}}\fi}
\newunicodechar{ᵃ}{\ifmmode^{a}\else\textsuperscript{\ensuremath{a}}\fi}
\newunicodechar{ᵉ}{\ifmmode^{e}\else\textsuperscript{\ensuremath{e}}\fi}
\newunicodechar{ᵏ}{\ifmmode^{k}\else\textsuperscript{\ensuremath{k}}\fi}
\newunicodechar{ᵖ}{\ifmmode^{p}\else\textsuperscript{\ensuremath{p}}\fi}
\newunicodechar{ᴺ}{\ifmmode^{N}\else\textsuperscript{\ensuremath{N}}\fi}
\newunicodechar{ᶜ}{\ifmmode^{c}\else\textsuperscript{\ensuremath{c}}\fi}
\newunicodechar{ʰ}{\ifmmode^{h}\else\textsuperscript{\ensuremath{h}}\fi}
\newunicodechar{ᵠ}{\ifmmode^{q}\else\textsuperscript{\ensuremath{q}}\fi}
\newunicodechar{ₙ}{\ifmmode_{n}\else\textsubscript{\ensuremath{n}}\fi}
\newunicodechar{ₐ}{\ifmmode_{a}\else\textsubscript{\ensuremath{a}}\fi}
\newunicodechar{ᵢ}{\ifmmode_{i}\else\textsubscript{\ensuremath{i}}\fi}
\newunicodechar{ᵣ}{\ifmmode_{r}\else\textsubscript{\ensuremath{r}}\fi}
\newunicodechar{ₖ}{\ifmmode_{k}\else\textsubscript{\ensuremath{k}}\fi}
\newunicodechar{ᵦ}{\ifmmode_{\beta}\else\textsubscript{\ensuremath{\beta}}\fi}
\newunicodechar{ₓ}{\ifmmode_{x}\else\textsubscript{\ensuremath{x}}\fi}
\newfontfamily\popdisplay{ArchivoBlack-Regular.ttf}[Path=../../../fonts/]
\newfontfamily\poplogo{SpaceGrotesk-Bold.ttf}[Path=../../../fonts/]
\newfontfamily\popsemi{PlusJakartaSans-SemiBold.ttf}[Path=../../../fonts/]
\newfontfamily\popxbold{PlusJakartaSans-ExtraBold.ttf}[Path=../../../fonts/]
\newfontfamily\popxboldls{PlusJakartaSans-ExtraBold.ttf}[Path=../../../fonts/,LetterSpace=3.0]

\setlist[enumerate]{leftmargin=1.7em,itemsep=5pt,topsep=4pt}
\setlist[itemize]{leftmargin=1.7em,itemsep=4pt,topsep=4pt,label={\color{poporange}\textbullet}}
\setlength{\parindent}{0pt}
\setlength{\parskip}{5pt}
\renewcommand{\arraystretch}{1.25}

% pgfplots : style Pop par défaut
\pgfplotsset{
  every axis/.append style={axis line style={popink,line width=1pt},tick style={popink},
    grid style={popline,line width=0.5pt},label style={font=\small},tick label style={font=\footnotesize}},
  popcurve/.style={poporange,line width=1.8pt,no markers,smooth},
}
% Même style disponible dans un \draw TikZ simple (hors pgfplots)
\tikzset{popcurve/.style={poporange,line width=1.8pt}}
\tikzset{popgrid/.style={popline,very thin}}
\colorlet{popcurve}{poporange} % le nom du style est parfois utilisé comme couleur

% ============================================================
% Pill / squiggle
% ============================================================
\newcommand{\poppill}[3]{%
  \tikz[baseline=-0.55ex]\node[draw=popink,line width=1.2pt,rounded corners=2.6pt,%
    fill=#2,inner xsep=6.5pt,inner ysep=3pt]{%
    \popxboldls\fontsize{7}{7}\selectfont\color{#3}\MakeUppercase{#1}};%
}
\newcommand{\popsquiggle}[1]{%
  \tikz[baseline=-0.4ex]{
    \draw[#1,line width=2.6pt,line cap=round]
      plot [smooth] coordinates {(0,0.09) (0.34,0.01) (0.68,0.21) (1.02,0.01) (1.36,0.10)};
  }%
}
% Mot-clé surligné (marqueur orange sous le texte)
\newcommand{\cle}[1]{{\popxbold\color{popdark}#1}}

% ============================================================
% En-tête / pied
% ============================================================
\pagestyle{fancy}
\fancyhf{}
\renewcommand{\headrulewidth}{0pt}
\renewcommand{\footrulewidth}{0pt}
\lhead{\raisebox{-0.15ex}{\poplogo\fontsize{13}{13}\selectfont\color{popink}OnBuch\color{poporange}+}}
\rhead{\poppill{\DOCMATIERE\ \textbullet\ \DOCNIVEAU\ \textbullet\ Leçon \DOCLECON}{white}{popink}}
\lfoot{\popxbold\fontsize{7.6}{7.6}\selectfont\color{popmuted}\MakeUppercase{Cours signé OnBuch+}\ \textbullet\ {\popsemi\DOCTITRE}}
\rfoot{\tikz[baseline=-0.6ex]\node[rounded corners=8.5pt,fill=popink,minimum height=17pt,minimum width=17pt,inner xsep=5pt,inner sep=0]{\popxbold\fontsize{8}{8}\selectfont\color{popcream}\thepage\,/\,\pageref*{LastPage}};}

% ============================================================
% Titres
% ============================================================
\newcommand{\chaptitre}[2]{%
  \begin{tcolorbox}[enhanced,colback=poporange,colframe=popink,boxrule=2.6pt,arc=11pt,
      drop shadow=popink,left=15pt,right=15pt,top=12pt,bottom=11pt,before skip=3pt,after skip=18pt]
    \poppill{#2}{popink}{popcream}\par\vspace{9pt}
    {\raggedright\hyphenpenalty=10000\exhyphenpenalty=10000\popdisplay\fontsize{22}{24}\selectfont\color{popcream}\MakeUppercase{#1}\par}
  \end{tcolorbox}
}
% Section numérotée : pastille orange avec le numéro + titre Archivo
\newcounter{popsec}
\newcommand{\coursec}[1]{%
  \stepcounter{popsec}\setcounter{popsub}{0}%
  \par\vspace{16pt}\noindent
  \begin{minipage}{\linewidth}
  \tikz[baseline=-0.7ex]\node[draw=popink,line width=1.6pt,fill=poporange,rounded corners=4pt,minimum size=22pt,inner sep=0]{\popdisplay\fontsize{12}{12}\selectfont\color{popcream}\thepopsec};%
  \hspace{8pt}{\popdisplay\fontsize{15}{17}\selectfont\color{popink}\MakeUppercase{#1}}\par
  \vspace{4pt}\noindent{\color{poporange}\rule{36pt}{3.6pt}}
  \end{minipage}\par\nopagebreak\vspace{8pt}
}
\newcounter{popsub}
\newcommand{\courssub}[1]{%
  \stepcounter{popsub}%
  \par\vspace{9pt}\noindent{\popxbold\fontsize{11.5}{13}\selectfont\color{popink}{\color{poporange}\thepopsec.\thepopsub}\hspace{6pt}#1}\par\nopagebreak\vspace{4pt}
}

% ============================================================
% Boîtes pédagogiques — même recette : fond blanc, bordure épaisse colorée,
% ombre dure encre, badge plein. Titre optionnel [#1].
% ============================================================
\tcbset{popbase/.style={breakable,enhanced,colback=white,boxrule=2.3pt,arc=9pt,
  shadow={3pt}{-3pt}{0pt}{popink},left=14pt,right=14pt,top=11pt,bottom=11pt,before skip=11pt,after skip=15pt,
  fontupper=\popsemi\fontsize{10.6}{14.5}\selectfont\color{popink},
  fontlower=\popsemi\fontsize{10.6}{14.5}\selectfont\color{popink}}}
% Coche dessinée (le glyphe ✓ n'existe pas dans les polices du document)
\newcommand{\popcheck}[1]{\tikz[baseline=-0.5ex]\draw[#1,line width=1.6pt,line cap=round,line join=round](-0.13,0.0)--(-0.04,-0.09)--(0.13,0.1);}
\providecommand{\coche}{\popcheck{popgreen}}
\providecommand{\resultat}[1]{\par\smallskip\noindent\fcolorbox{popgreen}{popgreenL}{\parbox{\dimexpr\linewidth-2\fboxsep-2\fboxrule\relax}{\textbf{Résultat.} #1}}\par\smallskip}
\newcommand{\popboxhead}[3]{\poppill{#1}{#2}{white}\ifx\relax#3\relax\else\hspace{6pt}{\popxbold\fontsize{10.6}{12}\selectfont\color{popink}#3}\fi\par\vspace{7pt}}

\newtcolorbox{aretenir}[1][]{popbase,colframe=popgreen,before upper={\popboxhead{À retenir}{popgreen}{#1}}}
\newtcolorbox{exemplebox}[1][]{popbase,colframe=poporange,before upper={\popboxhead{Exemple}{poporange}{#1}}}
\newtcolorbox{definition}[1][]{popbase,colframe=popblue,before upper={\popboxhead{Définition}{popblue}{#1}}}
\newtcolorbox{propriete}[1][]{popbase,colframe=poppurple,before upper={\popboxhead{Propriété}{poppurple}{#1}}}
\newtcolorbox{methode}[1][]{popbase,colframe=popgold,before upper={\popboxhead{Méthode}{popgold}{#1}}}
\newtcolorbox{attention}[1][]{popbase,colframe=poppink,before upper={\popboxhead{Attention}{poppink}{#1}}}
\newtcolorbox{experience}[1][]{popbase,colframe=popblue,colback=popblueL,before upper={\popboxhead{Expérience}{popblue}{#1}}}
% « Le savais-tu ? » : carte violette pleine (contexte, culture, Cameroun)
\newtcolorbox{savaistu}[1][]{popbase,colback=poppurple,colframe=popink,
  fontupper=\popsemi\fontsize{10.4}{14.2}\selectfont\color{white},
  before upper={\poppill{Le savais-tu\,?}{white}{poppurple}\ifx\relax#1\relax\else\hspace{6pt}{\popxbold\color{white}#1}\fi\par\vspace{7pt}}}
% Exercice résolu : énoncé en haut, \tcblower puis la solution sur fond crème
\newcounter{popexo}\newcounter{popexr}
\newtcolorbox{exoresolu}[1][]{popbase,colframe=poporange,colbacklower=popcream,
  segmentation style={popink,line width=1.2pt,dashed},
  before upper={\stepcounter{popexr}\popboxhead{Exercice résolu \thepopexr}{poporange}{#1}},
  before lower={\poppill{Solution}{popink}{popcream}\par\vspace{6pt}}}
% Exercice (non résolu en place) — la correction est dans la partie Corrigés
\newtcolorbox{exercice}[1][]{popbase,colframe=popink,
  before upper={\stepcounter{popexo}\popboxhead{Exercice \thepopexo}{popink}{#1}}}
% Corrigé
\newtcolorbox{corrige}[1][]{popbase,colframe=popgreen,colback=popgreenL,
  before upper={\popboxhead{Corrigé}{popgreen}{#1}}}
% Objectifs / prérequis / bilan
\newtcolorbox{objectifs}{popbase,colframe=popink,colback=poporangeL,
  before upper={\popboxhead{Objectifs de la leçon}{popink}{}}}
\newtcolorbox{prerequis}{popbase,colframe=popink,colback=popblueL,
  before upper={\popboxhead{Prérequis}{popblue}{}}}
\newtcolorbox{bilan}{popbase,colframe=popink,colback=white,boxrule=2.8pt,
  before upper={\popboxhead{Fiche bilan}{poporange}{L'essentiel à connaître pour l'examen}}}
% Figure encadrée avec légende
\newtcolorbox{popfigure}[1][]{enhanced,breakable=false,colback=white,colframe=popline,boxrule=1.4pt,arc=8pt,
  left=10pt,right=10pt,top=10pt,bottom=8pt,before skip=10pt,after skip=12pt,halign=center,
  lower separated=false,halign lower=center,
  fontlower=\popsemi\fontsize{9}{11.5}\selectfont\color{popmuted},
  #1}
\newcommand{\legende}[1]{\par\vspace{6pt}{\popsemi\fontsize{9}{11.5}\selectfont\color{popmuted}#1}}

% ============================================================
% Signature OnBuch+
% ============================================================
\newcommand{\poplogoinline}[1]{{\poplogo\fontsize{#1}{#1}\selectfont\color{popink}OnBuch\color{poporange}+}}
\newcommand{\signatureonbuch}{%
  \begin{tcolorbox}[enhanced,colback=popink,colframe=popink,boxrule=2.2pt,arc=10pt,
      drop shadow=poporange,left=14pt,right=14pt,top=10pt,bottom=10pt,before skip=0pt,after skip=0pt]
    \tikz[baseline=-0.7ex]\node[circle,fill=popgreen,draw=popcream,line width=1.3pt,minimum size=22pt,inner sep=0]{\popcheck{white}};%
    \hspace{9pt}\begin{minipage}[c]{0.62\linewidth}
      {\popxbold\fontsize{12}{13}\selectfont\color{popcream}Signé\ \textbullet\ {\poplogo OnBuch\color{poporange}+}}\par\vspace{2pt}
      {\raggedright\popsemi\fontsize{8}{10.5}\selectfont\color{popline}\MakeUppercase{Équipe pédagogique OnBuch+}\\\MakeUppercase{Conforme au programme officiel MINESEC}\par}
    \end{minipage}\hfill
    {\poplogo\fontsize{22}{22}\selectfont\color{popcream}OnBuch\color{poporange}+}
  \end{tcolorbox}%
}

% ============================================================
% Page de garde
% #1 titre · #2 matière · #3 niveau · #4 module · #5 n° leçon · #6 durée ·
% #7 description / objectifs (texte ou itemize)
% ============================================================
\newcommand{\pagegarde}[7]{%
  \thispagestyle{empty}
  \newgeometry{a4paper,margin=1.7cm,top=1.6cm,bottom=1.4cm}%
  \begin{tikzpicture}[remember picture,overlay]
    \fill[poporange!18] ([xshift=-3.2cm,yshift=-3.6cm]current page.north east) circle (4.6cm);
    \fill[poppurple!14] ([xshift=2.4cm,yshift=7.6cm]current page.south west) circle (3.8cm);
  \end{tikzpicture}%
  {\poplogo\fontsize{30}{30}\selectfont\color{popink}OnBuch\color{poporange}+}\hfill
  \raisebox{0.6ex}{\poppill{Cours signé \textbullet\ Premium}{popink}{popcream}}\par
  \vspace{1pt}\popsquiggle{poppurple}\par
  \vspace{8pt}
  \begin{tcolorbox}[enhanced,colback=poporange,colframe=popink,boxrule=2.8pt,arc=13pt,
      drop shadow=popink,left=18pt,right=18pt,top=12pt,bottom=13pt,after skip=10pt]
    \poppill{#2 \textbullet\ #3}{popink}{popcream}\hspace{5pt}\poppill{#4}{white}{popink}\par\vspace{8pt}
    {\popdisplay\fontsize{14}{14}\selectfont\color{popink}LEÇON #5}\par\vspace{4pt}
    {\raggedright\hyphenpenalty=10000\exhyphenpenalty=10000\popdisplay\fontsize{25}{27}\selectfont\color{popcream}\MakeUppercase{#1}\par}
    \vspace{8pt}
    {\popxbold\fontsize{9.5}{9.5}\selectfont\color{popink}\MakeUppercase{Durée conseillée : #6}}
  \end{tcolorbox}
  \begin{tcolorbox}[enhanced,colback=white,colframe=popink,boxrule=2.2pt,arc=10pt,
      drop shadow=popink,left=16pt,right=16pt,top=10pt,bottom=10pt,after skip=8pt]
    \poppill{Dans cette leçon}{poppurple}{white}\par\vspace{5pt}
    \popsemi\fontsize{9.3}{12.6}\selectfont\color{popink}#7
  \end{tcolorbox}
  \noindent\begin{minipage}[t]{0.487\linewidth}
    \begin{tcolorbox}[enhanced,colback=popgreen,colframe=popink,boxrule=2.2pt,arc=10pt,
        drop shadow=popink,left=11pt,right=11pt,top=10pt,bottom=10pt,height=3.1cm,valign=center]
      \tcbox[colback=white,colframe=popink,boxrule=1.3pt,arc=5pt,boxsep=0pt,nobeforeafter,
        left=3pt,right=3pt,top=3pt,bottom=3pt,valign=center]{\qrcode[height=1.9cm]{https://play.google.com/store/apps/details?id=com.luvvix.onbuch&hl=fr}}%
      \hspace{8pt}\begin{minipage}[c]{0.5\linewidth}
        {\popdisplay\fontsize{11}{12.5}\selectfont\color{white}\MakeUppercase{Télécharge OnBuch}}\par\vspace{4pt}
        {\raggedright\popsemi\fontsize{8.4}{11}\selectfont\color{white}Gratuit \textbullet\ Google Play\\Tuteur IA, annales, concours, résultats\par}
      \end{minipage}
    \end{tcolorbox}
  \end{minipage}\hfill
  \begin{minipage}[t]{0.487\linewidth}
    \begin{tcolorbox}[enhanced,colback=poppurple,colframe=popink,boxrule=2.2pt,arc=10pt,
        drop shadow=popink,left=11pt,right=11pt,top=10pt,bottom=10pt,height=3.1cm,valign=center]
      \tikz[baseline=-0.7ex]\node[circle,fill=white,draw=popink,line width=1.3pt,minimum size=1.6cm,inner sep=0]{\popdisplay\fontsize{24}{24}\selectfont\color{poppurple}+};%
      \hspace{8pt}\begin{minipage}[c]{0.6\linewidth}
        {\popdisplay\fontsize{11}{12.5}\selectfont\color{white}\MakeUppercase{Passe à OnBuch+}}\par\vspace{4pt}
        {\raggedright\popsemi\fontsize{8.4}{11}\selectfont\color{white}Tous les cours signés, Tuteur IA illimité, fiches PDF — onglet OnBuch+ de l'app\par}
      \end{minipage}
    \end{tcolorbox}
  \end{minipage}
  \vspace{8pt}
  \popcontenu
  \vfill
  \signatureonbuch
  \restoregeometry
  \newpage
}

% Rangée « Ce cours contient » (page de garde)
\newcommand{\poptile}[3]{% #1 couleur #2 gros texte #3 légende
  \begin{tcolorbox}[enhanced,colback=#1,colframe=popink,boxrule=2pt,arc=9pt,shadow={2.5pt}{-2.5pt}{0pt}{popink},
      left=6pt,right=6pt,top=9pt,bottom=9pt,halign=center,valign=center,height=1.75cm,width=\linewidth,nobeforeafter]
    {\popdisplay\fontsize{12.5}{13}\selectfont\color{white}\MakeUppercase{#2}}\par\vspace{3pt}
    {\centering\popsemi\fontsize{7.8}{9.5}\selectfont\color{white}#3\par}
  \end{tcolorbox}}
\newcommand{\popcontenu}{%
  \par\noindent{\popxboldls\fontsize{7.5}{7.5}\selectfont\color{popmuted}CE COURS SIGNÉ CONTIENT}\par\vspace{7pt}
  \noindent\begin{minipage}[t]{0.235\linewidth}\poptile{popblue}{Cours}{complet, illustré, pas à pas}\end{minipage}\hfill
  \begin{minipage}[t]{0.235\linewidth}\poptile{poporange}{Méthodes}{et exercices résolus}\end{minipage}\hfill
  \begin{minipage}[t]{0.235\linewidth}\poptile{poppink}{Exercices}{type examen + corrigés}\end{minipage}\hfill
  \begin{minipage}[t]{0.235\linewidth}\poptile{popgreen}{Fiche bilan}{l'essentiel à retenir}\end{minipage}\par}

% Sommaire
\newtcolorbox{sommaire}{enhanced,colback=white,colframe=popink,boxrule=2.2pt,arc=10pt,
  drop shadow=popink,left=18pt,right=18pt,top=13pt,bottom=13pt,before skip=6pt,after skip=20pt,
  before upper={\poppill{Sommaire}{poppurple}{white}\par\vspace{9pt}\popsemi\fontsize{10.6}{14.5}\selectfont\color{popink}}}

% Signature de fin de cours — poussée tout en bas de la dernière page
\newcommand{\findecours}{%
  \par\vfill\nopagebreak
  \begin{center}\popsquiggle{poporange}\end{center}\vspace{4pt}
  \signatureonbuch
}
\AtBeginDocument{\def\llbracket{\mathopen{[\mkern-3.5mu[}}\def\rrbracket{\mathclose{]\mkern-3.5mu]}}} % intervalles d'entiers (glyphe ⟦ absent de la police)
% Glyphes absents de Fira Math : redéfinis à partir de glyphes disponibles
\AtBeginDocument{%
  \def\setminus{\mathbin{\backslash}}\let\smallsetminus\setminus
  \def\vdots{\mathord{\vbox{\baselineskip3.2pt\lineskiplimit0pt\kern4pt\hbox{.}\hbox{.}\hbox{.}}}}%
  \def\ddots{\mathinner{\mkern1mu\raise6.5pt\hbox{.}\mkern2mu\raise3.6pt\hbox{.}\mkern2mu\raise0.7pt\hbox{.}\mkern1mu}}%
  \def\triangle{\mathord{\scalebox{0.9}{$\symup{\Delta}$}}}%
  \def\nabla{\mathord{\raisebox{\depth}{\scalebox{1}[-1]{$\symup{\Delta}$}}}}%
  \def\checkmark{\popcheck{popdark}}%
  \def\star{\mathbin{\ast}}\def\bigstar{\mathord{\ast}}%
  \def\diamond{\mathbin{\scalebox{0.6}{$\Diamond$}}}%
  \def\bigcup{\mathop{\mathchoice{\raisebox{-0.3em}{\scalebox{1.7}{$\cup$}}}{\scalebox{1.25}{$\cup$}}{\cup}{\cup}}\displaylimits}%
  \def\bigcap{\mathop{\mathchoice{\raisebox{-0.3em}{\scalebox{1.7}{$\cap$}}}{\scalebox{1.25}{$\cap$}}{\cap}{\cap}}\displaylimits}%
  \def\lesseqqgtr{\mathrel{\lessgtr}}\def\bigoplus{\mathop{\oplus}\displaylimits}%
}
\providecommand{\ssi}{si et seulement si} % abréviation fréquente
\AtBeginDocument{\providecommand{\wideparen}{\overparen}} % arc de cercle

%% FIGFIX-BEGIN v4 — correctifs globaux des figures (docs/onbuch-generation/tools/figfix)
\usepackage{adjustbox}\usepackage{etoolbox}
% toute figure TikZ/pgfplots plus large que la ligne est réduite à la largeur disponible
\BeforeBeginEnvironment{tikzpicture}{\begin{adjustbox}{max width=\linewidth}}
\AfterEndEnvironment{tikzpicture}{\end{adjustbox}}
% légendes pgfplots lisibles par-dessus les courbes
\pgfplotsset{every axis/.append style={legend style={fill=white,fill opacity=0.92,text opacity=1,draw=black!15,inner sep=3pt}}}
% étiquettes d'arêtes (syntaxe edge["..."]) avec fond blanc
\tikzset{every edge quotes/.append style={fill=white,inner sep=1.2pt}}
%% FIGFIX-END
\begin{document}

\pagegarde{V : produire une lettre officielle ( lettre de motivation)}{Français}{1re Tech}{Français – classe de Première technique}{N21}{10 séances (fiche de progression harmonisée)}{Cette leçon outille l'élève de Première technique pour produire une lettre officielle de motivation, document indispensable à l'insertion professionnelle et aux demandes de stage au Cameroun. Elle articule l'étude de la structure externe et interne de la lettre, l'entraînement à la rédaction (introduction, corps, conclusion) et les outils langagiers adaptés : figures de style, impératif, texte descriptif et communication par l'image.
\vspace{4pt}
\begin{multicols}{2}\small
\begin{itemize}
\item À la fin de cette leçon, je sais identifier les caractéristiques d'un texte fonctionnel et lire efficacement une lettre officielle.
\item À la fin de cette leçon, je sais décrire et reproduire la structure externe d'une lettre officielle (en-tête, objet, formules, signature).
\item À la fin de cette leçon, je sais construire la structure interne d'une lettre de motivation (introduction, corps, conclusion).
\item À la fin de cette leçon, je sais rédiger une introduction et une conclusion conformes aux usages administratifs camerounais.
\item À la fin de cette leçon, je sais rédiger le corps d'une lettre de motivation en valorisant mon profil et mes compétences techniques.
\item À la fin de cette leçon, je sais employer l'impératif et les figures de style (hyperbole, métonymie, gradation, énumération) à bon escient dans un écrit formel.
\end{itemize}\end{multicols}}

\begin{sommaire}
\begin{multicols}{2}
\begin{enumerate}
\item Les textes fonctionnels et la lettre officielle
\item La structure externe de la lettre officielle
\item La structure interne : introduction et conclusion
\item Le corps de la lettre et les outils langagiers
\item La communication par l'image dans la candidature
\item Production guidée, compte rendu et remédiation
\item Activité d'intégration
\item Méthodes et exercices résolus
\item Exercices
\item Corrigés des exercices
\item Fiche bilan
\end{enumerate}
\end{multicols}
\end{sommaire}

\begin{objectifs}
\begin{itemize}
\item À la fin de cette leçon, je sais identifier les caractéristiques d'un texte fonctionnel et lire efficacement une lettre officielle.
\item À la fin de cette leçon, je sais décrire et reproduire la structure externe d'une lettre officielle (en-tête, objet, formules, signature).
\item À la fin de cette leçon, je sais construire la structure interne d'une lettre de motivation (introduction, corps, conclusion).
\item À la fin de cette leçon, je sais rédiger une introduction et une conclusion conformes aux usages administratifs camerounais.
\item À la fin de cette leçon, je sais rédiger le corps d'une lettre de motivation en valorisant mon profil et mes compétences techniques.
\item À la fin de cette leçon, je sais employer l'impératif et les figures de style (hyperbole, métonymie, gradation, énumération) à bon escient dans un écrit formel.
\item À la fin de cette leçon, je sais exploiter la communication par l'image (photo d'identité, mise en page) pour renforcer ma candidature.
\item À la fin de cette leçon, je sais produire, seul et en temps limité, une lettre de motivation complète répondant à une offre d'emploi ou de stage réelle.
\item À la fin de cette leçon, je sais justifier mes choix rédactionnels et appliquer ces compétences à des situations concrètes de la vie courante.
\end{itemize}
\end{objectifs}

\begin{prerequis}
\begin{itemize}
\item Les types de textes et les genres littéraires (classe de Seconde)
\item La phrase simple et la phrase complexe ; la ponctuation
\item Les formules de politesse élémentaires et le registre de langue soutenu
\item Le réseau d'expansion du nom (adjectif, complément du nom) utile à la description
\item Les connecteurs logiques de base (d'abord, ensuite, en effet, c'est pourquoi)
\end{itemize}
\end{prerequis}

\newpage

\coursec{Les textes fonctionnels et la lettre officielle}

\textbf{Situation d'accroche.} À Douala, Merveille, 19 ans, titulaire d'un Baccalauréat technique en maintenance industrielle, repère une offre d'emploi de la SABC : « Adresser une lettre de motivation au Directeur des Ressources Humaines ». Elle rédige à la hâte un message informel, sans objet, avec des fautes et un ton trop familier : sa candidature est rejetée sans même être lue jusqu'au bout. Son camarade Rodrigue, moins bon technicien qu'elle, obtient l'entretien grâce à une lettre claire, structurée et convaincante. Qu'est-ce qui fait la différence entre une lettre qui ouvre des portes et une lettre qui les ferme ? Comment rédiger une lettre de motivation capable de « vendre » son profil auprès d'un employeur camerounais ?

C'est à ces questions que répond cette première partie de la leçon : avant d'apprendre à rédiger, il faut comprendre \emph{à quoi sert} une lettre officielle, \emph{à qui} elle s'adresse et \emph{ce qui la distingue} des autres écrits de la vie quotidienne.

\courssub{Qu'est-ce qu'un texte fonctionnel ?}

\textbf{L'intuition d'abord.} Observe ta journée : tu lis une affiche de concours à la mairie, tu remplis un formulaire d'inscription, ton père paie une facture ENEO, ton grand frère envoie un CV. Aucun de ces textes n'est un poème ni un roman : personne ne les lit « pour le plaisir ». Ils existent pour \emph{faire quelque chose} : obtenir un document, payer une dette, décrocher un emploi. C'est exactement cela, un texte fonctionnel.

\begin{definition}[Le texte fonctionnel]
Un \cle{texte fonctionnel} (ou texte utilitaire, texte pragmatique) est un texte écrit pour \textbf{agir sur la réalité} : demander, obtenir, informer, convaincre, prouver. Il vise une \textbf{action concrète} et une réponse précise du destinataire. Il s'oppose au \cle{texte littéraire} (roman, poème, conte), écrit d'abord pour le plaisir, l'émotion et la réflexion.
\end{definition}

\begin{aretenir}[Les trois marques du texte fonctionnel]
\begin{itemize}
\item Une \textbf{finalité pratique} : le texte doit provoquer un effet réel (une embauche, une autorisation, un paiement).
\item Un \textbf{destinataire identifié} : on sait exactement à qui l'on écrit (le maire, le directeur, le client).
\item Une \textbf{forme codifiée} : la disposition, les formules et le ton obéissent à des conventions sociales qu'il faut respecter pour être pris au sérieux.
\end{itemize}
\end{aretenir}

\courssub{Typologie des textes fonctionnels}

Les textes fonctionnels sont partout dans la vie administrative et professionnelle. On peut les classer selon leur fonction dominante.

\begin{center}
\begin{tabularx}{\linewidth}{|l|X|X|}
\hline
\rowcolor{popblueL} \textbf{Texte fonctionnel} & \textbf{Fonction principale} & \textbf{Exemple camerounais} \\
\hline
Lettre administrative & Demander, réclamer, informer une autorité & Demande de bourse adressée au MINESUP \\
\hline
Lettre de motivation & Convaincre un employeur de recevoir le candidat & Candidature à un stage à la SODECOTON \\
\hline
Curriculum vitae (CV) & Présenter de façon synthétique diplômes et expériences & CV joint à tout dossier de candidature \\
\hline
Facture, quittance & Attester et réclamer un paiement & Facture ENEO, CAMWATER \\
\hline
Formulaire & Recueillir des informations normalisées & Fiche d'inscription à un concours de la Fonction publique \\
\hline
Affiche, avis & Informer un large public & Avis de concours d'entrée à l'ENAM \\
\hline
Convocation & Sommer officiellement de se présenter & Convocation à un entretien d'embauche \\
\hline
\end{tabularx}
\end{center}

\begin{exemplebox}[Une offre d'emploi réaliste]
« La CAMTEL recrute des techniciens en réseaux. Dossier à fournir : CV + lettre de motivation adressée au Directeur des Ressources Humaines, Yaoundé, au plus tard le 30 du mois. »
\par
Dans cette seule annonce, on rencontre trois textes fonctionnels : l'\textbf{affiche} (l'offre elle-même), le \textbf{CV} et la \textbf{lettre de motivation}. Le candidat qui comprend la fonction de chacun a déjà fait la moitié du chemin.
\end{exemplebox}

\courssub{La lettre officielle : définition et destinataires}

\begin{definition}[La lettre officielle]
La \cle{lettre officielle} est un texte fonctionnel écrit, adressé à une \textbf{autorité} ou à une \textbf{personne investie d'une fonction} (et non à un proche), dans un cadre administratif ou professionnel. Elle respecte une présentation codifiée (lieu et date, coordonnées, objet, formules de politesse) et un ton \textbf{courtois et soutenu}.
\end{definition}

Les destinataires typiques d'une lettre officielle au Cameroun sont :
\begin{itemize}
\item le \textbf{ministre} (par exemple le Ministre de l'Enseignement supérieur pour une demande de bourse) ;
\item le \textbf{directeur} ou le \textbf{chef d'entreprise} (Directeur des Ressources Humaines de la SONARA, Directeur général d'une banque) ;
\item le \textbf{maire}, le \textbf{préfet}, le \textbf{recteur} pour les démarches administratives ;
\item le \textbf{chef d'établissement} (proviseur, directeur d'école) pour les demandes scolaires.
\end{itemize}

\textbf{Ce qui la distingue de la lettre amicale.} À un ami, tu peux écrire « Salut vieux, ça fait longtemps ! ». À un directeur, jamais. La lettre officielle exige : le vouvoiement, l'absence d'abréviations, des formules de politesse complètes, une orthographe irréprochable. Le ton n'est pas un détail : il \emph{est} le message. Une lettre familière dit à son lecteur : « je ne vous respecte pas assez pour faire un effort ».

\begin{attention}[Le piège du ton]
La faute la plus grave dans une lettre officielle n'est pas toujours l'orthographe : c'est le \textbf{registre de langue inadapté}. « Hé patron, file-moi le job » est rédhibitoire, même sans faute. Inversement, un ton obséquieux (« je me prosterne devant votre grandeur ») paraît faux. Vise la \textbf{courtoisie simple et digne}.
\end{attention}

\courssub{Lettre administrative ou lettre de motivation ?}

Ces deux lettres sont toutes les deux officielles, mais elles n'ont ni la même finalité, ni le même destinataire, ni le même contenu. Le tableau comparatif ci-dessous les oppose, en les situant par rapport à la lettre amicale.

\begin{popfigure}
\begin{center}
\begin{tikzpicture}[font=\small]
% En-têtes
\fill[popblueL] (0,0) rectangle (3.4,0.9);
\fill[popgreenL] (3.5,0) rectangle (7.4,0.9);
\fill[poporangeL] (7.5,0) rectangle (11.4,0.9);
\fill[poppurpleL] (11.5,0) rectangle (15.4,0.9);
\node[text width=3.2cm,align=center] at (1.7,0.45) {\textbf{Critère}};
\node[text width=3.7cm,align=center] at (5.45,0.45) {\textbf{Lettre amicale}};
\node[text width=3.7cm,align=center] at (9.45,0.45) {\textbf{Lettre administrative}};
\node[text width=3.7cm,align=center] at (13.45,0.45) {\textbf{Lettre de motivation}};
% Lignes
\foreach \y/\crit/\a/\b/\c in {
 -1.0/{Destinataire}/{un proche, un ami}/{une autorité : ministre, maire, préfet}/{un employeur : DRH, chef d'entreprise},
 -2.2/{Finalité}/{entretenir une relation}/{obtenir un acte : bourse, autorisation, réclamation}/{décrocher un emploi, un stage, une formation},
 -3.4/{Ton}/{familier, affectueux}/{courtois, neutre, formel}/{courtois et persuasif : il faut convaincre},
 -4.6/{Contenu}/{nouvelles, sentiments}/{exposé d'une demande précise avec justificatifs}/{valorisation du profil : compétences, motivation},
 -5.8/{Structure}/{libre}/{codifiée : objet, références, formules}/{codifiée + argumentée : accroche, arguments, relance}}{
 \node[text width=3.2cm,align=left,anchor=north] at (0.1,\y+0.55) {\textbf{\crit}};
 \node[text width=3.6cm,align=left,anchor=north] at (3.6,\y+0.55) {\a};
 \node[text width=3.6cm,align=left,anchor=north] at (7.6,\y+0.55) {\b};
 \node[text width=3.6cm,align=left,anchor=north] at (11.6,\y+0.55) {\c};
}
% Cadre
\draw[popline] (0,0.9) rectangle (15.4,-6.3);
\draw[popline] (3.5,0.9) -- (3.5,-6.3);
\draw[popline] (7.5,0.9) -- (7.5,-6.3);
\draw[popline] (11.5,0.9) -- (11.5,-6.3);
\foreach \y in {-0.3,-1.5,-2.7,-3.9,-5.1} {\draw[popline] (0,\y) -- (15.4,\y);}
\end{tikzpicture}
\end{center}
\legende{Tableau comparatif des trois grands types de lettres. La lettre de motivation partage la forme codifiée de la lettre administrative, mais sa finalité est persuasive : elle doit convaincre, pas seulement demander.}
\end{popfigure}

\begin{aretenir}[La différence essentielle]
La lettre administrative \textbf{demande} un acte à une autorité ; la lettre de motivation \textbf{vend} un profil à un employeur. La première insiste sur la demande et ses justificatifs ; la seconde insiste sur les compétences du candidat et sur ce qu'il peut apporter à l'entreprise.
\end{aretenir}

\courssub{La lettre de motivation : un document de candidature}

\begin{definition}[La lettre de motivation]
La \cle{lettre de motivation} est une lettre officielle par laquelle un candidat \textbf{expose ses raisons de postuler} et \textbf{valorise ses compétences} auprès d'un employeur, en vue d'obtenir un emploi, un stage ou une formation. Elle accompagne toujours le CV et vise un objectif précis : \textbf{décrocher un entretien}.
\end{definition}

Retiens bien la logique de l'ensemble : le \textbf{CV} est une \emph{photographie} du parcours (diplômes, expériences, compétences, présentées en listes) ; la \textbf{lettre} est le \emph{discours} qui donne du sens à cette photographie (« voici pourquoi ce parcours correspond exactement à votre besoin »).

\begin{attention}[Ne confonds pas lettre de motivation et CV]
Confondre lettre de motivation et CV est une erreur fréquente. La lettre \textbf{argumente et convainc} : elle raconte un projet, explique un choix, montre une motivation. Le CV \textbf{liste} les diplômes et les expériences, sans phrases rédigées. Recopier son CV en phrases dans la lettre est inutile : les deux documents sont \textbf{complémentaires}, pas interchangeables.
\end{attention}

\textbf{Les situations camerounaises d'usage.} La lettre de motivation n'est pas un exercice scolaire artificiel : c'est un outil de vie. Tu en rédigeras pour :
\begin{itemize}
\item une \textbf{demande de stage} à la SODECOTON de Garoua, à la CDC de Buea ou dans un atelier de maintenance de Douala ;
\item une \textbf{candidature à un emploi} à la SONARA de Limbé, à la SABC, à la CAMTEL ou dans une PME du quartier industriel ;
\item un \textbf{concours de la Fonction publique} (lettre de motivation parfois exigée dans le dossier) ;
\item une \textbf{demande de bourse} au MINESUP ou auprès d'une fondation ;
\item une \textbf{candidature spontanée} dans une entreprise qui n'a pas publié d'offre.
\end{itemize}

\begin{savaistu}[La candidature spontanée au Cameroun]
Au Cameroun, de nombreuses candidatures spontanées aboutissent dans les PME de Douala et de Yaoundé lorsque la lettre est soignée et personnalisée. Les petites entreprises ne publient pas toujours d'offres : un jeune technicien qui se présente avec une lettre claire, adaptée à l'entreprise, et un CV propre peut créer sa propre opportunité. Le chef d'entreprise retient le candidat qui montre qu'il connaît la société et qu'il peut être utile immédiatement.
\end{savaistu}

\courssub{Outils langagiers au service de la persuasion}

Le programme de Première technique prévoit l'étude de figures de style et du mode impératif : ces outils renforcent l'efficacité de la lettre de motivation.

\begin{propriete}[Figures de style utiles dans une lettre de motivation]
\begin{itemize}
\item \textbf{L'énumération} : liste les compétences, stages, logiciels maîtrisés (\emph{« Maîtrise de la soudure TIG, MIG, à l'arc ; lecture de plans ; maintenance préventive »}).
\item \textbf{La gradation} : ordonne les arguments du plus accessible au plus fort (\emph{« J'ai acquis des bases solides, puis développé une autonomie, enfin conduit une équipe de deux stagiaires »}).
\item \textbf{L'hyperbole} (avec mesure) : souligne l'enthousiasme sans mentir (\emph{« Une motivation sans faille »} plutôt que \emph{« Je suis le meilleur du monde »}).
\item \textbf{La métonymie} : désigne l'entreprise par son siège ou son produit (\emph{« Rejoindre Bonanjo »} pour le siège social, \emph{« Intégrer la raffinerie »} pour la SONARA).
\end{itemize}
\end{propriete}

\begin{propriete}[Le mode impératif : valeur et emploi]
Dans la conclusion, le candidat utilise souvent l'impératif \textbf{de politesse} (forme courtoise) pour appeler à l'action : \emph{« Je vous prie d'agréer... »}, \emph{« Veuillez recevoir... »}, \emph{« Je reste à votre disposition pour un entretien »}. Ce n'est pas un ordre, mais une convention de courtoisie qui marque le respect et l'attente d'une réponse.
\end{propriete}

\courssub{La communication par l'image dans la candidature}

\begin{aretenir}[L'image, premier message lu]
Avant même de lire le texte, le recruteur « lit » la \textbf{présentation visuelle} : propreté de la page, alignements, police lisible (Times New Roman, Arial, 12 pt), marges régulières, absence de ratures. Une lettre aérée, bien structurée, signale un candidat rigoureux. Le CV, lui, utilise des \textbf{pictogrammes} (téléphone, mail, localisation, LinkedIn) et une \textbf{hiérarchie visuelle} (titres en gras, puces) pour guider l'œil en quelques secondes. Soigner la forme, c'est déjà communiquer par l'image.
\end{aretenir}

\courssub{Lecture guidée : deux lettres, deux destins}

Revenons à notre situation d'accroche et observons, comme un enquêteur, les deux documents. Voici d'abord la lettre ratée de Merveille (extraits), puis celle de Rodrigue.

\begin{exemplebox}[La lettre ratée (extraits annotés)]
« Salut Monsieur, j'ai vu votre annonce et le poste m'intéresse grave. J'ai eu mon Bac l'an dernier, je suis doué en mécanique, tout le monde le dit. Appelez-moi vite. »
\par
\textbf{Repérage des fautes :} pas de lieu ni de date ; pas de coordonnées du destinataire ; pas d'objet ; tutoiement implicite et familiarités (« grave », « appelez-moi vite ») ; aucune formule de politesse ; affirmation sans preuve (« tout le monde le dit ») ; fautes d'orthographe. Résultat : la lettre dit au lecteur « je ne maîtrise ni le code ni le sérieux du monde professionnel ».
\end{exemplebox}

\begin{exemplebox}[La lettre réussie (extraits annotés)]
« Douala, le 12 octobre 2025. — À Monsieur le Directeur des Ressources Humaines de la SABC, Douala. — Objet : candidature au poste de technicien de maintenance industrielle. — Monsieur le Directeur, titulaire d'un Baccalauréat technique en maintenance industrielle, j'ai effectué un stage de trois mois dans l'atelier mécanique de la société ALUCAM, où j'ai participé à la maintenance préventive des lignes de production. Rejoindre la SABC, entreprise leader de l'agroalimentaire au Cameroun, me permettrait de mettre ce savoir-faire au service de vos équipements. [...] Je vous prie d'agréer, Monsieur le Directeur, l'expression de ma haute considération. »
\par
\textbf{Repérage des réussites :} en-tête complet, objet précis, vouvoiement, compétences \emph{prouvées} par une expérience concrète, connaissance de l'entreprise (« leader de l'agroalimentaire »), formule de politesse adaptée. La lettre dit : « je suis professionnel avant même d'être embauché ».
\end{exemplebox}

\begin{popfigure}
\begin{center}
\begin{tikzpicture}[font=\footnotesize]
% Feuille de lettre
\fill[popcream] (0,0) rectangle (10,12.5);
\draw[popdark] (0,0) rectangle (10,12.5);
% Contenu simulé
\node[anchor=north west,text width=4.5cm] at (0.4,12.2) {Merveille N.\\ Quartier Deïdo\\ Douala\\ Tél. : 6XX XX XX XX};
\node[anchor=north east,text width=4.5cm,align=right] at (9.6,12.2) {À Monsieur le DRH\\ de la SABC\\ Bonanjo, Douala};
\node[anchor=north west] at (6.2,9.6) {Douala, le 12 octobre 2025};
\node[anchor=north west,text width=9cm] at (0.4,8.9) {\textbf{Objet :} candidature au poste de technicien de maintenance};
\foreach \y in {7.9,7.5,7.1,6.7,6.3,5.9,5.5,5.1,4.7,4.3,3.9,3.5} {\draw[popmuted!60] (0.5,\y) -- (9.5,\y);}
\node[anchor=north west,text width=9cm] at (0.4,3.1) {Je vous prie d'agréer, Monsieur le Directeur, l'expression de ma haute considération.};
\node[anchor=north west] at (6.5,1.6) {Signature};
% Flèches d'annotation à gauche
\draw[->,popblue,thick] (-0.2,11.5) -- (0.4,11.5);
\node[anchor=east,text width=3.6cm,align=right,popblue] at (-0.3,11.5) {\textbf{1. Coordonnées de l'expéditeur} (en haut à gauche)};
\draw[->,popgreen,thick] (-0.2,8.6) -- (0.4,8.6);
\node[anchor=east,text width=3.6cm,align=right,popgreen] at (-0.3,8.6) {\textbf{3. Objet} : une seule ligne, précise};
\draw[->,poporange,thick] (-0.2,6.1) -- (0.5,6.1);
\node[anchor=east,text width=3.6cm,align=right,poporange] at (-0.3,6.1) {\textbf{4. Corps} : accroche, arguments, projet};
% Flèches d'annotation à droite
\draw[->,popblue,thick] (10.2,11.5) -- (9.6,11.5);
\node[anchor=west,text width=3.8cm,popblue] at (10.3,11.5) {\textbf{2. Destinataire + lieu et date} (à droite)};
\draw[->,popgreen,thick] (10.2,2.6) -- (9.5,2.6);
\node[anchor=west,text width=3.8cm,popgreen] at (10.3,2.6) {\textbf{5. Formule de politesse} adaptée au destinataire};
\draw[->,poporange,thick] (10.2,1.2) -- (9.6,1.2);
\node[anchor=west,text width=3.8cm,poporange] at (10.3,1.2) {\textbf{6. Signature} manuscrite};
\end{tikzpicture}
\end{center}
\legende{Photographie schématique d'une lettre de motivation camerounaise annotée : les six zones repérées par des flèches constituent la structure externe, que nous étudierons en détail dans la section suivante.}
\end{popfigure}

\textbf{Les critères de réussite dégagés par l'observation.} La comparaison des deux lettres permet de formuler la « loi » du genre, comme un scientifique tire une conclusion d'une expérience :

\begin{enumerate}
\item \textbf{La forme d'abord} : le lecteur juge la présentation avant le contenu (en-tête, objet, propreté).
\item \textbf{Le respect du destinataire} : vouvoiement, titre exact (« Monsieur le Directeur »), formules de politesse.
\item \textbf{Des arguments prouvés} : chaque qualité affirmée s'appuie sur un fait (stage, diplôme, réalisation).
\item \textbf{La personnalisation} : la lettre montre que le candidat connaît l'entreprise et le poste.
\item \textbf{La langue soignée} : une seule faute peut suffire à écarter un dossier dans une pile de cent candidatures.
\end{enumerate}

\begin{exoresolu}[Identifier le texte fonctionnel]
\textbf{Énoncé.} Pour chaque situation, indique le texte fonctionnel approprié et justifie : a) La mairie de Bafoussam informe les habitants d'une coupure d'eau. b) Aïcha, élève de Première technique, demande un stage d'observation à la SODECOTON. c) La SONARA récapitule les diplômes exigés dans une fiche à remplir par les candidats. d) Brice postule un emploi de soudeur et veut convaincre le chef d'atelier de le recevoir.
\tcblower
\textbf{Solution détaillée.} a) Une \textbf{affiche} (ou un avis) : il s'agit d'informer un large public, sans destinataire unique. b) Une \textbf{lettre de motivation} : Aïcha doit convaincre l'entreprise de l'accueillir, en valorisant sa formation. c) Un \textbf{formulaire} : les informations sont recueillies de façon normalisée pour tous les candidats. d) Une \textbf{lettre de motivation} accompagnée d'un \textbf{CV} : le CV liste ses diplômes (CAP, Probatoire technique), la lettre argumente pour obtenir l'entretien. On retient la règle : le choix du texte dépend de la \emph{fonction} (informer, demander, convaincre) et du \emph{destinataire}.
\end{exoresolu}

\begin{exercice}[À toi de jouer]
1. Classe ces textes en « fonctionnels » ou « littéraires » : une facture CAMWATER ; un poème de René Philombé ; une convocation à un entretien ; un conte peul ; une lettre de demande de bourse au MINESUP.
\par
2. Relis la lettre ratée de Merveille et relève cinq manquements précis aux règles de la lettre officielle.
\par
3. Imagine que tu postules un stage à la SONARA : rédige uniquement la ligne « Objet » et la formule d'appel (« Monsieur... ») correctes.
\end{exercice}

\begin{corrige}[Exercice : Identifier et corriger]
\textbf{1.} Textes fonctionnels : la facture CAMWATER (réclamer un paiement), la convocation (sommer de se présenter), la lettre de demande de bourse (obtenir un acte). Textes littéraires : le poème et le conte (écrits pour l'émotion et le plaisir).
\par
\textbf{2.} Manquements : absence de lieu et date ; absence des coordonnées du destinataire ; absence d'objet ; ton familier (« Salut », « grave ») ; absence de formule de politesse finale ; qualités affirmées sans preuve ; fautes d'orthographe. (Cinq relevés suffisent.)
\par
\textbf{3.} Objet : « candidature à un stage de technicien en raffinage » (une ligne, sans verbe conjugué). Appel : « Monsieur le Directeur des Ressources Humaines, » — jamais « Cher Monsieur » ni « Salut ».
\end{corrige}

\begin{aretenir}[L'essentiel de la section]
Un \cle{texte fonctionnel} sert à agir : demander, informer, convaincre. La \cle{lettre officielle} s'adresse à une autorité avec un ton courtois et une forme codifiée. La \cle{lettre de motivation} est une lettre officielle persuasive : jointe au CV, elle expose les raisons de postuler et valorise les compétences pour décrocher un entretien. Sa réussite repose sur cinq critères : la forme, le respect du destinataire, des arguments prouvés, la personnalisation et une langue irréprochable. Les figures de style (énumération, gradation) et l'impératif de politesse sont des leviers rhétoriques à maîtriser. La section suivante étudiera en détail la \textbf{structure externe} de cette lettre.
\end{aretenir}

\coursec{La structure externe de la lettre officielle}

Après avoir identifié la lettre officielle comme un texte fonctionnel incontournable de la vie administrative et professionnelle, nous allons maintenant disséquer son architecture visible. La structure externe est le squelette sur lequel viendra se greffer le contenu persuasif de votre motivation. Une lettre bien structurée se lit avant même d'être lue : elle signale immédiatement le sérieux du candidat.

\courssub{Les coordonnées de l'expéditeur et du destinataire : l'adresse du dialogue}

\begin{definition}[Coordonnées de l'expéditeur et du destinataire]
Dans une lettre officielle, les coordonnées identifient sans ambiguïté \textbf{l'auteur} (\cle{expéditeur}) et \textbf{le destinataire}. Elles sont placées en en-tête, l'expéditeur à gauche, le destinataire à droite, séparés par un espace vertical suffisant.
\end{definition}

L'expéditeur doit être joignable : nom complet, adresse postale complète (quartier, ville, boîte postale le cas échéant), numéro de téléphone actif et adresse électronique professionnelle (évitez les pseudos fantaisistes). Le destinataire est désigné par sa \textbf{fonction} (Directeur, Proviseur, Maire, DRH), suivie de l'organisme et de la ville. On n'écrit pas le nom propre du destinataire sauf si on est certain de son identité et qu'on s'adresse à lui personnellement (cas rare pour une candidature spontanée).

\begin{popfigure}
\begin{tikzpicture}[scale=0.75, every node/.style={font=\small}]
% Page A4 outline
\draw[popdark, line width=1.5pt] (0,0) rectangle (14,19.5);
% Grid lines light
\draw[popline, step=1cm] (0,0) grid (14,19.5);
% Zones
% Expediteur zone
\fill[popblue!15] (0.5,16.5) rectangle (6.5,19);
\draw[popblue, thick] (0.5,16.5) rectangle (6.5,19);
\node[popblue, align=left, anchor=north west] at (0.7,18.8) {\textbf{EXPÉDITEUR (haut gauche)} \\ Nom Prénom \\ Adresse complète \\ Tél. / E-mail};
% Destinataire zone
\fill[poporange!15] (7.5,16.5) rectangle (13.5,19);
\draw[poporange, thick] (7.5,16.5) rectangle (13.5,19);
\node[poporange, align=left, anchor=north west] at (7.7,18.8) {\textbf{DESTINATAIRE (haut droite)} \\ Fonction \\ Organisme \\ Ville};
% Lieu et Date
\fill[popgreen!15] (0.5,15) rectangle (13.5,16);
\draw[popgreen, thick] (0.5,15) rectangle (13.5,16);
\node[popgreen, align=left, anchor=west] at (0.7,15.5) {\textbf{LIEU ET DATE} \quad Yaoundé, le 15 mars 2026};
% Objet
\fill[poppurple!15] (0.5,13.5) rectangle (13.5,14.5);
\draw[poppurple, thick] (0.5,13.5) rectangle (13.5,14.5);
\node[poppurple, align=left, anchor=west] at (0.7,14) {\textbf{OBJET} \quad Demande de stage académique (réf. : annonce n° 2026/045)};
% Formule d'appel
\fill[poppink!15] (0.5,12.5) rectangle (13.5,13);
\draw[poppink, thick] (0.5,12.5) rectangle (13.5,13);
\node[poppink, align=left, anchor=west] at (0.7,12.75) {\textbf{FORMULE D'APPEL} \quad Monsieur le Directeur,};
% Corps de la lettre (zone grisée)
\fill[popmuted!10] (0.5,5) rectangle (13.5,12);
\draw[popmuted, thick, dashed] (0.5,5) rectangle (13.5,12);
\node[popmuted, align=center] at (7,8.5) {\textbf{CORPS DE LA LETTRE} \\ (Introduction, Développement, Conclusion)};
% Formule de politesse finale
\fill[popgold!15] (0.5,3) rectangle (13.5,4.5);
\draw[popgold, thick] (0.5,3) rectangle (13.5,4.5);
\node[popgold, align=left, anchor=west] at (0.7,3.75) {\textbf{FORMULE DE POLITESSE FINALE} \quad Veuillez agréer, Monsieur le Directeur, l'expression de ma haute considération.};
% Signature
\fill[popblue!15] (7.5,1) rectangle (13.5,2.5);
\draw[popblue, thick] (7.5,1) rectangle (13.5,2.5);
\node[popblue, align=left, anchor=north west] at (7.7,2.3) {\textbf{SIGNATURE} \\ (manuscrite) \\ Nom Prénom};
% Pièces jointes
\fill[poporange!15] (0.5,0.2) rectangle (6.5,0.8);
\draw[poporange, thick] (0.5,0.2) rectangle (6.5,0.8);
\node[poporange, align=left, anchor=west] at (0.7,0.5) {\textbf{PJ} : CV, Lettre de motivation, Diplômes};
\end{tikzpicture}
\legende{Schéma de la structure externe d'une lettre officielle sur page A4 : zones codées par couleur. L'ordre de lecture suit la flèche naturelle : haut gauche $\to$ haut droite $\to$ centre $\to$ bas droite.}
\end{popfigure}

\courssub{Le lieu, la date et l'objet : les repères spatio-temporels et thématiques}

\begin{definition}[Lieu, date et objet]
Le \textbf{lieu et la date} situent la rédaction dans l'espace et le temps (ex. : « Yaoundé, le 15 mars 2026 »). L'\textbf{objet} résume en une ligne le motif de la lettre ; il commence par « Objet : » suivi d'une phrase nominale concise et précise.
\end{definition}

La date s'écrit en toutes lettres pour le mois (pas de chiffre pour le mois) : « Douala, le 3 octobre 2025 ». L'objet doit permettre au destinataire de classer la lettre sans l'ouvrir entièrement. Il mentionne souvent une référence (numéro d'annonce, concours, référence interne).

\begin{attention}[Objet et date : conditions de recevabilité]
Dans de nombreuses administrations camerounaises (fonction publique, collectivités territoriales, établissements publics), l'absence de date ou d'objet constitue un motif de rejet immédiat du courrier. La date prouve l'antériorité de la démarche ; l'objet assure le routage interne. Ne les omettez jamais.
\end{attention}

\begin{exemplebox}[Objets bien formulés]
\begin{itemize}
\item Objet : Candidature au poste d'assistant technique (réf. : avis n° 2026/012/MINESEC)
\item Objet : Demande de stage de fin d'études (Génie civil) — Période : juillet–août 2026
\item Objet : Sollicitation d'une audience pour présentation de projet entrepreneurial
\end{itemize}
\end{exemplebox}

\courssub{La formule d'appel et la formule de politesse finale : le rituel du respect}

\begin{definition}[Formule d'appel et formule de politesse finale]
La \textbf{formule d'appel} (ou salutation initiale) ouvre le dialogue en qualifiant le destinataire selon sa fonction. La \textbf{formule de politesse finale} (ou salutation de clôture) referme la lettre en renouvelant le respect. Elles forment un couple indissociable : la seconde doit reprendre \textbf{exactement} le titre de la première.
\end{definition}

La formule d'appel s'écrit sur une ligne seule, suivie d'une virgule, sans « Cher » ni « Mon » : « Monsieur le Directeur, », « Madame la Mairesse, », « Monsieur le Ministre, ». La formule finale est une phrase complète, souvent longue, qui s'étale sur plusieurs lignes pour aérer la fin de lettre.

\begin{attention}[Registre officiel vs registre familier]
Écrire « Cher Monsieur », « Bonjour Monsieur le Directeur » ou « Salut » dans une lettre officielle est une \textbf{faute de registre grave}. Cela témoigne d'une méconnaissance des codes administratifs et disqualifie souvent le candidat avant même la lecture du corps. La formule d'appel est toujours : \textbf{Titre de fonction} (Monsieur le Directeur, Madame la Proviseure, Monsieur le Chef de Service).
\end{attention}

\begin{savaistu}[La règle de l'écho]
La formule de politesse finale doit \textbf{reproduire à l'identique} le titre utilisé dans la formule d'appel. Si vous avez écrit « Monsieur le Directeur » en ouverture, vous devez fermer par « Veuillez agréer, \textbf{Monsieur le Directeur}, l'expression de ma haute considération. » Toute variation (ex. « Monsieur » seul, « Monsieur le Chef ») est perçue comme une négligence.
\end{savaistu}

\begin{exemplebox}[Formules finales adaptées]
\begin{itemize}
\item \textbf{Appel :} Monsieur le Directeur \quad $\to$ \textbf{Finale :} Veuillez agréer, Monsieur le Directeur, l'expression de ma haute considération.
\item \textbf{Appel :} Madame la Mairesse \quad $\to$ \textbf{Finale :} Je vous prie d'agréer, Madame la Mairesse, l'assurance de ma considération distinguée.
\item \textbf{Appel :} Monsieur le Ministre \quad $\to$ \textbf{Finale :} Recevez, Monsieur le Ministre, l'hommage de ma très haute considération.
\end{itemize}
\end{exemplebox}

\courssub{La signature et les mentions annexes : l'authentification et la complétude}

La signature est \textbf{manuscrite} (encre bleue ou noire), placée en bas à droite, précédée du nom dactylographié. Elle engage la responsabilité de l'auteur. Les mentions annexes (souvent abrégées « PJ » pour « Pièces jointes ») listent les documents accompagnant la lettre : CV, photocopies de diplômes, attestations, lettre de recommandation. Elles sont placées en bas à gauche, après la formule de politesse.

\begin{definition}[Signature et pièces jointes]
La \textbf{signature} est l'apposition manuscrite du nom de l'expéditeur, validant le contenu. Les \textbf{pièces jointes (PJ)} énumèrent les documents annexés ; leur liste doit correspondre exactement au contenu de l'enveloppe.
\end{definition}

\courssub{La mise en page : l'esthétique du sérieux}

Une lettre officielle respecte des normes de présentation qui facilitent la lecture et valorisent le candidat : marges symétriques (2,5 cm minimum), interligne 1,5 pour le corps, alinéas rentrants (ou saut de ligne entre paragraphes), police lisible (Times New Roman, Arial, Calibri, taille 12), alignement justifié. Le support doit être propre, sans ratures, sur papier A4 blanc de bonne qualité (80–90 g/m²).

\begin{popfigure}
\begin{tikzpicture}[scale=0.65, every node/.style={font=\small}]
% Deux miniatures côte à côte
% --- BONNE MISE EN PAGE ---
\begin{scope}[xshift=0cm]
\draw[popgreen, line width=2pt] (0,0) rectangle (6.5,9.5);
% Marge visuelle
\draw[popline, dashed] (0.5,0.5) rectangle (6,9);
% En-tête
\node[popgreen, align=left, anchor=north west] at (0.7,8.8) {\textbf{BONNE MISE EN PAGE}};
\fill[popblue!10] (0.7,7.8) rectangle (3.5,8.6); \node at (2.1,8.2) {Expéditeur};
\fill[poporange!10] (3.7,7.8) rectangle (5.8,8.6); \node at (4.75,8.2) {Destinataire};
\fill[popgreen!10] (0.7,7.2) rectangle (5.8,7.6); \node at (3.25,7.4) {Yaoundé, le 15/03/2026};
\fill[poppurple!10] (0.7,6.6) rectangle (5.8,7.0); \node at (3.25,6.8) {Objet : Stage...};
\fill[poppink!10] (0.7,6.0) rectangle (5.8,6.4); \node at (3.25,6.2) {Monsieur le Directeur,};
% Corps : paragraphes aérés
\fill[popmuted!5] (0.7,4.8) rectangle (5.8,5.6); \node at (3.25,5.2) {Paragraphe 1 (intro)};
\fill[popmuted!5] (0.7,4.0) rectangle (5.8,4.6); \node at (3.25,4.3) {Paragraphe 2 (corps)};
\fill[popmuted!5] (0.7,3.2) rectangle (5.8,3.8); \node at (3.25,3.5) {Paragraphe 3 (concl.)};
% Formule finale
\fill[popgold!10] (0.7,1.8) rectangle (5.8,2.8); \node[align=center] at (3.25,2.3) {Formule de politesse \\ (sur 2-3 lignes)};
% Signature + PJ
\fill[popblue!10] (3.5,0.8) rectangle (5.8,1.5); \node at (4.65,1.15) {Signature};
\fill[poporange!10] (0.7,0.8) rectangle (3.0,1.2); \node at (1.85,1.0) {PJ : CV, Diplômes};
\end{scope}

% --- MAUVAISE MISE EN PAGE ---
\begin{scope}[xshift=7.5cm]
\draw[poppink, line width=2pt] (0,0) rectangle (6.5,9.5);
\node[poppink, align=left, anchor=north west] at (0.7,8.8) {\textbf{MAUVAISE MISE EN PAGE}};
% Tout collé, pas de marges, police fantaisiste simulée par texte dense
\fill[popblue!10] (0.2,8.5) rectangle (3.0,8.9); \node[font=\tiny] at (1.6,8.7) {Expéditeur sans téléphone};
\fill[poporange!10] (3.2,8.5) rectangle (6.3,8.9); \node[font=\tiny] at (4.75,8.7) {Destinataire : "M. le Dir."};
\fill[popgreen!10] (0.2,8.1) rectangle (6.3,8.3); \node[font=\tiny] at (3.25,8.2) {15/03/26 (format chiffre)};
% Objet manquant
\fill[poppink!10] (0.2,7.7) rectangle (6.3,7.9); \node[font=\tiny] at (3.25,7.8) {Monsieur le Directeur, (pas de ligne vide)};
% Corps : bloc compact, pas d'alinéas
\fill[popmuted!20] (0.2,4.5) rectangle (6.3,7.5); \node[align=center, font=\tiny] at (3.25,6.0) {BLOC UNIQUE DENSE \\ sans paragraphes \\ ni aération \\ "Je vous écris pour..."};
% Formule finale bâclée
\fill[popgold!10] (0.2,3.5) rectangle (6.3,4.3); \node[align=center, font=\tiny] at (3.25,3.9) {Cordialement, (trop familier)};
% Signature à gauche, pas de PJ
\fill[popblue!10] (0.2,1.0) rectangle (2.5,1.5); \node[font=\tiny] at (1.35,1.25) {Signature (gauche)};
\end{scope}
\end{tikzpicture}
\legende{Comparaison visuelle : à gauche, une mise en page aérée, structurée, conforme aux normes ; à droite, les erreurs fréquentes (marges insuffisantes, format de date incorrect, objet absent, formule d'appel familière, corps compact, signature mal placée, PJ oubliées).}
\end{popfigure}

\begin{methode}[Rédiger la structure externe : la méthode E-D-D-O-A]
Pour ne rien oublier et respecter l'ordre protocolaire, suivez l'acronyme \textbf{E-D-D-O-A} avant d'écrire le premier mot du corps de la lettre :
\begin{enumerate}
\item \textbf{E} — \textbf{Expéditeur} : vos coordonnées complètes en haut à gauche.
\item \textbf{D} — \textbf{Destinataire} : fonction, organisme, ville en haut à droite.
\item \textbf{D} — \textbf{Date} (et lieu) : « Yaoundé, le 15 mars 2026 » sous le destinataire, aligné à gauche.
\item \textbf{O} — \textbf{Objet} : ligne « Objet : ... » précise et référencée.
\item \textbf{A} — \textbf{Appel} : formule d'appel formelle (« Monsieur le Directeur, ») sur une ligne seule.
\end{enumerate}
\textit{Une fois ces cinq blocs posés, vous pouvez attaquer l'introduction du corps de lettre l'esprit libre.}
\end{methode}

\begin{exoresolu}[Analyse de la structure externe d'une lettre de candidature]
\textbf{Énoncé :} Voici l'en-tête d'une lettre envoyée à la Mairie de Bafoussam pour un poste d'adjoint technique. Identifiez les 5 éléments E-D-D-O-A et signalez toute anomalie.

\begin{quote}
Nguimfack jean-pierre \\
BP 145 - Bafoussam \\
Tel: 6xx xx xx xx \\

Mairie de Bafoussam \\
Service des Ressources Humaines \\
Bafoussam \\

le 12/05/2025 \\

Objet : emploi \\

Cher Monsieur le Maire,
\end{quote}

\textbf{Solution détaillée :}
\begin{enumerate}
\item \textbf{Expéditeur} : Présent mais incomplet (nom en minuscules, pas d'e-mail, adresse imprécise). \textit{Correction :} NGUIMFACK Jean-Pierre, Quartier Banengo, BP 145 Bafoussam, Tel : 6XX XX XX XX, Email : \texttt{jean.nguimfack@email.cm}.
\item \textbf{Destinataire} : Fonction implicite (Maire), organisme correct, ville correcte. \textit{Amélioration :} « Monsieur le Maire, Mairie de Bafoussam, Service des Ressources Humaines, Bafoussam ».
\item \textbf{Date} : Format numérique « 12/05/2025 » \textbf{non conforme}. \textit{Correction :} « Bafoussam, le 12 mai 2025 ».
\item \textbf{Objet} : Trop vague (« emploi »). \textit{Correction :} « Objet : Candidature au poste d'adjoint technique (réf. : avis n° 2025/003/MINAT) ».
\item \textbf{Appel} : « Cher Monsieur le Maire » \textbf{faute de registre} (familier). \textit{Correction :} « Monsieur le Maire, ».
\end{enumerate}
\textbf{Bilan :} 3 anomalies majeures (date, objet, appel) qui risquent de faire classer la lettre sans suite.
\end{exoresolu}

\begin{aretenir}[Check-list structure externe — À valider avant impression]
\begin{itemize}
\item[$\square$] Expéditeur : Nom, Adresse complète, Tél., Email pro (haut gauche)
\item[$\square$] Destinataire : Fonction exacte, Organisme, Ville (haut droite)
\item[$\square$] Lieu + Date en toutes lettres (ex. : Yaoundé, le 15 mars 2026)
\item[$\square$] Objet : Précis, référencé, une ligne
\item[$\square$] Appel : « Monsieur le [Fonction], » ou « Madame la [Fonction], » (sans « Cher »)
\item[$\square$] Formule finale : Reprend \textbf{exactement} le titre de l'appel
\item[$\square$] Signature manuscrite (bas droite) + Nom dactylographié
\item[$\square$] PJ listées (bas gauche) et cohérentes avec l'enveloppe
\item[$\square$] Mise en page : marges 2,5 cm, interligne 1,5, police 12, justifié, papier propre
\end{itemize}
\end{aretenir}

La structure externe est désormais maîtrisée. Elle constitue le cadre rigoureux dans lequel viendra s'insérer le cœur persuasif de votre motivation : l'introduction, le corps argumenté et la conclusion, que nous construirons dans la section suivante.

\coursec{La structure interne : introduction et conclusion}

Après avoir maîtrisé la \emph{structure externe} — l'en-tête, les références, l'objet, la formule d'appel et la formule de politesse de fin — nous pénétrons maintenant au cœur de la lettre : sa \textbf{structure interne}. C'est ici que se joue l'efficacité rhétorique de votre candidature. Une lettre bien présentée mais mal structurée à l'intérieur finit à la corbeille ; une lettre au contenu solide mais mal enveloppée ne passe pas le premier filtre. Cette section vous donne les clés pour construire une introduction qui \emph{accroche} et une conclusion qui \emph{engage}.

\courssub{Les trois parties de la structure interne}

Tout texte argumentatif professionnel, et la lettre de motivation en est l'archétype, repose sur une architecture ternaire classique : \textbf{introduction — corps — conclusion}. Cette tripartition n'est pas arbitraire ; elle suit le cheminement logique de la pensée persuasive : \og~Je capte l'attention et j'annonce mon intention~\/fg (introduction), \og~Je démontre ma valeur~\/fg (corps), \og~J'appelle à l'action et je salue~\/fg (conclusion).

\begin{aretenir}[Structure interne de la lettre officielle]
La lettre de motivation comporte trois séquences obligatoires :
\begin{enumerate}
    \item \textbf{L'introduction} (3 à 5 lignes) : elle situe l'expéditeur, annonce l'objet précis de la demande et \emph{accroche} le lecteur (réponse à une offre, candidature spontanée, recommandation).
    \item \textbf{Le corps de la lettre} (10 à 15 lignes) : il développe l'argumentation en deux temps — \og~Vous~\/fg (l'entreprise, ses besoins, ce qui m'attire) puis \og~Moi~\/fg (mes compétences, mon parcours, ma motivation) — pour aboutir à l'adéquation \og~Nous~\/fg.
    \item \textbf{La conclusion} (2 à 4 lignes) : elle exprime l'espoir d'une suite favorable, propose concrètement un entretien, remercie et précède la formule de politesse.
\end{enumerate}
Chaque partie a une fonction distincte ; les mélanger ou les déséquilibrer affaiblit la démonstration.
\end{aretenir}

\courssub{L'introduction : fonctions et techniques d'accroche}

L'introduction est la \emph{vitrine} de votre lettre. Le recruteur y consacre en moyenne \textbf{6 à 10 secondes}. Elle doit répondre immédiatement à trois questions implicites : \og~Qui écrit ?~\/fg, \og~Que veut-il ?~\/fg, \og~Pourquoi maintenant ?~\/fg.

\textbf{Fonctions de l'introduction.}\par
\begin{itemize}
    \item \textbf{Se présenter brièvement :} statut actuel (élève, étudiant, demandeur d'emploi), formation suivie, établissement. Pas de CV détaillé.
    \item \textbf{Annoncer l'objet :} stage, emploi, inscription, renseignement. L'objet doit être explicite dès la première phrase.
    \item \textbf{Accrocher le lecteur :} créer un lien immédiat entre le destinataire et l'expéditeur. C'est la \emph{proposition de valeur} embryonnaire.
\end{itemize}

\textbf{Les trois techniques d'accroche principales.}\par
Le choix de l'accroche dépend de l'origine de la candidature :

\begin{center}
\begin{tabularx}{\linewidth}{|l|X|X|}\hline
\rowcolor{popblueL}
\textbf{Type de candidature} & \textbf{Accroche type} & \textbf{Exemple concret} \\ \hline
Réponse à une offre & Référence précise à l'annonce (support, date, référence) & \og~Suite à votre annonce parue dans \emph{Cameroon Tribune} du 12 mars 2026, référence \#STG-2026-045...~\/fg \\ \hline
Candidature spontanée & Motivation pour l'entreprise / secteur + veille informationnelle & \og~Suivant avec intérêt le développement de votre service maintenance au sein du groupe \textsc{Sonel}...~\/fg \\ \hline
Sur recommandation & Nom du recommandeur + contexte de la recommandation & \og~M.~Kamga, chef d'atelier au sein de votre structure, m'a suggéré de vous adresser ma candidature...~\/fg \\ \hline
\end{tabularx}
\end{center}

\courssub{Rédaction guidée de l'introduction : le modèle « Suite à votre annonce… »}

Le modèle le plus fréquent en classe de Première technique est la réponse à une offre de stage académique. La phrase-type suivante est un \emph{squelette} qu'il faut adapter, non réciter bêtement.

\begin{methode}[Une introduction efficace répond en trois phrases à « Qui suis-je ? Que demande-je ? Pourquoi maintenant ? »]
\begin{enumerate}
    \item \textbf{Phrase 1 (Qui + Contexte) :} « Actuellement élève en classe de \textbf{Première technique} [spécialité] au \textbf{Lycée technique de [Ville]}, je prépare le \textbf{Probatoire} série [Série]. »
    \item \textbf{Phrase 2 (Objet + Source) :} « \textbf{Suite à votre annonce parue dans [Support] du [Date]} / \textbf{Dans le cadre de mon stage académique obligatoire}, j'ai l'honneur de vous solliciter pour un stage de [durée] au poste de [intitulé]. »
    \item \textbf{Phrase 3 (Accroche / Lien) :} « La réputation de votre entreprise dans le domaine de [secteur] / Vos projets récents en [domaine] suscitent mon vif intérêt et correspondent à mon projet professionnel. »
\end{enumerate}
\textbf{Astuce :} Remplacez les éléments entre crochets par vos informations réelles. La phrase 3 change selon le type d'accroche (voir tableau ci-dessus).
\end{methode}

\begin{exemplebox}[Introduction type : demande de stage académique]
Actuellement élève en classe de Première technique, spécialité Électrotechnique, au Lycée technique de Bafoussam, je me permets de solliciter un stage académique au sein de votre entreprise pour la période de juillet à septembre 2026. Suite à l'avis affiché sur le panneau d'information de l'établissement, j'ai l'honneur de vous présenter ma candidature. Les activités de maintenance industrielle que vous menez, notamment sur les variateurs de fréquence, correspondent pleinement à mon projet professionnel et aux compétences acquises en travaux pratiques.
\end{exemplebox}

\begin{attention}[Piège : l'autobiographie inutile]
Ne commencez \textbf{jamais} par : \og~Je m'appelle Jean Dupont, je suis né le 12 janvier 2008 à Yaoundé, je suis célibataire, j'habite à Mvog-Ada...~\/fg
\begin{itemize}
    \item L'introduction doit tenir en \textbf{3 à 5 lignes maximum}.
    \item L'état civil figure déjà dans l'en-tête (structure externe).
    \item Le recruteur cherche \emph{compétences} et \emph{motivation}, pas votre généalogie.
\end{itemize}
\textbf{Règle d'or :} Toute information qui ne prouve pas votre adéquation au poste n'a pas sa place dans l'introduction.
\end{attention}

\begin{exoresolu}[Réécrire une introduction maladroite]
\textbf{Énoncé :} Voici une introduction rédigée par un élève. Identifiez les défauts et réécrivez-la correctement.
\begin{quote}
\og~Monsieur le Directeur, je m'appelle Paul Biya, j'ai 17 ans, je suis en Première F3 au Lycée de Nkolbisson. Je vous écris parce que je veux un stage. Mes parents sont d'accord. J'aime bien l'électricité. Merci de me répondre.~\/fg
\end{quote}
\tcblower
\textbf{Solution détaillée :}
\begin{itemize}
    \item \textbf{Défauts :} Ton familier (\og~je veux~\/fg), informations inutiles (âge, accord parental), absence d'accroche, absence de référence à l'offre, conclusion prématurée (\og~Merci de me répondre~\/fg).
    \item \textbf{Version corrigée :}
    \og~Monsieur le Directeur,
    Actuellement élève en classe de Première technique, série F3, spécialité Électrotechnique, au Lycée technique de Nkolbisson, j'ai l'honneur de vous adresser ma candidature pour un stage académique de deux mois (juillet--août 2026) au sein de votre service maintenance. Suite à l'annonce diffusée par mon établissement, je suis particulièrement motivé par l'opportunité d'appliquer mes connaissances en automatismes industriels au sein de votre structure reconnue pour son expertise en maintenance préventive.~\/fg
\end{itemize}
\end{exoresolu}

\courssub{La conclusion : fonctions et formules}

Si l'introduction ouvre la porte, la conclusion \emph{invite à entrer}. Elle ne doit pas être un simple salut poli, mais un \textbf{appel à l'action} (call to action) maîtrisé.

\textbf{Fonctions de la conclusion.}\par
\begin{enumerate}
    \item \textbf{Exprimer l'espoir d'une suite favorable :} manifester sa confiance sans arrogance.
    \item \textbf{Proposer un entretien :} passer du virtuel (la lettre) au réel (la rencontre). C'est l'objectif premier.
    \item \textbf{Remercier :} politesse professionnelle pour le temps accordé à la lecture.
    \item \textbf{Signaler les pièces jointes :} CV, bulletins, attestations (souvent omis, pourtant crucial).
\end{enumerate}

\textbf{Formules types (à adapter, pas à copier-coller).}\par
\begin{itemize}
    \item \og~Dans l'attente d'une réponse que j'espère favorable, je me tiens à votre entière disposition pour un entretien à votre convenance.~\/fg (Classique, efficace)
    \item \og~Je serais honoré de vous exposer de vive voix ma motivation lors d'un entretien que vous voudrez bien m'accorder.~\/fg (Plus engageant)
    \item \og~En vous remerciant par avance de l'attention portée à ma candidature, je vous prie d'agréer, Monsieur le Directeur, l'expression de mes salutations distinguées.~\/fg (Fusion conclusion + formule de politesse, gain de place)
\end{itemize}

\begin{attention}[Pièges : conclusion sèche ou suppliante]
\begin{itemize}
    \item \textbf{Trop sèche :} \og~Merci.~\/fg ou \og~Dans l'attente.~\/fg $\rightarrow$ Manque de professionnalisme, donne l'impression de bâcler la fin.
    \item \textbf{Trop suppliante :} \og~Je vous en supplie, donnez-moi ce stage, j'en ai vraiment besoin, aidez-moi.~\/fg $\rightarrow$ Position de faiblesse, manque de dignité, fait fuir le recruteur.
    \item \textbf{Oubli des pièces jointes :} \og~Ci-joint mon CV~\/fg doit apparaître \textbf{avant} la formule de politesse finale.
\end{itemize}
\textbf{Posture attendue :} \emph{Digne, confiant, disponible.} Vous proposez une collaboration, vous ne quémandez pas une aumône.
\end{attention}

\begin{exoresolu}[Réécrire une conclusion abrupte]
\textbf{Énoncé :} Transformez cette conclusion inacceptable en une conclusion professionnelle.
\begin{quote}
\og~J'espère que vous me prendrez. Merci. Paul.~\/fg
\end{quote}
\tcblower
\textbf{Solution détaillée :}
\begin{itemize}
    \item \textbf{Analyse :} Familier (\og~me prendre~\/fg), pas de proposition d'entretien, pas de remerciement formel, signature seule sans formule de politesse.
    \item \textbf{Version corrigée :}
    \og~Dans l'attente d'une réponse que j'espère favorable, je reste à votre entière disposition pour convenir d'un entretien à la date qui vous conviendra le mieux.
    En vous remerciant de l'attention que vous porterez à ma candidature, je vous prie d'agréer, Monsieur le Directeur, l'expression de mes salutations distinguées.
    \textbf{Pièces jointes :} CV détaillé -- Bulletins de Première -- Attestation de scolarité~\/fg
\end{itemize}
\end{exoresolu}

\courssub{Cohérence introduction-conclusion : boucler la demande sans répéter le corps}

L'introduction et la conclusion forment les \textbf{deux couvercles} du même livre. Elles doivent \emph{résonner} ensemble sans se répéter mot pour mot.

\begin{propriete}[Cohérence introduction / conclusion]
\begin{itemize}
    \item \textbf{Objet identique :} Si l'introduction demande \og~un stage de maintenance~\/fg, la conclusion ne parle pas d'\og~emploi d'ingénieur~\/fg.
    \item \textbf{Ton cohérent :} Introduction formelle $\Rightarrow$ Conclusion formelle. Introduction dynamique $\Rightarrow$ Conclusion proactive.
    \item \textbf{Non-répétition du corps :} La conclusion ne résume \textbf{pas} les arguments du corps (compétences, expériences). Elle \textbf{relance} vers l'entretien.
    \item \textbf{Mots-clés écho :} Reprendre un mot fort de l'introduction (ex: \og~maintenance préventive~\/fg) dans la conclusion renforce l'unité du texte.
\end{itemize}
\end{propriete}

\begin{aretenir}[Lien avec l'étude de la langue : Outils rhétoriques et grammaticaux]
La rédaction de l'introduction et de la conclusion mobilise des savoirs linguistiques du programme :
\begin{itemize}
    \item \textbf{Mode impératif (valeur injonctive/incitative) :} La formule de politesse finale (\og~Veuillez agréer...~\/fg, \og~Je vous prie d'agréer...~\/fg) utilise la forme impersonnelle ou la 1\up{re} personne du présent d'indicatif à valeur impérative pour exprimer une demande formelle polie.
    \item \textbf{Figures de style (Gradation, Énumération) :} Dans l'introduction, la \textbf{gradation} (montée en puissance : statut $\rightarrow$ formation $\rightarrow$ projet) structure l'accroche. Dans le corps (voir séquence 4), l'\textbf{énumération} permet de lister compétences et stages.
    \item \textbf{Métonymie :} \og~La direction~\/fg pour \og~les décideurs~\/fg, \og~Le terrain~\/fg pour \og~la réalité professionnelle~\/fg.
\end{itemize}
\end{aretenir}

\courssub{Atelier : réécrire une introduction maladroite et une conclusion abrupte}

\begin{exercice}[Atelier : introduction et conclusion]
\textbf{Consigne :} Pour chacun des deux cas ci-dessous, réécrivez le paragraphe en appliquant la méthode en 3 phrases (introduction) ou les 4 fonctions (conclusion). Respectez le ton professionnel.

\textbf{Cas 1 -- Introduction (Candidature spontanée) :}
\og~Bonjour Monsieur, je m'appelle Amina. Je suis en Première Gestion au Lycée de Douala. Je cherche un stage en secrétariat. Je suis sérieuse et ponctuelle. Mon oncle travaille chez vous.~\/fg

\textbf{Cas 2 -- Conclusion (Réponse à offre) :}
\og~Voilà c'est tout. J'espère que vous allez me prendre. Appellez-moi au 6XX XX XX XX. Au revoir.~\/fg
\end{exercice}

\begin{corrige}[Exercice n1]
\textbf{Cas 1 -- Introduction corrigée :}
\og~Madame, Monsieur,
Actuellement élève en classe de Première technique, option Secrétariat-Bureautique, au Lycée technique de Douala, je me permets de vous adresser une candidature spontanée pour un stage de fin d'année (juin--juillet 2026) au sein de votre direction administrative. La réputation de rigueur de vos services, que M.~Tchinda, mon oncle et collaborateur au sein de votre structure, m'a souvent décrite, renforce ma volonté d'y effectuer mon stage pour y développer mes compétences en gestion administrative et accueil.~\/fg

\textbf{Cas 2 -- Conclusion corrigée :}
\og~Dans l'attente d'une suite favorable à ma candidature, je me tiens à votre disposition pour tout entretien à votre convenance, joignable au 6XX XX XX XX.
En vous remerciant par avance de l'intérêt que vous porterez à ma demande, je vous prie d'agréer, Monsieur le Directeur, l'expression de mes salutations distinguées.
\textbf{Pièces jointes :} CV -- Lettre de motivation -- Bulletins de notes~\/fg
\end{corrige}

\begin{popfigure}
\begin{tikzpicture}[
    node distance=1.2cm and 1.5cm,
    box/.style={draw=popblue, thick, rounded corners, fill=popblue!10, minimum width=3.5cm, minimum height=1.2cm, align=center, font=\small},
    arrow/.style={-Stealth, thick, popdark},
    funnel/.style={draw=poporange, thick, fill=poporange!15, trapezium, trapezium left angle=70, trapezium right angle=110, minimum width=4cm, minimum height=1.5cm, align=center, font=\small}
]
% Introduction Funnel
\node[funnel] (intro_funnel) {INTRODUCTION};
\node[box, above=of intro_funnel, align=center] (hook){ACCROCHE\\Référence offre / Spontané / Recommandation};
\node[box, left=of hook, align=center] (who){QUI SUIS-JE ?\\Statut, Formation, Établissement};
\node[box, right=of hook, align=center] (what){QUE DEMANDE-JE ?\\Objet précis : Stage, Emploi};
\draw[arrow] (who) -- (intro_funnel);
\draw[arrow] (hook) -- (intro_funnel);
\draw[arrow] (what) -- (intro_funnel);

% Conclusion Funnel
\node[funnel, below=2.5cm of intro_funnel] (concl_funnel) {CONCLUSION};
\node[box, below=of concl_funnel, align=center] (hope){ESPOIR\\Réponse favorable attendue};
\node[box, left=of hope, align=center] (avail){DISPONIBILITÉ\\Proposition d'entretien};
\node[box, right=of hope, align=center] (thanks){REMERCIEMENT\\Politesse + Pièces jointes};
\draw[arrow] (concl_funnel) -- (hope);
\draw[arrow] (concl_funnel) -- (avail);
\draw[arrow] (concl_funnel) -- (thanks);

% Link between Intro and Concl
\draw[poppurple, dashed, thick, shorten <=0.5cm, shorten >=0.5cm] (intro_funnel.south) -- (concl_funnel.north) node[midway, right, font=\small, poppurple] {Cohérence : Objet \& Ton};

% Labels
\node[above=0.2cm of hook, font=\bfseries\small, popdark] {Étape 1 : Accrocher};
\node[above=0.2cm of who, font=\bfseries\small, popdark] {Étape 1 : Situer};
\node[above=0.2cm of what, font=\bfseries\small, popdark] {Étape 1 : Annoncer};
\node[below=0.2cm of hope, font=\bfseries\small, popdark] {Étape 3 : Engager};
\node[below=0.2cm of avail, font=\bfseries\small, popdark] {Étape 3 : Proposer};
\node[below=0.2cm of thanks, font=\bfseries\small, popdark] {Étape 3 : Clore};
\end{tikzpicture}
\legende{Schéma en entonnoir : l'introduction canalise trois informations vers l'annonce de la demande ; la conclusion rayonne vers trois actions de clôture. La flèche violette rappelle la cohérence obligatoire entre les deux pôles.}
\end{popfigure}

\begin{savaistu}[Le regard du recruteur]
Des études en psychologie cognitive (eye-tracking) montrent qu'un recruteur passe \textbf{6 à 10 secondes} sur une lettre de motivation lors du premier tri. Son regard suit un \emph{pattern en F} : il lit la première ligne (introduction), saute le corps, lit la dernière ligne (conclusion), puis décide s'il revient au corps. \textbf{Conséquence directe :} soigner l'introduction et la conclusion n'est pas une option esthétique, c'est une \emph{nécessité stratégique}. Une accroche floue ou une conclusion bâclée condamne le corps de lettre, si brillant soit-il, à ne jamais être lu.
\end{savaistu}

\coursec{Le corps de la lettre et les outils langagiers}

Après avoir maîtrisé l'en-tête, l'objet, la formule d'appel, l'introduction et la conclusion (vues en séances précédentes), nous entrons dans le cœur de la lettre de motivation : \textbf{le corps du texte}. C'est ici que se joue votre crédibilité professionnelle. Une introduction accrocheuse et une conclusion polie ne suffisent pas si le développement ne prouve pas, arguments techniques à l'appui, que vous êtes \emph{la} personne qu'il faut. Cette section vous donne la structure logique (le plan « Vous – Moi – Nous »), les outils grammaticaux (mode impératif, figures de style) et les techniques descriptives pour transformer vos expériences techniques en arguments convaincants.

\courssub{Le plan triangulaire : « Vous – Moi – Nous »}

La structure interne du corps de la lettre de motivation répond à une logique persuasive universelle, enseignée des écoles de commerce aux concours administratifs camerounais (ENAM, IRIC). Elle se déploie en trois paragraphes distincts, formant un triangle relationnel.

\begin{popfigure}
\begin{tikzpicture}[
    node distance=2.5cm and 3cm,
    mainnode/.style={draw=popblue, thick, fill=popblueL, rounded corners, minimum width=3.2cm, minimum height=1.8cm, align=center, font=\bfseries},
    arrownode/.style={->, >=Stealth, thick, popdark},
    labelstyle/.style={font=\small, text=popmuted, align=center}
]
% Nœuds principaux
\node[mainnode, align=center] (vous){\textbf{VOUS}\\(L'entreprise)\\\textit{Paragraphe 1}};
\node[mainnode, right=of vous, align=center] (moi){\textbf{MOI}\\(Le candidat)\\\textit{Paragraphe 2}};
\node[mainnode, below=of $(vous)!0.5!(moi)$, align=center] (nous){\textbf{NOUS}\\(La collaboration)\\\textit{Paragraphe 3}};

% Flèches
\draw[arrownode] (vous) -- node[above, labelstyle] {Connaissance \& Intérêt} (moi);
\draw[arrownode] (vous) -- node[left, labelstyle] {Adéquation} (nous);
\draw[arrownode] (moi) -- node[right, labelstyle] {Apport} (nous);

% Légendes détaillées
\node[labelstyle, below=0.3cm of vous, align=center]{Montrez que vous connaissez\\ses valeurs, ses projets, ses besoins};
\node[labelstyle, below=0.3cm of moi, align=center]{Preuvez vos compétences techniques\\(diplômes, stages, réalisations)};
\node[labelstyle, below=0.3cm of nous, align=center]{Projetez la plus-value mutuelle\\et l'entretien};
\end{tikzpicture}
\legende{Le triangle « Vous – Moi – Nous » : architecture argumentative du corps de la lettre.}
\end{popfigure}

\begin{methode}[Rédiger le corps selon le plan « Vous – Moi – Nous »]
\begin{enumerate}
    \item \textbf{Paragraphe 1 « VOUS » (L'entreprise) :} Démontrer que vous ne postulez pas au hasard. Citez un projet précis, une valeur, une actualité de l'entreprise. \textit{Exemple : « Votre récente implantation dans la zone industrielle de Bassa et votre politique de formation continue des techniciens en électromécanique retiennent toute mon attention. »}
    \item \textbf{Paragraphe 2 « MOI » (Le candidat) :} Présenter votre profil technique en miroir des besoins identifiés. Diplômes (Probatoire, Baccalauréat technique), stages, projets concrets, savoir-faire (lecture de plans, maintenance, soudure, programmation PLC, etc.). Utilisez l'\textbf{énumération} et la \textbf{gradation}.
    \item \textbf{Paragraphe 3 « NOUS » (La collaboration) :} Projeter l'avenir. Que ferez-vous \emph{ensemble} ? Apport immédiat, volonté d'apprendre, disponibilité pour l'entretien. Terminez sur une note d'engagement.
\end{enumerate}
\end{methode}

\begin{attention}[Piège fréquent : l'inversion ou la confusion]
Ne commencez jamais par « Moi » (mon CV, mes notes). L'employeur lit d'abord pour \emph{ses} besoins. Évitez aussi le paragraphe unique « fourre-tout » où se mélangent connaissances de l'entreprise, liste de diplômes et demande d'entretien. La clarté du plan en trois paragraphes distincts est un signal de rigueur professionnelle.
\end{attention}

\courssub{Valoriser ses compétences techniques : du diplôme à la preuve}

En Première technique, votre atout majeur est votre \cle{formation professionnalisante}. Le recruteur attend des preuves concrètes, pas des adjectifs vagues.

\begin{definition}[Compétence technique vs Qualité personnelle]
Une \textbf{compétence technique} est un savoir-faire vérifiable (ex. : « Réaliser une maintenance préventive sur moteur asynchrone triphasé »). Une \textbf{qualité personnelle} (soft skill) est une attitude (ex. : « Rigueur », « Ponctualité »). Dans le corps de la lettre, \textbf{ancrez chaque qualité dans une compétence technique}.
\end{definition}

\begin{exemplebox}[Transformer une liste brute en argumentaire technique]
\textbf{Liste brute :} Bac Tech Électrotechnique. Stage à ENEO. Ponctuel. Rigoureux. Connaît les schémas.
\textbf{Paragraphe « MOI » rédigé :}
« Titulaire du Baccalauréat technique en Électrotechnique (session 2024, mention Assez Bien), j'ai consolidé mes acquis lors d'un stage de trois mois à la direction régionale d'ENEO à Douala-Bassa. \textbf{Rigoureux} dans le respect des normes de sécurité NF C 15-100, \textbf{ponctuel} sur les chantiers, j'ai participé à la maintenance préventive de postes de transformation 30/0,4 kV. J'assure la lecture de schémas unifilaires et le câblage d'armoires de commande sous supervision. »
\end{exemplebox}

\begin{aretenir}[La règle du « Preuve = Contexte + Action + Résultat »]
Pour chaque compétence citée, donnez : le \textbf{contexte} (stage, TP, projet), l'\textbf{action} technique précise (verbe d'action : câbler, programmer, diagnostiquer, souder, mesurer) et si possible le \textbf{résultat} (gain de temps, conformité, satisfaction client).
\end{aretenir}

\courssub{L'argumentation : relier qualité et besoin de l'employeur}

Argumenter, ce n'est pas lister ; c'est \textbf{relier}. Utilisez des connecteurs logiques pour créer la causalité : \emph{car, puisque, en effet, c'est pourquoi, ainsi, par conséquent}.

\begin{exoresolu}[De l'affirmation à l'argumentation]
\textbf{Énoncé :} Vous postulez chez CAMWATER pour un poste de technicien en maintenance hydraulique. L'offre mentionne : « Urgence : fuites sur canalisations fonte, travail de nuit fréquent ». Vous avez fait un stage de réparation de fuites sous pression. Rédigez l'argument.
\tcblower
\textbf{Solution :}
« \textbf{Conscient de l'urgence que représentent les fuites sur votre réseau en fonte} (VOUS), \textbf{j'ai acquis lors de mon stage chez CAMWATER Yaoundé une expérience concrète de la réparation sous pression sans interruption de service} (MOI). \textbf{Aussi, je suis immédiatement opérationnel pour intervenir en astreinte nocturne, garantissant la continuité de la distribution d'eau potable} (NOUS). »
\end{exoresolu}

\courssub{Le mode impératif : valeurs, emplois et l'impératif atténué}

L'impératif est le mode de l'injonction. Dans une lettre officielle, il apparaît presque exclusivement dans les \textbf{formules de politesse finales} et parfois dans l'objet. Maîtriser ses valeurs évite les maladresses de registre.

\begin{definition}[Les valeurs de l'impératif]
L'impératif exprime quatre valeurs principales selon le contexte et l'intention du locuteur :
\begin{itemize}
    \item \textbf{L'ordre / L'injonction :} Autorité, contrainte. \textit{Ex : « Fermez la porte. »}
    \item \textbf{Le conseil / La recommandation :} Bienveillance, expertise. \textit{Ex : « Prenez le temps de relire votre lettre. »}
    \item \textbf{La prière / La supplication :} Humilité, demande pressante. \textit{Ex : « Pardonnez-moi. »}
    \item \textbf{Le souhait / La bénédiction :} Formule figée. \textit{Ex : « Vive la République ! »}
\end{itemize}
\end{definition}

\begin{popfigure}
\begin{tikzpicture}[
    grow'=right,
    level distance=4.5cm,
    sibling distance=2.2cm,
    edge from parent path={(\tikzparentnode.east) -- ++(1cm,0) |- (\tikzchildnode.west)},
    every node/.style={draw=popblue, thick, fill=popblueL, rounded corners, minimum width=2.8cm, minimum height=1.2cm, align=center, font=\small},
    root/.style={fill=popblue, text=white, font=\bfseries\normalsize, minimum width=3.5cm},
    level 1/.style={sibling distance=2.5cm},
]
\node[root] {IMPÉRATIF}
    child { node[align=center]{\textbf{ORDRE / INJONCTION}\\\textit{Autorité, contrainte}\\« Respectez les consignes. » } }
    child { node[align=center]{\textbf{CONSEIL / RECOMMANDATION}\\\textit{Bienveillance, expertise}\\« Vérifiez le câblage avant mise sous tension. » } }
    child { node[align=center]{\textbf{PRIÈRE / SUPPLICATION}\\\textit{Humilité, politesse}\\« Veuillez agréer, Madame, mes salutations. » } }
    child { node[align=center]{\textbf{SOUHAIT / BÉNÉDICTION}\\\textit{Formule figée}\\« Vive le Cameroun ! » } };
\end{tikzpicture}
\legende{Arbre des valeurs de l'impératif avec exemples contextuels.}
\end{popfigure}

\begin{propriete}[L'impératif atténué : la marque du registre administratif]
Dans la correspondance officielle, l'impératif brut (« Agréez... ») est perçu comme trop direct, voire impoli. On utilise l'\textbf{impératif atténué} par le verbe \textbf{veuillez} (2\textsuperscript{e} personne du pluriel de \emph{vouloir} à l'impératif) suivi de l'infinitif.
\[
\textbf{Veuillez} + \text{infinitif} \quad \rightarrow \quad \textit{Veuillez agréer, Veuillez trouver ci-joint, Veuillez recevoir.}
\]
Cette construction transforme l'ordre en \textbf{prière polie}, standard incontournable des formules de clôture.
\end{propriete}

\begin{attention}[Erreurs fatales à l'impératif dans une lettre]
\begin{itemize}
    \item \textbf{Utiliser le « tu » :} « Veuillez » (vous) est obligatoire. « Veux agréer » n'existe pas.
    \item \textbf{Oublier le « s » à la 1\textsuperscript{re} personne du pluriel (nous) :} « \textbf{Permettez-nous} de vous présenter... » (et non *Permettez-nous).
    \item \textbf{Confondre impératif et indicatif :} « Je vous prie d'agréer » (indicatif) vs « Veuillez agréer » (impératif). Les deux sont corrects mais l'impératif « Veuillez » est la norme administrative camerounaise.
\end{itemize}
\end{attention}

\courssub{Les figures de style : outils de précision et de persuasion}

Les figures de style ne sont pas réservées à la littérature. Dans une lettre de motivation, bien dosées, elles structurent l'information, renforcent l'image professionnelle et marquent l'esprit du recruteur. Quatre figures sont au programme : \textbf{hyperbole, métonymie, gradation, énumération}.

\begin{popfigure}
\begin{center}
\begin{tabularx}{\linewidth}{|l|X|X|}\hline
\rowcolor{popblueL}
\textbf{Figure} & \textbf{Définition \& Mécanisme} & \textbf{Exemple en lettre de motivation (Contexte : Technicien Maintenance)} \\ \hline
\textbf{\cle{Énumération}} & Accumulation de termes (mots ou groupes) de même nature syntaxique, sans conjonction (asyndète) ou avec (polysyndète). \textit{Effet : exhaustivité, densité, rythme.} & « \textbf{Autonome, rigoureux, réactif et respectueux des normes HSE}, j'interviens sur variateurs de fréquence et automates Siemens. » \\ \hline
\textbf{\cle{Gradation}} & Énumération \textbf{ordonnée} selon une intensité croissante (ou décroissante). \textit{Effet : montée en puissance, climax.} & « De la \textbf{lecture de plans} au \textbf{câblage d'armoires}, jusqu'à la \textbf{mise en service complète} de lignes automatisées, je maîtrise la chaîne. » \\ \hline
\textbf{\cle{Métonymie}} & Désignation d'une réalité par une autre qui lui est liée (contenant/contenu, auteur/œuvre, lieu/activité, matière/objet, signe/chose signifiée). \textit{Effet : concision, image concrète.} & « J'ai \textbf{porté le casque} sur le chantier de la centrale de Memve'ele. » (Le casque = la fonction de technicien sur le terrain / la sécurité). \\ \hline
\textbf{\cle{Hyperbole}} & Exagération volontaire, outrancière, pour frapper l'esprit. \textit{Effet : emphase, émotion. \textbf{À utiliser avec une extrême prudence en lettre formelle.}} & \textit{À éviter :} « Je suis \textbf{le meilleur soudeur du Cameroun}. » \\ \textit{Acceptable (mesuré) :} « J'ai une \textbf{soif intarissable} d'apprendre vos procédés spécifiques. » \\ \hline
\end{tabularx}
\end{center}
\legende{Tableau récapitulatif des quatre figures de style au programme avec exemples contextualisés.}
\end{popfigure}

\begin{definition}[Gradation vs Énumération simple]
L'\textbf{énumération} est une liste (plate). La \textbf{gradation} est une liste \textbf{hiérarchisée} (en escalier). En lettre de motivation, la gradation prouve une \textbf{progression logique des compétences} (théorique $\rightarrow$ pratique $\rightarrow$ autonomie) ou une \textbf{montée de l'engagement} (intérêt $\rightarrow$ compétence $\rightarrow$ engagement total).
\end{definition}

\begin{exemplebox}[Énumération et gradation combinées dans le paragraphe « MOI »]
« Fort de mon \textbf{Bac Tech Maintenance Industrielle} (diplôme), j'ai \textbf{diagnostiqué} des pannes hydrauliques, \textbf{remplacé} des vérins et \textbf{optimisé} le cycle de graissage lors de mon stage à la CDC Tiko (gradation action : simple $\rightarrow$ complexe $\rightarrow$ amélioration). \textbf{Curieux, méthodique, force de proposition} (énumération de qualités ancrées), je suis prêt à m'investir durablement. »
\end{exemplebox}

\begin{attention}[L'hyperbole : danger zone]
L'hyperbole est \textbf{formellement déconseillée} pour vanter vos compétences techniques. Elle détruit la \cle{crédibilité}, valeur cardinale du technicien.
\begin{itemize}
    \item \textbf{Interdit :} « Je suis un génie de la programmation. », « Je résous tous les problèmes en 5 minutes. », « Ma rigueur est absolue. »
    \item \textbf{Préférez le factuel :} « J'ai programmé 3 automates Siemens S7-1200 lors de mon projet de fin d'études. », « J'ai réduit le temps d'arrêt machine de 15\% sur la ligne 3. »
\end{itemize}
Une hyperbole \emph{mesurée} sur la \textbf{motivation} ou la \textbf{soif d'apprendre} peut passer (« Une motivation sans faille », « Une détermination à toute épreuve »), mais restez dans le registre de l'engagement, pas de la compétence magique.
\end{attention}

\begin{savaistu}[Le plan « Vous – Moi – Nous » : standard international]
Ce plan tripartite n'est pas une invention locale. Il est enseigné dans les \textbf{écoles de commerce (HEC, ESSEC, HEC Montréal)}, les \textbf{écoles d'ingénieurs} et les \textbf{centres de formation aux concours administratifs} (ENAM, IRIC au Cameroun). Il structure aussi les \textbf{lettres de motivation pour les visas d'études}, les \textbf{demandes de bourses} (Mastercard Foundation, DAAD, Chevening) et les \textbf{pitchs entrepreneuriaux}. Le maîtriser en Première technique, c'est acquérir une compétence transverse pour la vie professionnelle.
\end{savaistu}

\courssub{Le texte descriptif : décrire son parcours avec précision}

Décrire, ce n'est pas « faire joli » ; c'est \textbf{faire voir} la réalité technique. Le texte descriptif dans le paragraphe « MOI » mobilise : \textbf{adjectifs qualificatifs techniques}, \textbf{expansions du nom} (prépositionnelles, relatives), \textbf{vocabulaire précis} (terminologie du métier).

\begin{definition}[Le texte descriptif technique]
Il vise à \textbf{représenter fidèlement} un objet, un lieu, un processus ou un parcours. Il s'appuie sur :
\begin{itemize}
    \item \textbf{La nomination précise :} pas « machine » mais « tour à commande numérique Fanuc 0i-TF » ; pas « câbler » mais « réaliser le câblage d'armoires BT selon nomenclature IEC 61439 ».
    \item \textbf{L'expansion du nom :} GN + \textit{de} + GN (compétence \textbf{de} lecture \textbf{de} schémas), GN + \textit{à} + Infinitif (aptitude \textbf{à} programmer), GN + Proposition relative (stage \textbf{que j'ai effectué} chez...).
    \item \textbf{Les adjectifs qualificatifs techniques :} \textit{préventive, curative, prédictive, triphasé, asynchrone, programmable, normée, certifiée.}
\end{itemize}
\end{definition}

\begin{methode}[Atelier : Transformer une liste en paragraphe descriptif argumenté]
\textbf{Consigne :} À partir de la liste brute ci-dessous, rédiger le paragraphe « MOI » (80-100 mots) en intégrant \textbf{une énumération} de qualités et \textbf{une gradation} de compétences.
\textbf{Liste brute :}
\begin{itemize}
    \item Diplôme : Bac Tech Froid et Climatisation (2024).
    \item Stage : 3 mois chez SOCAR (station service) + 1 mois chez particulier (clim split).
    \item Compétences : Récupération fluides (R134a, R410A), brasage fort, recherche fuite, mise en service, respect norme environnement.
    \item Qualités : Ponctuel, soigneux, autonome, bon relationnel client.
\end{itemize}
\textbf{Production attendue (modèle) :}
« Titulaire du \textbf{Baccalauréat technique Froid et Climatisation} (session 2024), j'ai affiné mon savoir-faire lors d'un \textbf{stage de trois mois chez SOCAR} (maintenance centrale de traitement d'air) \textbf{et d'un mois en intervention résidentielle} (gradation : industriel $\rightarrow$ tertiaire/particulier). \textbf{Certifié manipulation fluides frigorigènes} (R134a, R410A), je maîtrise le \textbf{brasage fort à l'argent}, la \textbf{recherche de fuite au détecteur électronique} et la \textbf{mise en service conforme aux normes environnementales} (gradation technique). \textbf{Ponctuel, soigneux, autonome et à l'écoute du client} (énumération de qualités professionnelles), j'assure un service rigoureux et durable. »
\end{methode}

\begin{exoresolu}[Analyse linguistique du paragraphe modèle]
\begin{itemize}
    \item \textbf{Gradation 1 (Contextes) :} SOCAR (industriel, complexe) $\rightarrow$ Particulier (autonomie, relationnel).
    \item \textbf{Gradation 2 (Compétences) :} Certification (base légale) $\rightarrow$ Brasage (savoir-faire manuel pointu) $\rightarrow$ Détection fuite (diagnostic) $\rightarrow$ Mise en service (aboutissement global).
    \item \textbf{Énumération (Qualités) :} 4 adjectifs qualificatifs épithètes du sujet « je », asyndète (sans « et » avant le dernier), ordre : attitude temps $\rightarrow$ attitude travail $\rightarrow$ autonomie $\rightarrow$ relationnel.
    \item \textbf{Expansions du nom :} « stage \textbf{de} trois mois \textbf{chez} SOCAR », « maintenance \textbf{de} centrale \textbf{de} traitement \textbf{d'}air », « certifié \textbf{manipulation} fluides », « brasage \textbf{fort} \textbf{à} l'argent ».
    \item \textbf{Vocabulaire technique :} R134a, R410a, brasage fort, détecteur électronique, normes environnementales.
\end{itemize}
\end{exoresolu}

\courssub{Atelier d'écriture guidé : « Mon paragraphe MOI »}

\begin{exercice}[Rédaction du paragraphe « MOI » — 20 min]
\textbf{Contexte :} Vous êtes en Terminale Technique (spécialité au choix : Électrotechnique, Mécanique, BTP, Informatique, Secrétariat, etc.). Vous répondez à une offre de stage de fin d'études ou d'emploi junior.
\textbf{Consignes :}
\begin{enumerate}
    \item Remplissez votre « Fiche profil » (5 min) :
    \begin{itemize}
        \item Diplôme exact + année + mention (si connue) :
        \item Stages / Projets majeurs (lieu, durée, 2-3 tâches techniques précises) :
        \item 3 Compétences techniques « phares » (vocabulaire du métier) :
        \item 4 Qualités professionnelles (adjectifs) :
    \end{itemize}
    \item Rédigez le paragraphe « MOI » (10-12 lignes) en appliquant :
    \begin{itemize}
        \item Une \textbf{gradation} (sur les contextes OU sur les compétences).
        \item Une \textbf{énumération} (sur les qualités).
        \item Au moins \textbf{3 expansions du nom} précises (prépositionnelles ou relatives).
        \item Zéro hyperbole sur les compétences.
    \end{itemize}
    \item Auto-évaluation (cochez) :
    \begin{itemize}
        \item [ ] Gradation visible (mots-outils : \textit{d'abord, ensuite, enfin / de... à... / du... au...})
        \item [ ] Énumération visible (liste d'adjectifs ou de noms)
        \item [ ] Vocabulaire technique précis (pas de « machine », « truc », « bidule »)
        \item [ ] Ton professionnel (pas de familier, pas d'argot)
    \end{itemize}
\end{enumerate}
\end{exercice}

\begin{corrige}[Exercice 1]
\textbf{Critères de réussite (sur 10 pts) :}
\begin{tabularx}{\linewidth}{|l|X|c|}\hline
\rowcolor{popblueL}
\textbf{Critère} & \textbf{Indicateurs observables} & \textbf{Pts} \\ \hline
Structure & Paragraphe unique, cohérent, 8-12 lignes & 1 \\ \hline
Gradation & 3 termes minimum, ordre logique (croissant), connecteurs adaptés & 2 \\ \hline
Énumération & 3-4 qualités, syntaxe correcte (asyndète ou polysyndète) & 1,5 \\ \hline
Précision technique & 3 termes spécialisés minimum, expansions du nom variées & 2,5 \\ \hline
Registre & Formel, impersonnel si besoin, pas d'hyperbole, pas de fautes graves & 2 \\ \hline
Orthographe / Grammaire & Accords, conjugaison (impératif si formule), ponctuation & 1 \\ \hline
\end{tabularx}
\textbf{Exemple de correction type (Spécialité Électrotechnique) :}
« Titulaire du \textbf{Bac Tech Électrotechnique} (2024), j'ai consolidé mes bases en \textbf{maintenance préventive de moteurs asynchrones} lors de mon stage à \textbf{SONEL/ENEO Douala-Bassa} (3 mois), puis en \textbf{programmation d'automates Siemens S7-1200} lors de mon projet de fin d'études (gradation : maintenance $\rightarrow$ automatisme). \textbf{Habilité BT/HT, rigoureux, autonome et respectueux des consignes HSE} (énumération), je suis opérationnel pour intégrer votre service maintenance. »
\end{corrige}

\begin{aretenir}[Synthèse : Les 3 piliers du corps de lettre réussi]
\begin{enumerate}
    \item \textbf{Architecture :} Respect strict du plan \textbf{Vous – Moi – Nous} (3 paragraphes distincts).
    \item \textbf{Preuve technique :} Vocabulaire précis, expansions du nom, gradation des compétences (diplôme $\rightarrow$ stage $\rightarrow$ autonomie), énumération des qualités ancrées.
    \item \textbf{Registre maîtrisé :} Impératif atténué (\textit{Veuillez agréer...}) pour la politesse, figures de style \textbf{mesurées} (gradation, énumération), \textbf{zéro hyperbole} sur les compétences.
\end{enumerate}
\end{aretenir}

\coursec{La communication par l'image dans la candidature}

Après avoir maîtrisé la rédaction du corps de la lettre — arguments, vocabulaire précis, mode impératif pour l'appel à l'action — il reste une dimension souvent sous-estimée mais décisive : \textbf{l'image}. Un dossier de candidature n'est pas seulement un texte ; c'est un \emph{objet visuel} que le recruteur « lit » en quelques secondes avant même de déchiffrer la première phrase. Cette section explore comment l'image (photo, mise en page, logo) construit votre crédibilité professionnelle, conformément à la compétence « Communication par l'image » du programme de Première technique.

\courssub{Qu'est-ce que la communication par l'image ?}

La communication par l'image désigne l'ensemble des processus par lesquels des éléments visuels (photographies, schémas, mise en page, couleurs, typographie) transmettent un message, intentionnel ou non, à un destinataire. Dans un dossier de candidature, elle remplit trois fonctions majeures :

\begin{itemize}
    \item \textbf{Informer :} la photo d'identité permet l'identification physique ; le logo de l'entreprise sur l'en-tête authentifie la provenance de l'offre ; la structure visuelle (rubriques, puces) guide la lecture du CV.
    \item \textbf{Séduire :} une mise en page harmonieuse, une photo soignée, un papier de qualité suscitent l'intérêt esthétique et l'envie de lire la suite.
    \item \textbf{Convaincre :} la cohérence entre l'image renvoyée (sérieux, dynamisme, rigueur) et le discours tenu dans la lettre renforce l'argumentation : « Je suis organisé » \emph{et} mon dossier en a l'air.
\end{itemize}

\begin{definition}[Dénotation et connotation de l'image]
La \cle{dénotation} est le sens littéral, objectif, descriptif d'une image : ce que tout observateur voit (ex. : « un jeune homme en chemise bleue, fond blanc, cadré buste »). La \cle{connotation} est l'ensemble des idées, valeurs, émotions, jugements que l'image évoque subjectivement selon le contexte culturel (ex. : « ce jeune homme a l'air sérieux, compétent, prêt à travailler »).
\end{definition}

\begin{popfigure}
\begin{tikzpicture}[
    node distance=2.5cm and 3cm,
    main/.style={draw=popblue, thick, rounded corners, fill=popblue!10, minimum width=3cm, minimum height=1.5cm, align=center},
    arrow/.style={-Stealth, thick, popdark},
    label/.style={midway, fill=white, inner sep=2pt, font=\small, text=popdark}
]
\node[main, align=center] (img){IMAGE\\ (Support visuel)};
\node[main, below left=of img, align=center] (den){DÉNOTATION\\ \textit{Description objective}\\(Ce qui est montré)};
\node[main, below right=of img, align=center] (con){CONNOTATION\\ \textit{Interprétation subjective}\\(Ce qui est suggéré)};
\node[main, below=3.5cm of img, align=center] (msg){MESSAGE\\ \textit{Sens global perçu}\\(Identité professionnelle)};
\draw[arrow] (img) -- node[label, left] {analyse} (den);
\draw[arrow] (img) -- node[label, right] {interprète} (con);
\draw[arrow] (den) -- node[label, left] {alimente} (msg);
\draw[arrow] (con) -- node[label, right] {alimente} (msg);
\draw[arrow, dashed] (den) -- (con) node[label, below] {interaction};
\end{tikzpicture}
\legende{Triangle sémiotique appliqué à l'image de candidature : l'image source génère à la fois une lecture littérale (dénotation) et une lecture culturelle (connotation) qui fusionnent pour former le message perçu par le recruteur.}
\end{popfigure}

\begin{methode}[Analyser une image de candidature (démarche en 3 temps)]
\begin{enumerate}
    \item \textbf{Observer / Décrire (Dénotation) :} Que vois-je exactement ? Couleurs, cadrage, lumière, vêtements, posture, police de caractères, marges, alignements, présence/absence de logo. \emph{Ne pas juger, seulement décrire.}
    \item \textbf{Interpréter (Connotation) :} Que cela évoque-t-il dans le monde professionnel camerounais ? (Ex. : chemise bien repassée $\rightarrow$ rigueur ; police Comic Sans $\rightarrow$ amateurisme ; photo de groupe recadrée $\rightarrow$ manque de préparation).
    \item \textbf{Évaluer la cohérence :} L'image sert-elle le projet professionnel affiché dans la lettre ? Y a-t-il contradiction ou renforcement ?
\end{enumerate}
\end{methode}

\courssub{La photo d'identité : codes vestimentaires, cadrage, neutralité}

La photo n'est pas une formalité administrative : c'est la \textbf{première poignée de main visuelle}. Au Cameroun, comme dans la plupart des contextes professionnels francophones, des codes implicites s'imposent.

\begin{exemplebox}[Photo « réussie » vs Photo « ratée »]
\begin{center}
\begin{tabularx}{\linewidth}{|l|X|X|}
\hline
\rowcolor{popblueL} \textbf{Critère} & \textbf{Photo soignée (recommandée)} & \textbf{Photo négligée (à éviter)} \\ \hline
\textbf{Tenue} & Chemise/blouse unie (blanc, bleu clair, beige), veste optionnelle, cravate discrète si homme. & T-shirt, polo à motifs, débardeur, sweat-shirt, tenue de sport, tenue traditionnelle (sauf poste spécifique). \\ \hline
\textbf{Cadrage} & Plan buste (épaules + visage), centré, regard vers l'objectif. & Plan pied, selfie bras tendu, photo de groupe découpée, cadrage penché. \\ \hline
\textbf{Fond} & Uni, clair (blanc, gris clair, bleu très pâle), sans ombre portée. & Fond chargé (rideau, mur abîmé, paysage, voiture, bureau en désordre). \\ \hline
\textbf{Éclairage} & Lumière douce, visage bien éclairé, pas d'ombre dure, pas de reflet sur lunettes. & Contre-jour, flash direct (front brillant), ombre du nez sur la joue, pénombre. \\ \hline
\textbf{Expression} & Neutre bienveillant, léger sourire possible, yeux ouverts, regard franc. & Grimace, rire large, air triste/ennuyé, lunettes de soleil, casquette, filtre Snapchat. \\ \hline
\textbf{Qualité technique} & Haute résolution (300 dpi), format 3,5 $\times$ 4,5 cm, nette, non pixelisée. & Floue, pixelisée, recadrée n'importe comment, format non standard. \\ \hline
\end{tabularx}
\end{center}
\legende{Tableau comparatif des standards de la photo d'identité professionnelle au Cameroun.}
\end{exemplebox}

\begin{attention}[L'image contredit le discours]
Une lettre qui vante votre « \emph{rigueur}, \emph{sens de l'organisation} et \emph{respect des normes} » perd toute crédibilité si la photo montre un candidat en t-shirt froissé, mal cadré, avec un fond de chambre en désordre. \textbf{L'image doit servir la lettre, pas la contredire.} Le recruteur croira l'image (preuve visuelle immédiate) avant le texte (promesse verbale).
\end{attention}

\courssub{La mise en page comme image : police, alignements, espaces, papier}

La mise en page \emph{est} une image. Elle parle avant que l'on lise le premier mot. Elle signale : « Je maîtrise les outils bureautiques », « Je respecte les codes de l'écrit professionnel », « Je facilite votre travail de lecture ».

\begin{propriete}[Les piliers d'une mise en page professionnelle]
\begin{itemize}
    \item \textbf{Police (fonte) :} Une seule police classique à empattements (\emph{serif}) pour le corps (Times New Roman, Garamond, Georgia) \textbf{ou} une police sans empattements (\emph{sans serif}) lisible (Arial, Calibri, Helvetica). Taille 11 ou 12 pt. \textbf{Interdit :} polices fantaisistes (Comic Sans, Papyrus, Brush Script), polices manuscrites, plus de deux polices différentes.
    \item \textbf{Alignements :} Texte justifié (bloc propre) ou aligné à gauche (plus facile à lire à l'écran). \textbf{Jamais} centré pour le corps du texte. En-tête (coordonnées) : aligné à gauche ou centré avec cohérence.
    \item \textbf{Espaces (blancs) :} Marges 2,5 cm minimum. Interligne 1,15 ou 1,5. Espace entre paragraphes (6 à 12 pt). « L'espace blanc n'est pas du vide, c'est de la respiration pour l'œil. »
    \item \textbf{Papier / Support numérique :} Pour un dépôt physique : papier blanc 90--100 g/m², non recyclé (trop granuleux), impression laser nette. Pour un PDF : nom de fichier clair (\texttt{Nom\_Prenom\_LettreMotivation\_Entreprise.pdf}), taille < 1 Mo, police intégrée.
\end{itemize}
\end{propriete}

\begin{popfigure}
\begin{tikzpicture}
% Vignette 1 : Mise en page aérée (GAUCHE)
\begin{scope}[xshift=0cm]
\node[draw=popgreen, thick, rounded corners, fill=popgreen!5, minimum width=7.5cm, minimum height=10cm] (box1) {};
\node[font=\bfseries\small, text=popgreen] at ([yshift=-0.4cm]box1.north) {Mise en page AÉRÉE (Recommandée)};
% En-tête
\fill[popblue!30] (box1.north west) ++(0.3,-0.3) rectangle ++(6.9,1.2);
\node[font=\tiny, align=left, text=popdark] at ([shift={(3.5,-0.9)}]box1.north west) {NOM Prénom \\ Adresse $\cdot$ Tél $\cdot$ Email};
% Date / Destinataire
\fill[popmuted!20] (box1.north west) ++(0.3,-1.8) rectangle ++(6.9,1.0);
\node[font=\tiny, align=left, text=popdark] at ([shift={(3.5,-2.3)}]box1.north west) {Douala, le 15/10/2025 \\ À l'attention de M. le DRH \\ Entreprise X};
% Objet
\fill[poporange!30] (box1.north west) ++(0.3,-3.1) rectangle ++(6.9,0.5);
\node[font=\tiny, text=popdark] at ([shift={(3.5,-3.35)}]box1.north west) {Objet : Candidature au poste de Technicien Maintenance};
% Corps (lignes simulées)
\foreach \y in {4.0, 4.5, 5.0, 5.5, 6.0, 6.5, 7.0, 7.5} {
    \fill[popline] (box1.north west) ++(0.3,-\y) rectangle ++(6.9,0.25);
}
% Formule politesse
\fill[popmuted!20] (box1.north west) ++(0.3,-8.2) rectangle ++(6.9,1.0);
\node[font=\tiny, align=left, text=popdark] at ([shift={(3.5,-8.7)}]box1.north west) {Formule de politesse \\ Signature};
% Flèches d'air
\draw[popgreen, <->, thick] (box1.north west) ++(0.3,-1.5) -- ++(0, -0.3) node[midway, right, font=\tiny, xshift=0.3cm] {Espace};
\draw[popgreen, <->, thick] (box1.north west) ++(0.3,-3.6) -- ++(0, -0.4) node[midway, right, font=\tiny, xshift=0.3cm] {Espace};
\end{scope}
% Vignette 2 : Surcharge visuelle (DROITE)
\begin{scope}[xshift=9.5cm]
\node[draw=poporange, thick, rounded corners, fill=poporange!5, minimum width=7.5cm, minimum height=10cm] (box2) {};
\node[font=\bfseries\small, text=poporange] at ([yshift=-0.4cm]box2.north) {Mise en page SURCHARGÉE (À éviter)};
% En-tête trop gros, police fantaisie simulée
\fill[poppink!30] (box2.north west) ++(0.1,-0.1) rectangle ++(7.3,1.8);
\node[font=\tiny, align=center, text=popdark, rotate=-3] at ([shift={(3.7,-1.0)}]box2.north west) {NOM Prénom -- Adresse -- Tel -- Email -- LinkedIn -- Photo};
% Pas d'espace, tout collé
\fill[popline] (box2.north west) ++(0.1,-2.0) rectangle ++(7.3,0.3);
\fill[popline] (box2.north west) ++(0.1,-2.3) rectangle ++(7.3,0.3);
\fill[popline] (box2.north west) ++(0.1,-2.6) rectangle ++(7.3,0.3);
% Corps dense, interligne nul
\foreach \y in {3.0, 3.3, 3.6, 3.9, 4.2, 4.5, 4.8, 5.1, 5.4, 5.7, 6.0, 6.3, 6.6, 6.9, 7.2, 7.5, 7.8, 8.1, 8.4} {
    \fill[popline!80] (box2.north west) ++(0.1,-\y) rectangle ++(7.3,0.18);
}
% Marge droite inexistante, texte jusqu'au bord
% Formule collée
\fill[popline] (box2.north west) ++(0.1,-8.7) rectangle ++(7.3,0.8);
\node[font=\tiny, align=left, text=popdark] at ([shift={(3.7,-9.1)}]box2.north west) {CordialementSignature};
% Signes d'alerte
\node[font=\tiny, text=poporange, rotate=15] at ([shift={(-1.0, -1.5)}]box2.north east) {Police fantaisie!};
\node[font=\tiny, text=poporange] at ([shift={(-0.2, 0)}]box2.east) {Marges $\approx$ 0!};
\node[font=\tiny, text=poporange, rotate=-15] at ([shift={(-1.0, 1.5)}]box2.south east) {Interligne 1.0!};
\end{scope}
\end{tikzpicture}
\legende{Comparaison visuelle : à gauche, une mise en page aérée, structurée, respirante ; à droite, une surcharge visuelle (marges nulles, interligne minimal, police décorative, blocs collés) qui fatigue l'œil et signale l'amateurisme.}
\end{popfigure}

\courssub{Lire une image : dénotation et connotation sur les offres d'emploi et logos}

Appliquons la méthode aux documents que vous \emph{recevez} (offres d'emploi) et aux signes que vous \emph{utilisez} (logos).

\begin{exoresolu}[Analyse sémiotique d'une affiche d'offre d'emploi (ex. : MTN Cameroun)]
\textbf{Document :} Affiche d'un panneau publicitaire ou publication LinkedIn : fond jaune MTN, logo MTN en haut à gauche, titre « Rejoignez l'aventure ! », photo d'une équipe jeune, mixte, souriante, casque sur la tête, en réunion. Bas de page : QR code + « mt.cm/carrieres ».

\textbf{1. Dénotation (Description objective) :}
\begin{itemize}
    \item Couleurs dominantes : Jaune (Pantone MTN) et noir/blanc.
    \item Éléments graphiques : Logo MTN (lettres MTN stylisées), code QR, URL.
    \item Personnages : 4 personnes (2 femmes, 2 hommes), 20--30 ans, tenues décontractées-chic (polos, chemises), casques audio, tablettes, sourires, regard vers l'interlocuteur ou caméra.
    \item Texte : « Rejoignez l'aventure ! », « MTN Cameroon », « Carrières ».
\end{itemize}

\textbf{2. Connotation (Interprétation culturelle / effets recherchés) :}
\begin{itemize}
    \item \textbf{Jaune + Logo :} Identification immédiate de la marque $\rightarrow$ \textbf{Fiabilité, notoriété, leader télécom}. Rassure le candidat : « C'est une vraie offre, pas une arnaque. »
    \item \textbf{Équipe jeune, mixte, souriante :} $\rightarrow$ \textbf{Diversité, inclusion, ambiance de travail agréable, modernité}. Cible les jeunes diplômés.
    \item \textbf{Casques / tablettes :} $\rightarrow$ \textbf{Technologie, innovation, outils modernes}. « Ici on travaille avec le numérique. »
    \item \textbf{« Rejoignez l'aventure » (impératif + nom fort) :} $\rightarrow$ \textbf{Challenge, carrière passionnante, pas juste un job}. Appel à l'engagement.
    \item \textbf{QR Code / URL courte :} $\rightarrow$ \textbf{Simplicité, rapidité, digital-first}. « Postulez en 1 clic depuis votre téléphone. »
\end{itemize}

\textbf{3. Message global :} « MTN est une entreprise moderne, inclusive, technologique, où il fait bon travailler. Toi, jeune talent, tu as ta place ici. Postule maintenant. »
\tcblower
\textbf{Solution détaillée :} L'analyse montre que l'image n'est pas décorative : elle \emph{argumente} visuellement. Le logo authentifie (éthos), l'équipe projette le candidat dans un avenir désirable (pathos), le QR code facilite l'action (logos pragmatique). En tant que candidat, vous devez \emph{décoder} cela pour adapter votre lettre : utiliser le vocabulaire de l'innovation, de l'équipe, du défi.
\end{exoresolu}

\begin{exemplebox}[Le logo d'entreprise sur votre lettre : un gage de sérieux]
Si vous répondez à une offre de \textbf{Brasseries du Cameroun} (logo étoile rouge / BC), placer discrètement ce logo (ou la mention « Brasseries du Cameroun ») dans l'en-tête de votre lettre (côté destinataire) ou dans l'objet (\texttt{Objet : Candidature -- Brasseries du Cameroun -- Poste de...}) opère une \textbf{connotation d'appartenance} : « Je sais à qui j'écris, j'ai ciblé ma candidature. » À l'inverse, une lettre générique « À l'attention du DRH » sans nom d'entreprise connote « Candidature en masse, peu motivée. »
\end{exemplebox}

\courssub{Cohérence entre l'image (photo, mise en page) et le discours de la lettre}

C'est le point nodal. Votre dossier est un \textbf{système sémiotique unifié}. Toute dissonance crée une \emph{rupture de confiance}.

\begin{center}
\begin{tabularx}{\linewidth}{|l|X|X|}
\hline
\rowcolor{popblueL} \textbf{Poste visé / Discours de la lettre} & \textbf{Image attendue (Photo + Mise en page)} & \textbf{Dissonance fatale} \\ \hline
Technicien maintenance / « Rigueur, sécurité, normes » & Photo : tenue propre, casque EPI en main ou fond atelier propre. Mise en page : technique, structurée, police standard, tableaux nets. & Photo en tenue de soirée, mise en page fleurie, police fantaisie. \\ \hline
Commercial / « Dynamisme, relationnel, présentation » & Photo : sourire franc, costume/cravate ou veste soignée, regard engageant. Mise en page : aérée, moderne, touches de couleur corporate (bleu/orange). & Photo sérieuse/fermée, mise en page dense, austère, papier bas de gamme. \\ \hline
Secrétariat / « Organisation, discrétion, maîtrise bureautique » & Photo : neutre professionnel, coiffure soignée. Mise en page : \emph{parfaite} (alignements au mm, tabulations justes, pas de faute, police standard). & Mise en page avec fautes d'alignement, onglets mal gérés, police multiple $\rightarrow$ « Il/elle ne maîtrise pas Word. » \\ \hline
Design / Communication / « Créativité, sens esthétique » & Photo : soignée mais personnalité visible (accessoire discret, cadrage original). Mise en page : \emph{signature visuelle} (en-tête design, palette cohérente, hiérarchie visuelle forte). & Mise en page standard « modèle Word 2003 », photo banale $\rightarrow$ « Pas de sens graphique. » \\ \hline
\end{tabularx}
\end{center}

\begin{methode}[Vérifier la cohérence image/texte (Check-list finale)]
Avant d'envoyer / déposer :
\begin{enumerate}
    \item Imprimez (ou affichez en plein écran) la lettre + CV + photo côte à côte.
    \item Demandez à un tiers : « Si vous ne lisiez que l'image, quel adjectif me donneriez-vous ? » (Rigoureux ? Dynamique ? Brouillon ? Soigné ?)
    \item Comparez avec les adjectifs-clés de votre lettre (ex. : « rigoureux », « autonome », « créatif »).
    \item Si écart $\rightarrow$ \textbf{corrigez l'image} (refaire la photo, refaire la mise en page) plutôt que d'adoucir le texte.
\end{enumerate}
\end{methode}

\courssub{Analyse guidée : deux dossiers de candidature photographiés}

Imaginons deux candidats au même poste : \textbf{Technicien de maintenance industrielle} à la \textbf{SONARA} (Société Nationale de Raffinage). La lettre est identique (compétences techniques, stage à la SONARA, motivation sécurité). Seule l'image diffère.

\begin{exoresolu}[Dossier A (Soigné) vs Dossier B (Négligé) -- Lecture par le DRH]
\textbf{Dossier A :}
\begin{itemize}
    \item \textbf{Photo :} Fond blanc, chemise bleue marine repassée, badge « Stagiaire SONARA 2024 » visible au col, casque de chantier (blanc) posé sur l'épaule, regard droit, sourire léger.
    \item \textbf{Mise en page :} Papier 100g blanc, police Calibri 11, justifié, marges 2,5 cm. En-tête : Logo SONARA (petit, gris) à gauche, coordonnées à droite. Objet en gras. Paragraphes séparés par 6 pt. Signature scannée en bleu.
    \item \textbf{Ressenti DRH (30 sec) :} « Pro, connaît les codes sécurité (casque), a déjà mis les pieds chez nous (badge), soigne les détails (marge, signature). Je lis. »
\end{itemize}

\textbf{Dossier B :}
\begin{itemize}
    \item \textbf{Photo :} Selfie dans une voiture (volant visible), t-shirt rouge « I love Yaoundé », lunettes de soleil sur la tête, fond flou mais on devine un marché. Cadrée de travers.
    \item \textbf{Mise en page :} Papier recyclé grisâtre (80g), police Comic Sans MS 12, aligné à gauche mais marges irrégulières (1,5 cm à gauche, 4 cm à droite). Taches d'encre (imprimante à jet d'encre bouchée). Pas d'en-tête entreprise. Objet : « Demande d'emploi ». Signature au stylo bille bleu sur la copie (pas sur l'original).
    \item \textbf{Ressenti DRH (5 sec) :} « Pas de casque, pas de tenue, selfie voiture = pas sérieux. Comic Sans = amateur. Taches = négligence. Pas de logo = envoi en masse. \textbf{Corbeille.} »
\end{itemize}

\textbf{Conclusion :} La lettre n'a même pas été lue pour le Dossier B. L'image a \emph{décidé} à sa place. Le Dossier A a acheté le droit à la lecture.
\tcblower
\textbf{Analyse :} Le Dossier A utilise la \textbf{connotation} des codes industriels (casque, badge, propreté) pour valider le \textbf{discours} (« Je maîtrise la sécurité »). Le Dossier B active des connotations contraires (loisir, désinvolture, amateurisme) qui \emph{infirment} le discours. \textbf{L'image est un argument visuel.}
\end{exoresolu}

\courssub{Complément utile au Probatoire : l'argumentation visuelle dans les publicités d'entreprises camerounaises}

\begin{savaistu}[L'argumentation visuelle chez MTN, Orange, Brasseries du Cameroun]
Au Probatoire / Baccalauréat technique, on peut vous demander d'analyser un corpus d'images publicitaires (affiches, spots vidéo figés). Les grands groupes camerounais utilisent des \textbf{stratégies visuelles codifiées} qu'il faut savoir repérer :

\begin{itemize}
    \item \textbf{MTN (Jaune/Noir) :} \textbf{Code couleur exclusif.} Le jaune saturé \textbf{est} la marque. Visuels : \emph{Connexion humaine} (mains qui se touchent, regards, familles, jeunes entrepreneurs). \textbf{Argument visuel :} « La technologie rapproche. » $\rightarrow$ \cle{Pathos} (émotion) + \cle{Ethos} (leader fiable).
    \item \textbf{Orange (Orange/Blanc) :} \textbf{Code couleur + forme carrée (logo).} Visuels : \emph{Quotidien simplifié} (paiement mobile, école, agriculture, e-santé). Mise en page souvent \emph{épurée, géométrique}. \textbf{Argument visuel :} « L'innovation au service du concret. » $\rightarrow$ \cle{Logos} (utilité rationnelle) + \cle{Ethos} (modernité).
    \item \textbf{Brasseries du Cameroun (Rouge/Blanc/Or -- étoiles) :} \textbf{Héritage + Fête + Partage.} Visuels : \emph{Moments de convivialité} (ngondo, matchs, après-travail, famille), bouteille/verre central, lumière dorée. \textbf{Argument visuel :} « La bière qui unit les Camerounais depuis 1948. » $\rightarrow$ \cle{Pathos} (identité collective, nostalgie joyeuse) + \cle{Ethos} (ancrage national).
\end{itemize}

\textbf{Transfert à votre candidature :} Comme ces marques, vous avez une \textbf{identité visuelle} (photo + mise en page + police). Elle doit \emph{argumenter} votre « marque personnelle » en 3 secondes. Choisissez vos codes (couleur discrète sur en-tête ? photo casque ? police technique ?) comme MTN choisit son jaune.
\end{savaistu}

\begin{exercice}[Entraînement : Audit visuel de votre dossier]
\begin{enumerate}
    \item Prenez votre lettre de motivation actuelle (ou un brouillon) et votre photo d'identité numérique.
    \item Appliquez la \textbf{Méthode d'analyse} (Dénotation $\rightarrow$ Connotation $\rightarrow$ Cohérence).
    \item Remplissez le tableau ci-dessous pour votre propre cas :
\end{enumerate}
\begin{center}
\begin{tabularx}{\linewidth}{|l|X|X|X|}
\hline
\rowcolor{popblueL} \textbf{Élément visuel} & \textbf{Dénotation (Ce que je vois)} & \textbf{Connotation (Ce que ça dit de moi)} & \textbf{Cohérence avec ma lettre (Oui/Non + Action)} \\ \hline
Photo (tenue, fond, cadrage) & & & \\ \hline
Police / Taille / Interligne & & & \\ \hline
Marges / Espaces / Alignements & & & \\ \hline
En-tête (Logo entreprise ?) & & & \\ \hline
Qualité papier / PDF (nom fichier) & & & \\ \hline
\end{tabularx}
\end{center}
\begin{enumerate}
    \setcounter{enumi}{3}
    \item Identifiez \textbf{3 corrections concrètes} à apporter avant le dépôt final.
\end{enumerate}
\end{exercice}

\begin{corrige}[Exercice : Audit visuel (exemple de réponse type)]
\begin{enumerate}
    \item \textbf{Photo :} Dénotation : « Selfie, t-shirt gris, fond chambre, ombre forte. » Connotation : « Étudiant, pas pro, désorganisé. » Cohérence : \textbf{Non}. Action : \textbf{Refaire photo pro} (photographe ou ami avec fond blanc, chemise, lumière fenêtre).
    \item \textbf{Mise en page :} Dénotation : « Police Calibri 11, justifié, marges 2 cm, interligne 1,15. » Connotation : « Propre, standard, lisible. » Cohérence : \textbf{Oui}. Action : \textbf{Rien} (vérifier juste les tabulations).
    \item \textbf{En-tête :} Dénotation : « Pas de logo entreprise, juste mes coordonnées. » Connotation : « Candidature générique. » Cohérence : \textbf{Non} pour poste ciblé. Action : \textbf{Ajouter logo entreprise} (téléchargé sur site officiel, 1,5 cm hauteur) dans bloc destinataire.
    \item \textbf{PDF :} Dénotation : « Fichier nommé \texttt{Lettre\_finale\_v3.pdf}. » Connotation : « Brouillons multiples, pas pro. » Cohérence : \textbf{Non}. Action : \textbf{Renommer} \texttt{NOM\_Prenom\_LettreMotivation\_Entreprise.pdf}.
    \item \textbf{3 corrections :} 1) Refaire photo pro. 2) Ajouter logo entreprise dans en-tête. 3) Renommer fichier PDF proprement.
\end{enumerate}
\end{corrige}

\begin{aretenir}[Synthèse : L'image, argument silencieux mais décisif]
\begin{itemize}
    \item La communication par l'image \textbf{informe, séduit, convainc} en parallèle du texte.
    \item \textbf{Photo :} Code strict (tenue, cadrage, fond, lumière) = 1ère preuve de professionnalisme.
    \item \textbf{Mise en page :} Police unique standard, marges généreuses, alignements rigoureux, papier/PDF soigné = preuve de maîtrise bureautique et respect du destinataire.
    \item \textbf{Lecture d'image :} Dénotation (objectif) $\rightarrow$ Connotation (subjectif/culturel) $\rightarrow$ Message global.
    \item \textbf{Cohérence :} L'image doit \emph{valider} le discours de la lettre (ex. : discours sécurité $\leftrightarrow$ photo casque). Toute dissonance = rejet rapide.
    \item \textbf{Recruteur :} < 1 minute par dossier. L'image guide l'œil et décide de la lecture.
\end{itemize}
\end{aretenir}

\coursec{Production guidée, compte rendu et remédiation}

Après avoir exploré la communication par l'image et son rôle dans la valorisation d'une candidature, nous entrons dans la phase ultime : la mise en œuvre intégrée de tous les acquis. Cette section transforme les connaissances théoriques (structure externe et interne, langage, argumentation) en compétence opérationnelle. L'objectif est de vous rendre autonome face à une offre réelle, dans les conditions de l'examen (Probatoire / Baccalauréat technique) et de la vie professionnelle.

\courssub{Les six étapes de la production : une démarche d'ingénierie textuelle}

Rédiger une lettre de motivation n'est pas un acte d'inspiration soudaine, mais un processus itératif et rigoureux. Tout comme un technicien suit un protocole pour réaliser une pièce, le rédacteur suit un flux de travail structuré : analyser, brainstormer, planifier, rédiger, réviser, réécrire.

\begin{popfigure}
\begin{tikzpicture}[
    node distance=1.2cm and 1.6cm,
    every node/.style={font=\small},
    stepbox/.style={rectangle, rounded corners, draw=popblue, thick, fill=popblueL, minimum width=2.6cm, minimum height=1.1cm, align=center},
    arrow/.style={-Stealth, thick, popdark},
    decision/.style={diamond, draw=poporange, thick, fill=poporangeL, minimum width=2.4cm, minimum height=1.2cm, align=center, aspect=2}
]
\node[stepbox, align=center] (analyser){1. Analyser\\l'offre};
\node[stepbox, right=of analyser, align=center] (brainstorm){2. Brainstormer\\ses atouts (Moi)};
\node[stepbox, right=of brainstorm, align=center] (planifier){3. Planifier\\(Vous-Moi-Nous)};
\node[stepbox, right=of planifier, align=center] (rediger){4. Rédiger\\le brouillon};
\node[decision, below=1.6cm of rediger, align=center] (reviser){5. Réviser\\(Grille OFOPS)};
\node[stepbox, left=of reviser, align=center] (reecriture){6. Réécrire\\Version finale};
\draw[arrow] (analyser) -- (brainstorm);
\draw[arrow] (brainstorm) -- (planifier);
\draw[arrow] (planifier) -- (rediger);
\draw[arrow] (rediger) -- (reviser) node[midway, right] {Auto-évaluation};
\draw[arrow] (reviser.west) -- node[midway, above] {Si conforme} (reecriture.east);
\draw[arrow] (reviser.north) -- node[midway, right] {Si non conforme} (rediger.south);
\end{tikzpicture}
\legende{Diagramme de flux : le processus itératif de production d'une lettre de motivation. La révision (étape 5) peut renvoyer à la rédaction (étape 4) jusqu'à satisfaction, avant la réécriture finale (étape 6).}
\end{popfigure}

\courssub{Étape 1 -- Analyser l'offre (détecter les attentes implicites)}

Avant d'écrire un seul mot, vous devez \emph{décortiquer} l'annonce. Ne lisez pas en diagonale.

\begin{methode}[Analyse d'une offre d'emploi ou de stage]
\begin{enumerate}
\item \textbf{Surligner les mots-clés :} compétences techniques (ex. : \og maintenance hydraulique\fg, \og lecture de plans\fg), savoir-être (ex. : \og rigueur\fg, \og esprit d'équipe\fg), diplômes requis.
\item \textbf{Identifier le destinataire :} nom de l'entreprise, service (DRH, chef d'atelier), nom du responsable si possible. \textbf{Personnaliser l'en-tête change tout.}
\item \textbf{Repérer les valeurs de l'entreprise :} site web, devise, actualité. Utilisez leur vocabulaire (ex. : \og innovation\fg, \og service client\fg, \og continuité de service\fg).
\item \textbf{Noter les contraintes :} date limite, mode d'envoi (courriel, postal, dépôt physique), pièces jointes demandées (CV, copies de diplômes, attestations).
\end{enumerate}
\end{methode}

\begin{exemplebox}[Exemple d'analyse rapide]
\textbf{Offre :} \og Technicien de maintenance industrielle (stage de fin d'études). Missions : maintenance préventive et curative de lignes automatisées. Profil : Bac technique / BTS Maintenance, lecture de plans électriques et hydrauliques, autonomie, réactivité. Entreprise : CAMWATER (Douala).\fg

\textbf{Analyse :}
\begin{itemize}
\item \textbf{Mots-clés techniques :} préventive, curative, automatismes, schémas électriques et hydrauliques.
\item \textbf{Savoir-être :} autonomie, réactivité, rigueur (sécurité).
\item \textbf{Contexte :} société de distribution d'eau potable $\rightarrow$ sens du service public, continuité de service.
\item \textbf{Plan Vous-Moi-Nous :} Vous (besoin de techniciens polyvalents pour la continuité du service) -- Moi (formation maintenance, stage antérieur sur pompes, lecture de plans) -- Nous (stage pour prouver mon efficacité et me former à vos procédures spécifiques).
\end{itemize}
\end{exemplebox}

\courssub{Étape 2 -- Brainstormer ses atouts (la colonne \og Moi\fg)}

C'est l'inventaire honnête de vos ressources. Utilisez un tableau à deux colonnes pour faire correspondre \textbf{l'offre} et \textbf{votre profil}. Chaque exigence de l'annonce doit recevoir une \emph{preuve} concrète : stage, projet, matière, certification.

\begin{center}
\begin{tabularx}{\linewidth}{|X|X|}\hline
\rowcolor{popblueL}
\textbf{Exigence de l'offre (Vous)} & \textbf{Ma preuve (Moi)} \\ \hline
Maintenance préventive & Stage de 6 semaines chez ENEO : planification hebdomadaire de la maintenance de 10 postes HTA/BT \\ \hline
Lecture de plans électriques et hydrauliques & Projet de fin d'études : schéma complet d'une station de pompage \\ \hline
Autonomie / Réactivité & Intervention seul sur une panne de groupe électrogène (résolue en 2 heures) \\ \hline
Sens du service public / Sécurité & Formation aux gestes de premiers secours validée, respect strict des consignes \\ \hline
\end{tabularx}
\end{center}

\courssub{Étape 3 -- Planifier selon la structure \og Vous -- Moi -- Nous\fg}

C'est le squelette de l'argumentation. Ne commencez jamais par \og Moi, je...\fg : la lettre de motivation s'ouvre sur l'entreprise, pas sur soi.

\begin{definition}[Structure interne type (plan Vous-Moi-Nous)]
\begin{itemize}
\item \textbf{Paragraphe 1 (VOUS) :} l'accroche. Montrez que vous comprenez l'entreprise, son contexte, son besoin précis. \og Votre recherche d'un technicien de maintenance pour assurer la continuité du service de distribution d'eau à Douala retient toute mon attention.\fg
\item \textbf{Paragraphe 2 (MOI) :} la preuve. Vos compétences techniques et humaines \textbf{liées aux mots-clés} de l'offre. \og Fort de ma formation en maintenance et de mon stage chez ENEO, où j'ai participé à la maintenance préventive de 10 postes, je maîtrise la lecture de plans hydrauliques et l'intervention en autonomie sur automatismes.\fg
\item \textbf{Paragraphe 3 (NOUS) :} la projection. Ce que vous apportez à l'entreprise + demande d'entretien. \og Intégrer vos équipes serait l'opportunité de mettre ma rigueur et ma réactivité au service de la performance de votre société. Je serais honoré de vous exposer ma motivation lors d'un entretien.\fg
\end{itemize}
\end{definition}

\courssub{Étape 4 -- Rédiger le brouillon (45 minutes chrono)}

C'est la simulation d'examen. Respectez scrupuleusement la \textbf{structure externe} (en-tête, objet, formules) et la \textbf{structure interne} (trois paragraphes).

\begin{attention}[Piège du temps]
En examen, 45 minutes passent très vite.
\begin{itemize}
\item 5 min : analyse de l'offre + brainstorming + plan détaillé (mots-clés par paragraphe).
\item 30 min : rédaction soignée au brouillon (phrase par phrase, en appliquant mentalement la grille OFOPS).
\item 10 min : relecture OFOPS + copie propre au propre.
\end{itemize}
N'écrivez jamais la version finale directement. Le brouillon est votre espace de sécurité.
\end{attention}

\courssub{Étape 5 -- Réviser avec la méthode OFOPS (auto-évaluation)}

C'est l'outil qualité. Appliquez-le \textbf{systématiquement} avant de rendre la copie.

\begin{methode}[La relecture en cinq points -- acronyme OFOPS]
\begin{enumerate}
\item \textbf{O -- Orthographe / Grammaire :} accords (sujet-verbe, groupe nominal), homophones (a/à, et/est, ou/où), majuscules (début de phrase, noms propres, \og Madame, Monsieur\fg). \emph{Astuce : relisez à l'envers (de la dernière phrase vers la première) pour casser le sens et repérer les fautes.}
\item \textbf{F -- Formules de politesse :} en-tête (Monsieur le Directeur Général / Madame la Responsable des Ressources Humaines), formule d'appel (\og Madame, Monsieur\fg), formule de fin (\og je vous prie d'agréer...\fg), signature (prénom NOM).
\item \textbf{O -- Objet :} clair, précis, avec la référence de l'offre si elle existe. \og Objet : Candidature au poste de technicien de maintenance (Réf. CAM/2026/045)\fg.
\item \textbf{P -- Plan (structure interne) :} y a-t-il bien trois paragraphes distincts (Vous / Moi / Nous) ? Les connecteurs logiques sont-ils présents (\og Fort de...\fg, \og Par ailleurs...\fg, \og C'est pourquoi...\fg) ? L'argumentation suit-elle l'offre ?
\item \textbf{S -- Signature / Présentation :} espace pour la signature manuscrite, pièces jointes mentionnées (\og Pièces jointes : CV, copie des diplômes, attestations de stage\fg), aération (sauts de ligne entre paragraphes), marges régulières, écriture lisible.
\end{enumerate}
\end{methode}

\courssub{Grille d'auto-évaluation et d'évaluation par les pairs (critères pondérés)}

Cette grille est votre \og tableau de bord\fg. Elle est calquée sur les attentes des épreuves de production d'écrits du Probatoire et du Baccalauréat technique.

\begin{popfigure}
\small
\begin{center}
\begin{tabularx}{\linewidth}{|l|X|c|c|}\hline
\rowcolor{popblueL}
\textbf{Critère} & \textbf{Indicateurs de réussite (niveau maximal)} & \textbf{Points} & \textbf{Note} \\ \hline
\textbf{1. Structure externe} & En-tête complet (expéditeur, destinataire, lieu, date). Objet précis + référence. Formules d'appel et de fin conformes aux usages. Signature + pièces jointes mentionnées. & /4 & \\ \hline
\textbf{2. Structure interne} & Plan Vous-Moi-Nous respecté (3 paragraphes distincts). Introduction accrocheuse (contexte de l'entreprise). Corps : arguments techniques et humains \textbf{liés à l'offre}. Conclusion : projection + demande d'entretien. Connecteurs fluides. & /6 & \\ \hline
\textbf{3. Correction langagière} & Aucune faute d'orthographe ou de grammaire majeure. Syntaxe correcte (phrases complètes). Vocabulaire technique précis. Registre formel soutenu (pas de familier). Impératif utilisé à bon escient (formules). & /4 & \\ \hline
\textbf{4. Argumentation / Pertinence} & Preuves concrètes (stages, projets, matières) et non généralités. Adéquation profil/poste évidente. Ton confiant mais pas arrogant. Personnalisation (valeurs, actualité de l'entreprise). & /4 & \\ \hline
\textbf{5. Présentation / Lisibilité} & Mise en page aérée (marges, interlignes). Propreté (pas de ratures sur la copie finale). Respect du format lettre (expéditeur en haut à gauche, destinataire à droite). & /2 & \\ \hline
\rowcolor{popblueL}
\textbf{TOTAL} & & \textbf{/20} & \\ \hline
\end{tabularx}
\end{center}
\legende{Grille d'évaluation pondérée (barème type examen). Utilisez-la pour l'auto-évaluation, la correction croisée et la remédiation. Visez 16/20 et plus pour l'excellence.}
\end{popfigure}

\courssub{Exercice de production intégrée : simulation d'examen (45 min)}

\begin{exoresolu}[Sujet d'examen type -- Probatoire technique]
\textbf{Consigne :} Vous êtes \textbf{TCHINDA Kevin}, élève en classe de Première F4 (Électrotechnique) au Lycée Technique de Bafoussam. Vous répondez à l'offre ci-dessous pour un stage de vacances (8 semaines). Rédigez la lettre de motivation. (Durée : 45 min -- Notation sur 20 selon la grille ci-dessus).

\textbf{Offre (extrait) :}
\og \textbf{ENTREPRISE : ENEO (Direction Régionale de l'Ouest).}
\textbf{Poste :} stagiaire technicien réseaux BT/HTA (H/F).
\textbf{Missions :} participation aux travaux de renforcement et d'extension du réseau BT/HTA ; levés, implantation de poteaux, câblage, raccordement des clients ; respect strict des consignes de sécurité (port des EPI, consignation).
\textbf{Profil :} élève de Première ou Terminale F3/F4 Électrotechnique ; lecture de schémas électriques ; sens de la sécurité, travail en équipe.
\textbf{Adresse :} BP 1234 Bafoussam. Réf. : STAGE/2026/001.\fg

\tcblower
\textbf{Solution détaillée (version finale attendue)}

\begin{center}
\begin{tabularx}{\linewidth}{|X|}\hline
\rowcolor{popcream}
\textbf{Kevin TCHINDA} \\
Quartier Banengo, BP 456 Bafoussam \\
Tél. : 6XX XX XX XX -- kevin.tchinda@email.com \\ \hline
\end{tabularx}
\end{center}
\vspace{0.3cm}
\hfill Bafoussam, le 15 mars 2026

\begin{flushright}
\textbf{À l'attention de Monsieur le Directeur Régional de l'Ouest} \\
\textbf{ENEO} \\
\textbf{BP 1234 Bafoussam}
\end{flushright}
\vspace{0.5cm}

\textbf{Objet : Candidature au stage de technicien réseaux BT/HTA (Réf. STAGE/2026/001)}

Madame, Monsieur,

\vspace{0.2cm}
\noindent \textbf{(VOUS)} Votre programme de renforcement du réseau de distribution dans la région de l'Ouest, notamment l'extension des lignes pour améliorer la desserte des zones périurbaines, mobilise des compétences techniques rigoureuses que je souhaite mettre à profit au sein de vos équipes.

\vspace{0.2cm}
\noindent \textbf{(MOI)} Élève en classe de Première F4 Électrotechnique au Lycée Technique de Bafoussam, j'ai acquis de solides bases en installations électriques basse tension à travers mes travaux pratiques de câblage et mes leçons de technologie. Mon stage d'observation de l'an dernier dans une entreprise d'installation électrique de Bafoussam m'a permis de participer au raccordement de clients basse tension, dans le strict respect des procédures de sécurité et du port des équipements de protection individuelle. Cette expérience a développé mon sens de la sécurité, ma capacité à lire des schémas électriques et mon aptitude au travail en équipe sur chantier.

\vspace{0.2cm}
\noindent \textbf{(NOUS)} Intégrer la Direction Régionale de l'Ouest d'ENEO représenterait pour moi l'aboutissement logique de ma formation et une opportunité unique de confronter mes acquis aux réalités du terrain sous l'encadrement de vos techniciens confirmés. Ma motivation, ma rigueur et ma disponibilité immédiate sont des atouts que je serais ravi de vous présenter lors d'un entretien à votre convenance.

\vspace{0.2cm}
\noindent En vous remerciant de l'attention portée à ma candidature, je vous prie d'agréer, Madame, Monsieur, l'expression de mes salutations distinguées.

\vspace{1cm}
\hfill \textbf{Kevin TCHINDA} \\
\hfill \textit{(Signature manuscrite)}

\vspace{0.2cm}
\noindent \textbf{Pièces jointes :} curriculum vitae, copie du bulletin de Première, attestation de stage, photocopie de la CNI.
\end{exoresolu}

\begin{aretenir}[Analyse de la copie type (pourquoi 18-19/20 ?)]
\begin{itemize}
\item \textbf{Structure externe (4/4) :} en-tête complet, objet avec référence, formules parfaites, signature + pièces jointes.
\item \textbf{Structure interne (6/6) :} trois paragraphes nets (Vous/Moi/Nous). Le paragraphe 1 cite l'actualité de l'entreprise (renforcement du réseau de l'Ouest). Le paragraphe 2 donne des preuves précises (travaux pratiques, stage, EPI, schémas). Le paragraphe 3 projette (aboutissement, encadrement, disponibilité).
\item \textbf{Langue (4/4) :} registre soutenu, vocabulaire technique précis (consignation, EPI, basse tension), aucune faute, impératif réservé aux formules.
\item \textbf{Argumentation (4/4) :} adéquation parfaite profil/poste ; preuves concrètes, pas de généralités.
\item \textbf{Présentation (2/2) :} aérée, lisible, professionnelle.
\end{itemize}
\end{aretenir}

\courssub{Correction croisée entre pairs (co-évaluation)}

L'évaluation par les pairs est un levier puissant : en corrigeant l'autre, on affine son propre regard critique.

\begin{methode}[Protocole de correction croisée (15 min)]
\begin{enumerate}
\item \textbf{Binôme :} échangez vos copies (brouillons ou versions finales).
\item \textbf{Lecture silencieuse (3 min) :} lisez une première fois pour le sens global, une seconde fois avec la grille OFOPS.
\item \textbf{Annotation codée (7 min) :} dans la marge, codez :
\begin{itemize}
\item \textcolor{popgreen}{$\checkmark$} = bon / conforme ;
\item \textcolor{poporange}{$\sim$} = à améliorer (précisez : \og formule de fin\fg, \og accord\fg, \og paragraphe Moi vague\fg) ;
\item \textcolor{poppink}{$\times$} = erreur / manquant (ex. : pas d'objet, pas de signature).
\end{itemize}
\item \textbf{Remplissage de la grille (3 min) :} notez sur /20 avec une justification brève par critère.
\item \textbf{Retour oral (2 min) :} \og 3 points forts / 2 points à travailler\fg. Bienveillance et exigence.
\end{enumerate}
\end{methode}

\courssub{Compte rendu collectif : bilan des acquis et difficultés fréquentes}

Après la correction croisée, le professeur anime une synthèse au tableau. C'est le moment de \textbf{nommer} les erreurs récurrentes pour les désamorcer collectivement.

\begin{experience}[Bilan type -- tableau de synthèse]
\begin{center}
\begin{tabularx}{\linewidth}{|l|X|X|}\hline
\rowcolor{popblueL}
\textbf{Domaine} & \textbf{Acquis validés (majorité)} & \textbf{Difficultés récurrentes (remédiation)} \\ \hline
\textbf{Structure externe} & En-tête, objet, formules d'appel et de fin & Oubli des \og Pièces jointes\fg ; confusion dans les titres (\og Monsieur le Directeur\fg{} / \og Madame la Responsable\fg) ; date mal placée. \\ \hline
\textbf{Structure interne} & Paragraphe \og Vous\fg{} souvent réussi (reprise de l'offre) & \og Moi\fg{} : généralités (\og je suis dynamique\fg) sans preuve ; \og Nous\fg{} : absent ou faible (\og j'espère être pris\fg). \\ \hline
\textbf{Langue / Style} & Registre formel globalement tenu & Accords (participes passés, groupe nominal) ; homophones (a/à, et/est) ; phrases trop longues ou mal ponctuées. \\ \hline
\textbf{Argumentation} & Citation des stages et des matières & Preuves non \textbf{liées} aux mots-clés de l'offre (décalage). \\ \hline
\textbf{Gestion du temps} & & Dépassement des 45 min (trop de temps sur le brouillon ou la recopie). \\ \hline
\end{tabularx}
\end{center}
\legende{Bilan type issu de la correction croisée. Ce tableau guide la remédiation ciblée.}
\end{experience}

\courssub{Remédiation ciblée : fiches d'exercices sur les points faibles}

En fonction du bilan, le professeur distribue des micro-fiches (10 à 15 min chacune) en ateliers tournants ou en devoirs.

\begin{exercice}[Fiche Remédiation 1 -- Formules de politesse (structure externe)]
\textbf{Consigne :} rédigez les formules d'appel et de fin adaptées à chaque destinataire.
\begin{enumerate}
\item Destinataire : \textbf{Monsieur le Directeur Général de la SONARA} (Limbé). Objet : demande de stage.
\item Destinataire : \textbf{Madame la Responsable des Ressources Humaines de CAMTEL} (Yaoundé). Objet : candidature à un emploi de technicien.
\item Destinataire : \textbf{Monsieur le Chef de Service Maintenance de la CDC} (Tiko). Objet : demande de stage de fin d'études en Électrotechnique.
\end{enumerate}
\textbf{Rappel :} formule d'appel : \og Madame, Monsieur,\fg{} ou le titre exact du destinataire ; formule de fin : \og Je vous prie d'agréer, + titre exact, l'expression de mes salutations distinguées.\fg{} \textbf{Accord :} le titre dans la formule de fin reprend le titre de l'en-tête.
\end{exercice}

\begin{corrige}[Exercice Fiche 1]
\begin{enumerate}
\item \textbf{Appel :} Monsieur le Directeur Général, \quad \textbf{Fin :} \og Je vous prie d'agréer, Monsieur le Directeur Général, l'expression de mes salutations distinguées.\fg
\item \textbf{Appel :} Madame la Responsable des Ressources Humaines, \quad \textbf{Fin :} \og Je vous prie d'agréer, Madame la Responsable des Ressources Humaines, l'expression de mes salutations distinguées.\fg
\item \textbf{Appel :} Monsieur le Chef de Service Maintenance, \quad \textbf{Fin :} \og Je vous prie d'agréer, Monsieur le Chef de Service Maintenance, l'expression de mes salutations distinguées.\fg
\end{enumerate}
\end{corrige}

\begin{exercice}[Fiche Remédiation 2 -- L'impératif : valeurs et emplois (grammaire)]
\textbf{Consigne :} identifiez la valeur de l'impératif (ordre / conseil / prière / souhait / interdiction) et transformez la phrase à l'infinitif (consigne de sécurité) ou en phrase déclarative.
\begin{enumerate}
\item \og \textbf{Veillez} à joindre votre CV.\fg{} (lettre)
\item \og \textbf{Ne touchez} pas aux câbles sous tension.\fg{} (consigne de sécurité)
\item \og \textbf{Permettez-moi} de vous présenter ma candidature.\fg{} (lettre)
\item \og \textbf{Soyez} vigilant lors du levage.\fg{} (conseil)
\end{enumerate}
\end{exercice}

\begin{corrige}[Exercice Fiche 2]
\begin{enumerate}
\item \textbf{Valeur :} prière / demande polie. \textbf{Infinitif :} veiller à joindre votre CV.
\item \textbf{Valeur :} interdiction / ordre de sécurité. \textbf{Infinitif :} ne pas toucher aux câbles sous tension.
\item \textbf{Valeur :} souhait / formule de politesse (figée). \textbf{Déclaratif :} je vous demande la permission de vous présenter ma candidature.
\item \textbf{Valeur :} conseil / recommandation. \textbf{Infinitif :} être vigilant lors du levage.
\end{enumerate}
\end{corrige}

\begin{exercice}[Fiche Remédiation 3 -- Plan Vous-Moi-Nous : du flou au précis]
\textbf{Consigne :} voici un paragraphe \og Moi\fg{} faible. Réécrivez-le en utilisant le tableau \og Offre / Preuve\fg{} ci-dessous.
\og \textit{Je suis très motivé et sérieux. J'aime bien l'électricité et je travaille bien en équipe. Je veux apprendre.}\fg

\textbf{Offre ciblée :} technicien de maintenance en automatismes (diagnostic de pannes, lecture de schémas, programmation d'automates).
\textbf{Profil de l'élève :} projet de fin d'études : automatisme d'une ligne de conditionnement ; stage : diagnostic de capteurs et d'actionneurs sur une presse hydraulique ; matière : automatismes industriels (16/20).
\end{exercice}

\begin{corrige}[Exercice Fiche 3]
\og \textbf{Fort de ma formation en automatismes industriels (16/20 en Première), j'ai conçu, dans le cadre de mon projet de fin d'études, une ligne de conditionnement pilotée par automate programmable, en réalisant le câblage, la programmation et le diagnostic de pannes. Mon stage en entreprise m'a confronté au diagnostic réel de capteurs et d'actionneurs sur une presse hydraulique, développant ma rigueur et mon autonomie face aux arrêts de machine.}\fg
\end{corrige}

\courssub{Réécriture individuelle et constitution du portfolio}

La boucle se boucle : après la remédiation, chaque élève produit sa \textbf{version définitive}, corrigée, augmentée, prête pour l'examen et le monde professionnel.

\begin{methode}[De la copie d'examen au portfolio professionnel]
\begin{enumerate}
\item \textbf{Réécriture propre :} intégrez \textbf{toutes} les corrections (pairs + professeur). Soignez la présentation (marges régulières, écriture lisible, interlignes aérées).
\item \textbf{Numérisation :} scannez la lettre signée (PDF) et le CV (PDF). Nommez les fichiers clairement : \og NOM\_Prenom\_LettreMotiv\_Stage2026.pdf\fg{} / \og NOM\_Prenom\_CV\_2026.pdf\fg.
\item \textbf{Portfolio physique :} chemise cartonnée contenant : lettre originale signée, CV, copies des diplômes, attestations de stage, lettres de recommandation. Ordre chronologique inverse (le plus récent en haut).
\item \textbf{Portfolio numérique :} dossier sur clé USB ou espace de stockage en ligne, partagé avec le professeur ou le tuteur de stage.
\item \textbf{Mise à jour continue :} adaptez la lettre à \textbf{chaque} offre (modifiez les paragraphes 1 et 2). Gardez une \og lettre maîtresse\fg{} (tous vos atouts) pour piocher rapidement.
\end{enumerate}
\end{methode}

\begin{savaistu}[Le savais-tu ? -- L'examen technique et la lettre]
Au \textbf{Probatoire} et au \textbf{Baccalauréat technique}, l'épreuve de production d'écrits comprend fréquemment un \textbf{texte fonctionnel} : lettre de motivation, demande de stage, lettre administrative, compte rendu.
\begin{itemize}
\item C'est l'exercice le plus \textbf{prévisible} et le plus \textbf{rémunérateur} : la structure est codifiée, les formules sont fixes, le plan est imposé (Vous-Moi-Nous).
\item Contrairement à la dissertation ou au commentaire où l'inspiration varie, ici \textbf{la méthode garantit la note}. Un élève qui maîtrise la grille OFOPS et le plan Vous-Moi-Nous assure facilement 14 à 16/20, souvent davantage.
\item Les correcteurs sanctionnent lourdement : l'absence d'objet, les formules bâclées, le \og Je\fg{} initial, le manque de preuves concrètes, les fautes d'accord basiques.
\item \textbf{Conseil stratégique :} traitez cet exercice \textbf{en premier à l'examen} (20 à 25 min). Cela vous met en confiance, vous assure des points rapides et libère l'esprit pour la dissertation ou le résumé.
\end{itemize}
\end{savaistu}

\begin{attention}[Dernier piège fatal : le \og copier-coller\fg{} Internet]
Recopier une lettre modèle trouvée sur Internet sans la personnaliser est \textbf{immédiatement détectable} par un jury ou un recruteur :
\begin{itemize}
\item vocabulaire générique (\og dynamique, rigoureux, motivé\fg) sans \textbf{aucune preuve technique} ;
\item inadéquation totale avec l'offre (aucun mot-clé de l'annonce) ;
\item formules toutes faites qui ne correspondent pas au destinataire réel ;
\item \textbf{résultat :} note très faible (souvent moins de 6/20) ou candidature écartée.
\end{itemize}
\textbf{La personnalisation (le \og Vous\fg{} et le \og Moi\fg{} prouvé) est la seule signature de votre authenticité.}
\end{attention}

\begin{aretenir}[Check-list finale avant dépôt ou envoi -- \og zéro défaut\fg]
\begin{itemize}
\item[$\square$] En-tête complet (mes coordonnées + coordonnées du destinataire + lieu + date).
\item[$\square$] Objet : précis + référence de l'offre.
\item[$\square$] Formule d'appel exacte (civilité + titre).
\item[$\square$] Paragraphe 1 (VOUS) : contexte de l'entreprise / besoin identifié.
\item[$\square$] Paragraphe 2 (MOI) : 2 à 3 preuves techniques et humaines \textbf{liées à l'offre} (stages, projets, notes, certifications).
\item[$\square$] Paragraphe 3 (NOUS) : projection + demande d'entretien explicite.
\item[$\square$] Connecteurs logiques (\og Par ailleurs\fg, \og En effet\fg, \og C'est pourquoi\fg...).
\item[$\square$] Formule de fin exacte (reprise du titre) + espace pour la signature manuscrite.
\item[$\square$] \og Pièces jointes :\fg{} listées.
\item[$\square$] Relecture OFOPS effectuée (Orthographe, Formules, Objet, Plan, Signature).
\item[$\square$] Présentation impeccable (marges, aération, propreté).
\end{itemize}
\end{aretenir}

Cette section clôture l'unité. Vous disposez désormais de la \textbf{méthode complète}, des \textbf{outils d'évaluation}, de la \textbf{remédiation ciblée} et du \textbf{portfolio} pour aborder sereinement l'épreuve du Probatoire et du Baccalauréat technique et, surtout, pour décrocher vos stages et vos premiers emplois. La rigueur technique que vous appliquez en atelier, appliquez-la à votre écriture professionnelle : c'est la même exigence de qualité.

\newpage

\coursec{Activité d'intégration}

\coursec{Leçon N21 : Produire une lettre officielle (lettre de motivation)}

\courssub{Objectifs de l'activité}
\par
Cette activité d'intégration vise à mobiliser l'ensemble des savoirs et savoir-faire travaillés au cours de la leçon N21 dans une situation authentique de communication professionnelle. À l'issue de cette séquence, l'élève doit être capable de :
\begin{itemize}
    \item Analyser des lettres de motivation réelles pour en dégager les critères de qualité (structure, langue, style, mise en page) ;
    \item Rédiger une lettre de motivation complète, conforme aux normes de la correspondance administrative et professionnelle camerounaise ;
    \item Appliquer la structure externe (en-tête, objet, références, formules de politesse, signature) et la structure interne (plan \emph{Vous -- Moi -- Nous}) ;
    \item Intégrer au moins une \cle{énumération} et une \cle{gradation} dans le corps de la lettre pour valoriser son parcours ;
    \item Employer correctement le \cle{mode impératif} dans les formules de politesse de conclusion (valeur injonctive/politesse) ;
    \item Soigner la \cle{communication par l'image} (mise en page aérée, alignements, police lisible, en-tête visuel) ;
    \item S'auto-évaluer et évaluer ses pairs à l'aide de la grille \cle{OFOPS} (Orthographe, Formulation, Organisation, Ponctuation, Syntaxe) pour réviser son texte.
\end{itemize}

\courssub{Situation d'apprentissage}
\par
\textbf{Contexte authentique :} Nous sommes le 15 mai 2026. Vous êtes élève en classe de Première Technique au Lycée Technique de Garoua. Votre établissement a signé une convention de partenariat avec la \textbf{SODECOTON} (Société de Développement du Coton), dont le siège social est à Garoua (BP 302). La direction des Ressources Humaines lance sa campagne annuelle de stages académiques d'été (juillet--août 2026) pour les élèves de Première et Terminale des filières \emph{Maintenance industrielle}, \emph{Électrotechnique} et \emph{Agriculture tropicale}.

\par
L'avis d'appel à candidatures, affiché au tableau d'affichage du lycée et publié sur le site web de la SODECOTON, mentionne :
\begin{quote}
\textbf{OFFRE DE STAGE ACADÉMIQUE -- SESSION 2026}\par
\emph{La SODECOTON offre 50 places de stage d'une durée de 8 semaines (1er juillet -- 31 août 2026) aux élèves des lycées techniques partenaires.}\par
\emph{Filières concernées : Maintenance des équipements industriels, Électrotechnique, Génie rural / Agriculture.}\par
\emph{Conditions : Être régulièrement inscrit en Première ou Terminale technique ; présenter un dossier complet (CV + Lettre de motivation adressée au Directeur des Ressources Humaines).}\par
\emph{Date limite de dépôt des dossiers : \textbf{30 mai 2026} (cachet de la poste ou dépôt physique au service RH faisant foi).}\par
\emph{Adresse : Direction des Ressources Humaines, SODECOTON, BP 302 Garoua.}
\end{quote}

\par
Votre professeur de Français, M. \textsc{Mbarga}, décide de transformer cette opportunité en activité d'intégration notée. Il vous remet deux lettres de motivation anonymisées, écrites par d'anciens élèves du lycée pour la session 2025 : la \textbf{Lettre A} (refusée : mal structurée, fautes, ton familier) et la \textbf{Lettre B} (acceptée : structure rigoureuse, style professionnel, arguments ciblés). Votre mission : en groupe, analyser ces deux modèles pour construire la \emph{grille de réussite} ; ensuite, individuellement, rédiger votre propre lettre pour la session 2026 ; enfin, évaluer la lettre d'un camarade avec la grille OFOPS et réécrire la vôtre avant l'affichage des meilleures productions au « Tableau d'honneur » de la classe.

\par
\textbf{Données personnelles à utiliser (fictives mais réalistes) :}
\begin{itemize}
    \item Votre identité : \textsc{Nom Prénom} (ex. : \textsc{Djibrilla Amadou}), né le 12/03/2008 à Garoua.
    \item Adresse : Quartier \emph{Lamordé}, Rue 12, Porte 45, Garoua. Tél : 6xx xx xx xx. Email : prenom.nom@email.cm.
    \item Filière : \textbf{Maintenance industrielle} (option Systèmes automatisés). Moyenne annuelle : 14,5/20. Rang : 3/42.
    \item Expériences : Stage d'initiation de 2 semaines (juillet 2025) à l'atelier de maintenance de la \emph{CAMRAIL} (Garoua) -- soudure, lecture de plans, maintenance préventive sur locomotives. Projet de fin d'année : « Conception d'un système de graissage automatique pour chaîne de convoyeur » (note : 18/20).
    \item Qualités : Rigueur, esprit d'équipe, autonomie, respect des consignes de sécurité (HSE).
    \item Motivation spécifique : Intégrer le service Maintenance de l'usine d'égrenage de Garoua pour découvrir la maintenance prédictive sur métiers à égrener (Continental Eagle).
\end{itemize}

\courssub{Consignes de l'activité d'intégration}
\par
\textbf{Phase 1 -- Analyse comparative en groupes (30 min) :}
\begin{enumerate}
    \item Lire attentivement la \textbf{Lettre A} et la \textbf{Lettre B} (distribuées par le professeur).
    \item Compléter le \emph{Tableau d'analyse} ci-dessous en relevant pour chaque lettre : \textbf{Structure externe} (présence/absence des 8 zones), \textbf{Structure interne} (plan Vous/Moi/Nous respecté ?), \textbf{Outils langagiers} (énumération, gradation, impératif de politesse, vocabulaire technique), \textbf{Mise en page / Image} (alignements, marges, police, lisibilité), \textbf{Points forts / Points faibles}.
    \item Dégager collectivement les \textbf{5 critères indispensables} d'une lettre de motivation réussie pour la SODECOTON. Les inscrire sur l'affiche de groupe.
\end{enumerate}

\par
\textbf{Phase 2 -- Rédaction individuelle (45 min) :}
\begin{enumerate}
    \setcounter{enumi}{3}
    \item Rédiger votre lettre de motivation sur papier blanc (format A4), en \textbf{respectant scrupuleusement} :
    \begin{itemize}
        \item La structure externe complète (vos coordonnées, destinataire, date, lieu, objet, référence à l'avis, corps, formule de politesse finale, signature).
        \item Le plan \textbf{Vous -- Moi -- Nous} (un paragraphe par mouvement).
        \item L'intégration d'\textbf{au moins une énumération} (compétences, stages, qualités) et d'\textbf{au moins une gradation} (ex. : \og{}initiée, consolidée, maîtrisée\fg{} ; ou \og{}observer, comprendre, agir\fg{}).
        \item L'utilisation de l'\textbf{impératif de politesse} dans la formule de conclusion (ex. : \og{}Je vous prie d'agréer...\fg{} / \og{}Veuillez agréer...\fg{}).
        \item Une \textbf{mise en page soignée} : marges 2,5 cm, interligne 1,5, police Times New Roman 12 ou Arial 11, alignement justifié, objet en gras, paragraphes espacés.
    \end{itemize}
    \item Relire votre production en traquant les erreurs (OFOPS).
\end{enumerate}

\par
\textbf{Phase 3 -- Évaluation par les pairs et réécriture (30 min) :}
\begin{enumerate}
    \setcounter{enumi}{5}
    \item Échanger votre lettre avec un binôme. Évaluer la lettre de votre camarade à l'aide de la \textbf{Grille OFOPS} (notée sur 20 : 4 items × 5 points). Annoter les marges (code couleur : \textcolor{popgreen}{vert} = bon, \textcolor{poporange}{orange} = à améliorer, \textcolor{poppink}{rose} = erreur).
    \item Restituer la grille à l'auteur. Prendre connaissance de l'évaluation reçue.
    \item \textbf{Réécrire} votre lettre en intégrant les corrections pertinentes (version finale).
    \item Remettre la version finale au professeur pour notation et sélection pour le Tableau d'honneur.
\end{enumerate}

\par
\textbf{Phase 4 -- Communication par l'image (15 min) :}
\begin{enumerate}
    \setcounter{enumi}{9}
    \item En groupe-classe, observer les 3 meilleures lettres affichées. Identifier les choix de \textbf{mise en page} qui renforcent l'image professionnelle (en-tête avec logo fictif, interlignage, alignement de la signature à droite, objet visible).
    \item Discuter : \og{}En quoi la forme (l'image du texte) influence-t-elle le recruteur avant même la lecture ?\fg{}.
\end{enumerate}

\courssub{Ressources théoriques et méthodologiques}

\begin{definition}[Figures de style au programme (Stylistique/Rhétorique)]
Le programme de Première Technique prévoit l'étude de quatre figures de style. Deux sont des outils productifs majeurs pour la lettre de motivation ; les deux autres enrichissent la réception de textes.
\begin{itemize}
    \item \textbf{Énumération :} Suite de termes (noms, verbes, adjectifs) sans conjonction (asyndète) ou avec conjonction (polysyndète) pour montrer l'étendue, la diversité. \emph{Outil clé du paragraphe « Moi ».}
    \item \textbf{Gradation :} Mots disposés dans un ordre croissant (ou décroissant) d'intensité pour marquer la progression, l'évolution. \emph{Outil clé pour montrer la montée en compétences (paragraphe « Moi ») ou la projection (paragraphe « Nous »).}
    \item \textbf{Hyperbole :} Exagération volontaire pour frapper l'esprit (ex. : \og{}Une montagne de travail\fg{}). À repérer en réception ; à utiliser avec parcimonie en production (risque de manque de crédibilité).
    \item \textbf{Métonymie :} Remplacement d'un mot par un autre qui entretient avec lui un rapport de contiguïté (contenant/contenu, cause/effet, auteur/œuvre, lieu/activité). Ex. : \og{}La SODECOTON a décidé\fg{} (le lieu pour l'institution/dirigeants). \emph{Fréquente en réception (presse, discours).}
\end{itemize}
\end{definition}

\begin{definition}[Structure externe type de la lettre officielle (norme administrative camerounaise)]
La lettre officielle comprend \textbf{8 zones obligatoires} disposées dans l'ordre suivant :
\begin{enumerate}
    \item \textbf{Expéditeur} (coordonnées complètes : nom, adresse, téléphone, email) -- en haut à gauche.
    \item \textbf{Destinataire} (qualité, nom de la structure, adresse) -- en haut à droite, sous l'expéditeur.
    \item \textbf{Lieu et Date} -- à droite, sous le destinataire (ex. : \og{}Garoua, le 20 mai 2026\fg{}).
    \item \textbf{Objet} (en gras, centré ou aligné à gauche, précisant le poste/stage et la référence) -- ligne séparée.
    \item \textbf{Référence} (optionnel mais recommandé : \og{}Réf. : Votre avis n°... du...\fg{} ou \og{}Réf. : Convention Lycée/SODECOTON 2024\fg{}).
    \item \textbf{Corps de la lettre} (structure interne : Introduction / Développement / Conclusion).
    \item \textbf{Formule de politesse finale} (adaptée au destinataire, utilisant l'impératif de politesse).
    \item \textbf{Signature} (manuscrite) précédée du nom tapé -- en bas à droite.
\end{enumerate}
\end{definition}

\begin{propriete}[Structure interne : Plan « Vous -- Moi -- Nous »]
\begin{itemize}
    \item \textbf{Vous (Introduction) :} Accroche -- Montrer qu'on connaît l'entreprise, l'offre, ses valeurs. Citer la source de l'info. Exprimer l'intérêt.
    \item \textbf{Moi (Développement) :} Argumentation -- Présenter son parcours (formation, stages, projets), ses compétences techniques (\emph{hard skills}) et qualités humaines (\emph{soft skills}). Utiliser \textbf{énumération} et \textbf{gradation}. Éviter la simple liste de CV ; \emph{lier} compétences et besoins de l'entreprise.
    \item \textbf{Nous (Conclusion) :} Projection -- Exprimer la volonté de contribuer, d'apprendre, de s'investir. Demander l'entretien. Formule de politesse impérative.
\end{itemize}
\end{propriete}

\begin{definition}[Figures de style productives : Énumération et Gradation]
\begin{itemize}
    \item \textbf{Énumération (asyndète) :} Suite de termes séparés par des virgules. \emph{Ex. : « Ma formation m'a permis d'acquérir des compétences en soudure TIG, en lecture de plans ISO, en maintenance préventive et en diagnostic de pannes hydrauliques. »}
    \item \textbf{Gradation (croissante) :} Verbes ou noms d'intensité croissante. \emph{Ex. : « J'ai pu \textbf{observer}, puis \textbf{comprendre}, enfin \textbf{agir} en autonomie sur la chaîne de convoyeur. »} ou \emph{Ex. : « Une rigueur \textbf{initiale}, \textbf{consolidée} par la pratique, \textbf{maîtrisée} aujourd'hui. »}
\end{itemize}
\end{definition}

\begin{propriete}[Mode impératif : valeur de politesse]
Dans les formules de politesse finales, l'impératif présent (2e personne du singulier ou pluriel) exprime une \textbf{injonction courtoise}, non un ordre. Le sujet \og{}vous\fg{} est souvent sous-entendu.
\begin{itemize}
    \item \og{}\textbf{Veuillez} agréer, Monsieur le Directeur, l'expression de ma considération distinguée.\fg{} (Impératif 2e pers. pl. de \emph{vouloir} -- forme figée).
    \item \og{}\textbf{Recevez}, Madame, Monsieur, mes salutations respectueuses.\fg{} (Impératif 2e pers. pl. de \emph{recevoir}).
    \item \og{}\textbf{Je vous prie de} bien vouloir accepter...\fg{} (Structure infinitive, équivalent poli très usité, évitant l'impératif direct).
\end{itemize}
\textbf{Attention :} Ne pas confondre avec l'impératif d'ordre (\og{}Fermez la porte\fg{}). Ici, c'est un \emph{rituel codifié} de la correspondance formelle.
\end{propriete}

\begin{methode}[Démarche de rédaction en 5 étapes]
\begin{enumerate}
    \item \textbf{Analyser l'offre} : Identifier l'entreprise, le service visé, les compétences techniques attendues, les mots-clés (HSE, GMAO, maintenance prédictive, etc.).
    \item \textbf{Structurer le brouillon (Plan Vous/Moi/Nous)} : Noter les arguments par paragraphe. Prévoir l'énumération (compétences) et la gradation (progression/projection).
    \item \textbf{Rédiger la structure externe} : Appliquer la checklist des 8 zones (coordonnées, destinataire, date, objet, référence).
    \item \textbf{Rédiger le corps} : Un paragraphe par mouvement. Vocabulaire technique précis. Phrases complexes. Intégrer figures de style.
    \item \textbf{Relire et corriger (OFOPS)} : Vérifier orthographe, formulation, organisation, ponctuation, syntaxe (accords, impératif politesse). Soigner la mise en page (image).
\end{enumerate}
\end{methode}

\begin{aretenir}[Grille OFOPS -- Évaluation sur 20 points]
\begin{center}
\begin{tabularx}{\linewidth}{|l|X|c|}\hline
\textbf{Critère} & \textbf{Descripteurs (5 pts chacun)} & \textbf{Note} \\ \hline
\textbf{O -- Orthographe} & 0-1 : > 5 fautes / 2-3 : 2-5 fautes / 4-5 : 0-1 faute (lexicale, grammaticale, accentuation). & /5 \\ \hline
\textbf{F -- Formulation} & 0-1 : Phrases maladroites, vocabulaire pauvre / 2-3 : Correct, vocabulaire courant / 4-5 : Précis, vocabulaire pro, figures de style utilisées. & /5 \\ \hline
\textbf{O -- Organisation} & 0-1 : Structure absente / 2-3 : Structure partielle (manque Vous/Moi/Nous) / 4-5 : Structure externe + interne parfaites, paragraphes logiques. & /5 \\ \hline
\textbf{P -- Ponctuation} & 0-1 : Absente ou fautive / 2-3 : Principaux signes ok / 4-5 : Maîtrisée (virgules énumération, deux-points, point-virgule). & /5 \\ \hline
\textbf{S -- Syntaxe} & 0-1 : Phrases agrammaticales / 2-3 : Phrases simples correctes / 4-5 : Phrases complexes, accords parfaits, impératif politesse ok. & /5 \\ \hline
\end{tabularx}
\end{center}
\textbf{Total : /20}. Seuil de réussite : 12/20. Mention « Tableau d'honneur » : $\ge$ 17/20.
\end{aretenir}

\begin{savaistu}[SODECOTON : acteur clé du Septentrion]
La Société de Développement du Coton (SODECOTON), créée en 1974, est une entreprise publique à capital majoritairement étatique. Elle encadre plus de 200 000 producteurs dans les régions du Nord, de l'Adamaoua et de l'Extrême-Nord. L'usine d'égrenage de Garoua (métiers Continental Eagle) est un site industriel majeur où la maintenance prédictive (analyse vibratoire, thermographie) est cruciale pour la disponibilité des équipements pendant la campagne d'égrenage (novembre--mars). Le stage y offre une immersion en environnement industriel réel, exigence du référentiel Maintenance industrielle.
\end{savaistu}

\begin{exoresolu}[Activité d'intégration : Candidature stage SODECOTON]
\textbf{Énoncé -- Situation et consignes}
\par
\emph{Contexte :} 15 mai 2026. Élève de 1re Maintenance industrielle (Lycée Technique Garoua). Partenariat SODECOTON (BP 302 Garoua). Stage été 2026 (8 semaines). Date limite : 30 mai 2026.
\par
\emph{Mission :} 1) Analyser Lettre A (refusée) vs Lettre B (acceptée) $\rightarrow$ Grille de réussite. 2) Rédiger sa lettre (Données : Djibrilla Amadou, 14,5/20, stage CAMRAIL, Projet graissage 18/20). 3) Évaluation pairs (Grille OFOPS) + Réécriture. 4) Analyse communication par l'image (Tableau d'honneur).

\tcblower

\textbf{Solution détaillée -- Étapes commentées}

\par
\textbf{Étape 1 -- Analyse des lettres modèles (Phase 1)}

\par
\textbf{Lettre A (Extrait -- Candidat refusé 2025) :}
\par
\emph{garoua le 25 mai 2025}
\par
\emph{Monsieur le Directeur de la SODECOTON}
\par
\emph{Objet : Demande de stage}
\par
\emph{Monsieur,}
\par
\emph{Je veux faire un stage chez vous cet été. Je suis en 1ère maintenance. J'aime bien les machines. L'an dernier j'ai fait un stage à la CAMRAIL c'était bien. Je suis fort, courageux et je viens à l'heure. J'espère que vous me prendrez.}
\par
\emph{Salutations.}
\par
\emph{Amadou D.}

\par
\textbf{Analyse critique (grille groupe) :}
\begin{itemize}
    \item \textbf{Structure externe :} Manque coordonnées expéditeur, adresse destinataire complète, date mal formatée (lieu seul), pas de référence à l'avis, objet trop vague, formule de politesse inexistante (juste \og{}Salutations\fg{}), signature incomplète.
    \item \textbf{Structure interne :} Pas de plan Vous/Moi/Nous. Introduction inexistante. Développement : phrases télégraphiques, pas d'argumentation. Conclusion : supplique (\og{}J'espère que vous me prendrez\fg{}).
    \item \textbf{Langage :} Familier (\og{}Je veux\fg{}, \og{}c'était bien\fg{}), pas de figure de style, pas d'impératif de politesse. Fautes d'orthographe/ponctuation (majuscules, points).
    \item \textbf{Mise en page :} Aléatoire, pas de paragraphes, alignement approximatif.
\end{itemize}

\par
\textbf{Lettre B (Extrait -- Candidat accepté 2025 -- Maintenance) :}
\par
\emph{\textsc{Tchinda Kevin}}
\par
\emph{Quartier Commercial, Rue 5, Porte 12}
\par
\emph{677 12 34 56 -- kevin.tchinda@email.cm}
\par
\vspace{0.2cm}
\par
\emph{À l'attention de Monsieur le Directeur des Ressources Humaines}
\par
\emph{SODECOTON -- BP 302 Garoua}
\par
\vspace{0.2cm}
\par
\emph{Garoua, le 22 mai 2025}
\par
\vspace{0.2cm}
\par
\emph{\textbf{Objet : Candidature au stage académique d'été 2025 -- Filière Maintenance industrielle (Réf. : Avis n° 2025/004/SOD/RH)}}
\par
\vspace{0.3cm}
\par
\emph{Monsieur le Directeur,}
\par
\emph{Votre entreprise, leader incontesté du secteur cotonnier au Cameroun, incarne l'excellence industrielle par la modernisation continue de ses unités d'égrenage. L'avis de stage publié sur le site de la SODECOTON et relayé par mon établissement, le Lycée Technique de Garoua, retient toute mon attention car il offre l'opportunité unique de confronter mes acquis scolaires à la réalité de la maintenance prédictive sur métiers Continental Eagle.}
\par
\emph{Actuellement élève de Première Maintenance industrielle (moyenne 15,2/20, major de promotion), j'ai consolidé mes bases théoriques en technologie mécanique, électrotechnique et automatismes. Mon stage d'initiation de deux semaines à l'atelier CAMRAIL de Garoua (juillet 2024) m'a permis d'\textbf{observer} les opérations de maintenance préventive sur locomotives, de \textbf{comprendre} la gestion des pièces de rechange (GMAO) et d'\textbf{agir} en binôme sur le remplacement de segments de piston -- une \textbf{gradation} qui illustre ma progression vers l'autonomie. Par ailleurs, mon projet de fin d'année, « Système de graissage automatique pour convoyeur », noté 17/20, a mobilisé des compétences en \textbf{soudure TIG, lecture de plans normalisés, programmation API Siemens LOGO! et dimensionnement de pompes} -- \textbf{énumération} technique ciblée.}
\par
\emph{Rigoureux, respectueux des normes HSE et doté d'un bon esprit d'équipe, je souhaite vivement intégrer votre service Maintenance pour contribuer, à mon niveau, à la disponibilité de l'outil de production tout en perfectionnant ma pratique du diagnostic vibratoire.}
\par
\emph{\textbf{Je vous prie d'agréer, Monsieur le Directeur, l'expression de ma considération distinguée.}}
\par
\vspace{0.5cm}
\par
\emph{\hfill Kevin TCHINDA}

\par
\textbf{Analyse critique (grille groupe) :}
\begin{itemize}
    \item \textbf{Structure externe :} Complète (8 zones). Objet précis avec référence. Date bien formatée.
    \item \textbf{Structure interne :} \textbf{Vous} (connaissance entreprise, source info, motivation ciblée) -- \textbf{Moi} (parcours, stage, projet, figures de style) -- \textbf{Nous} (projection, apport mutuel, demande implicite entretien).
    \item \textbf{Langage :} Soutenu, vocabulaire technique précis (GMAO, API, HSE, diagnostic vibratoire). \textbf{Gradation} (\og{}observer, comprendre, agir\fg{}). \textbf{Énumération} asyndète (compétences projet). \textbf{Impératif de politesse} (\og{}Je vous prie d'agréer...\fg{}).
    \item \textbf{Mise en page :} Aérée, alignée, police lisible, objet en gras, signature à droite.
\end{itemize}

\par
\textbf{Critères de réussite dégagés par la classe (affichés) :}
\begin{enumerate}
    \item Respect strict de la structure externe (8 zones).
    \item Plan Vous/Moi/Nous visible (3 paragraphes distincts).
    \item Vocabulaire technique pertinent + au moins 1 énumération + 1 gradation.
    \item Formule de politesse finale avec impératif (Veuillez / Je vous prie d'agréer).
    \item Mise en page professionnelle (marges, interligne, alignement, zéro faute OFOPS).
\end{enumerate}

\par
\textbf{Étape 2 -- Rédaction individuelle guidée (Phase 2) -- Production de l'élève \textsc{Djibrilla Amadou}}

\par
\textbf{Plan de rédaction (brouillon) :}
\begin{itemize}
    \item \textbf{Vous :} SODECOTON = leader coton, modernisation usine Garoua. Avis vu au lycée + site web. Intérêt pour maintenance prédictive sur métiers Continental Eagle.
    \item \textbf{Moi :} 1ère Maintenance (moy 14,5, 3/42). Stage CAMRAIL 2025 (soudure, lecture plans, maintenance préventive loco). Projet fin année : graissage auto convoyeur (18/20) $\rightarrow$ compétences : soudure, plans, API, pneumatique. Qualités : rigueur, HSE, équipe. \textbf{Gradation :} \og{}découvrir, approfondir, maîtriser\fg{} (ou \og{}observer, analyser, intervenir\fg{}). \textbf{Énumération :} compétences projet.
    \item \textbf{Nous :} Envie d'intégrer service Maintenance Garoua. Apporter motivation + acquis. Demander entretien. Formule politesse.
\end{itemize}

\par
\textbf{Version finale corrigée (après auto-évaluation OFOPS et réécriture) :}

\begin{center}
\begin{minipage}{\linewidth}
\setlength{\parindent}{0pt}
\setlength{\parskip}{4pt}
\textsc{Djibrilla Amadou}\par
Quartier Lamordé, Rue 12, Porte 45 -- Garoua\par
Tél : 6 99 88 77 66 -- Email : djibrilla.amadou@email.cm\par
\vspace{0.5cm}
\hfill
\begin{minipage}{0.6\linewidth}
\raggedleft
À l'attention de \textbf{Monsieur le Directeur des Ressources Humaines}\par
\textbf{SODECOTON} -- BP 302 Garoua\par
\vspace{0.3cm}
Garoua, le 20 mai 2026
\end{minipage}
\vspace{0.5cm}
\textbf{Objet : Candidature au stage académique d'été 2026 -- Filière Maintenance industrielle (Réf. : Avis de stage Session 2026 / Convention Lycée/SODECOTON)}\par
\vspace{0.5cm}
Monsieur le Directeur,\par
\vspace{0.2cm}
Votre société, pilier de l'économie nationale et fer de lance de l'industrialisation du Septentrion, engage depuis plusieurs années une politique ambitieuse de modernisation de son outil de production, notamment au sein de l'usine d'égrenage de Garoua. L'offre de stage académique portée à ma connaissance par l'administration de mon établissement, le Lycée Technique de Garoua, et consultée sur votre site institutionnel, suscite en moi un vif intérêt : elle représente l'occasion privilégiée de découvrir la maintenance prédictive appliquée aux métiers à égrener Continental Eagle, au cœur même de votre processus de transformation du coton-graine.\par
\vspace{0.2cm}
Actuellement élève de Première Technique, filière Maintenance industrielle (option Systèmes automatisés), je figure parmi les premiers de ma promotion avec une moyenne annuelle de 14,5/20. Ma formation m'a doté de solides bases en technologie mécanique, en électrotechnique et en automatismes programmables. Une première immersion professionnelle de deux semaines à l'atelier de maintenance de la CAMRAIL (Garoua, juillet 2025) m'a permis d'\textbf{observer} les procédures de maintenance préventive sur moteurs diesel, d'\textbf{analyser} les schémas hydrauliques et pneumatiques, et d'\textbf{intervenir} en binôme sur le remplacement d'organes de transmission -- une \textbf{gradation} témoignant de mon passage de l'observation à l'action technique encadrée. Par ailleurs, la réalisation de mon projet de fin d'année, « Conception d'un système de graissage automatique pour chaîne de convoyeur », sanctionnée par la note de 18/20, a mobilisé des compétences concrètes en \textbf{soudure TIG et MIG, lecture de plans mécaniques normalisés ISO, programmation d'automates Siemens LOGO!, dimensionnement de circuits pneumatiques et rédaction de modes opératoires} -- \textbf{énumération} qui reflète ma polyvalence technique.\par
\vspace{0.2cm}
Conscient des exigences de sécurité (HSE) et de rigueur qui régissent vos installations, je suis animé par la volonté de mettre mon dynamisme, ma capacité d'adaptation et mon sens du collectif au service de la performance de votre service Maintenance. Intégrer vos équipes serait pour moi l'opportunité de \textbf{consolider} mes acquis, de \textbf{perfectionner} ma pratique du diagnostic de pannes et de \textbf{construire} mon projet professionnel d'ingénieur de maintenance -- seconde \textbf{gradation} marquant la projection vers l'excellence.\par
\vspace{0.2cm}
Dans l'attente d'une réponse que j'espère favorable, \textbf{je vous prie d'agréer, Monsieur le Directeur, l'expression de ma considération distinguée.}\par
\vspace{1cm}
\hfill \textsc{Djibrilla Amadou}
\end{minipage}
\end{center}

\par
\textbf{Commentaire du corrigé (justifications pédagogiques) :}
\begin{itemize}
    \item \textbf{Zone 1 (Expéditeur) :} Complète, alignée à gauche. Email professionnel (prénom.nom).
    \item \textbf{Zone 2 (Destinataire) :} Qualité exacte (\og{}Directeur des Ressources Humaines\fg{}), nom structure, BP. Aligné à droite (simulé par \texttt{\textbackslash hfill} + minipage).
    \item \textbf{Zone 3 (Lieu/Date) :} \og{}Garoua, le 20 mai 2026\fg{} (avant limite 30 mai). À droite.
    \item \textbf{Zone 4 (Objet) :} En gras, précis : stage + filière + référence à l'avis/convention. Mots-clés pour le tri RH.
    \item \textbf{Paragraphe 1 (VOUS) :} Accroche valorisante (\og{}pilier, fer de lance\fg{} -- \emph{métaphore/métonymie} acceptée). Preuve de recherche (site, lycée). Cible technique précise (maintenance prédictive, Continental Eagle). \textbf{Objectif :} montrer qu'on ne candidate pas au hasard.
    \item \textbf{Paragraphe 2 (MOI -- 1er bloc) :} Situation scolaire (niveau, filière, rang). Stage CAMRAIL : \textbf{gradation} \og{}observer / analyser / intervenir\fg{} (verbes d'action croissants). Projet fin d'année : note + \textbf{énumération} asyndète de 5 compétences techniques ciblées (soudure, plans, API, pneumatique, rédaction). \textbf{Vocabulaire technique} pertinent.
    \item \textbf{Paragraphe 3 (MOI/Nous -- Qualités + Projection) :} Qualités (HSE, rigueur, équipe) liées au contexte. Projection : \textbf{gradation} \og{}consolider / perfectionner / construire\fg{} (progression vers l'autonomie professionnelle). Transition vers \og{}Nous\fg{}.
    \item \textbf{Conclusion (NOUS) :} Formule type avec \textbf{impératif de politesse} \og{}je vous prie d'agréer\fg{} (figé, correct). Signature manuscrite simulée (nom en capitales à droite).
    \item \textbf{Mise en page :} Marges respectées, interligne 1,5 (simulé par \texttt{\textbackslash parskip}), paragraphes séparés, alignement justifié, objet en gras. \textbf{Zéro faute OFOPS visée.}
\end{itemize}

\par
\textbf{Étape 3 -- Simulation évaluation par les pairs (Grille OFOPS) -- Exemple sur la lettre ci-dessus}

\begin{center}
\begin{tabularx}{\linewidth}{|l|X|c|}\hline
\textbf{Critère} & \textbf{Commentaire évaluateur (binôme)} & \textbf{Note} \\ \hline
\textbf{O -- Orthographe} & Aucune faute détectée. Accords parfaits (compétences \textbf{mobilisé\textcolor{popgreen}{e}}, gradation \textbf{perm\textcolor{popgreen}{is}}). Majuscules correctes (SODECOTON, CAMRAIL, ISO). & 5/5 \\ \hline
\textbf{F -- Formulation} & Registre soutenu constant. Vocabulaire technique précis (API, pneumatique, HSE présents). Métaphore \og{}fer de lance\fg{} élégante. Gradations et énumération bien intégrées, non forcées. & 5/5 \\ \hline
\textbf{O -- Organisation} & Structure externe : 8/8 zones présentes. Structure interne : 3 paragraphes nets (Vous/Moi/Nous). Transitions fluides. Objet clair. & 5/5 \\ \hline
\textbf{P -- Ponctuation} & Virgules énumération parfaites. Deux-points après objet. Point-virgule absent mais non nécessaire. Ponctuation forte (points) aux fins de phrases complexes. & 5/5 \\ \hline
\textbf{S -- Syntaxe} & Phrases complexes (subordonnées relatives, complétives). Impératif politesse correct (\og{}je vous prie d'agréer\fg{}). Accords participes passés (mobilisée, permis -- COD avant). & 5/5 \\ \hline
\end{tabularx}
\end{center}
\par
\textbf{Total : 20/20} -- \textbf{Sélectionnée pour le Tableau d'honneur.}

\par
\textbf{Étape 4 -- Bilan communication par l'image (Phase 4)}
\par
Les trois lettres affichées présentent toutes : un en-tête avec un \emph{logo fictif} (cercle + sigle SODECOTON dessiné au stylo ou inséré image), une police unique (Times New Roman 12), des marges symétriques, l'objet en gras souligné, la signature encadrée d'un trait fin. La classe conclut : \og{}La forme soignée signale le sérieux du candidat avant la lecture ; une lettre brouillonne suggère un technicien négligent, ce qui est rédhibitoire en maintenance où la précision sauve des vies.\fg{}.

\end{exoresolu}

\courssub{Bilan de l'activité et remédiation}
\par
Cette activité d'intégration a permis de \textbf{valider l'acquisition opérationnelle} des savoirs de la leçon N21 dans une situation de communication réelle, ancrée dans le tissu économique local (SODECOTON, Garoua). Les points suivants ont été mis en évidence :

\begin{enumerate}
    \item \textbf{Transfert effectif des connaissances :} Les élèves ont spontanément mobilisé la structure externe (8 zones), le plan Vous/Moi/Nous, les figures de style (gradation, énumération) et l'impératif de politesse sans réexplication du professeur, prouvant l'intériorisation des modèles.
    \item \textbf{Rôle décisif de l'analyse comparative :} Le face-à-face Lettre A / Lettre B a agi comme un \emph{révélateur} des critères implicites de la correspondance professionnelle. Les élèves ont eux-mêmes formulé la grille de réussite (5 critères), ce qui renforce l'appropriation.
    \item \textbf{Exigence de la langue technique :} L'obligation d'intégrer du vocabulaire de spécialité (API, pneumatique, HSE, GMAO, Continental Eagle) a forcé les élèves à puiser dans leurs cours techniques (Maintenance, Électrotechnique), réalisant ainsi l'interdisciplinarité visée par l'enseignement technique.
    \item \textbf{Efficacité de la grille OFOPS :} L'évaluation par les pairs a transformé les élèves en \emph{lecteurs critiques}. Le codage couleur (vert/orange/rose) a rendu la révision concrète. La réécriture systématique a fait progresser la note moyenne de la classe de 11,2/20 (brouillon) à 15,8/20 (version finale).
    \item \textbf{Communication par l'image :} La phase finale a montré que la \emph{mise en page} n'est pas un ornement mais un \emph{argument de compétence} : en maintenance, la rigueur graphique signale la rigueur technique.
    \item \textbf{Ancrage camerounais :} L'utilisation de données réelles (BP 302 Garoua, filières SODECOTON, convention lycée, date limite 30 mai) a suscité un engagement fort : 87\% de la classe a déposé son dossier réel auprès du service RH du lycée pour la session 2026.
\end{enumerate}

\par
\textbf{Pistes de remédiation pour les élèves en difficulté (note < 12/20) :}
\begin{itemize}
    \item Atelier ciblé \og{}Formules de politesse : impératif vs indicatif\fg{} (exercices de transformation).
    \item Fiche mémo \og{}Structure externe : checklist 8 zones\fg{} à coller dans le cahier.
    \item Entraînement guidé à l'énumération/gradation à partir de leur CV technique (transformer une liste CV en phrases de lettre).
    \item Relecture à voix haute pour repérer les ruptures de ton (familier vs soutenu).
\end{itemize}

\par
Cette activité constitue ainsi une \textbf{évaluation certificative} naturelle des compétences \og{}Produire une lettre officielle\fg{} et \og{}Communiquer par l'image\fg{}, tout en préparant concrètement l'insertion professionnelle des élèves de Première Technique.

\newpage

\coursec{Méthodes et exercices résolus}

\coursec{V : Produire une lettre officielle (lettre de motivation)}

\courssub{Méthodes et exercices résolus}

\begin{methode}[Identifier un texte fonctionnel et les caractéristiques de la lettre de motivation]
    \textbf{Objectif :} Distinguer les textes fonctionnels des textes littéraires et repérer les spécificités de la lettre de motivation.
    \begin{enumerate}
        \item \textbf{Analyser la finalité :} Le texte vise-t-il une action concrète (obtenir un emploi, un stage, une inscription) ? $\rightarrow$ Texte fonctionnel.
        \item \textbf{Repérer le destinataire :} Une personne morale (entreprise, administration, établissement) ou physique identifiée (DRH, Directeur).
        \item \textbf{Vérifier la structure codifiée :} Présence obligatoire d'en-têtes (expéditeur, destinataire, lieu, date, objet), formules d'appel et de politesse codifiées, signature.
        \item \textbf{Identifier le ton :} Formel, respectueux, persuasif, sans familiarité ni émotion excessive.
        \item \textbf{Conclure :} Si tous ces critères sont réunis, c'est une lettre officielle (ici, de motivation).
    \end{enumerate}
\end{methode}

\begin{exoresolu}[Analyse de la nature d'un texte]
    \textbf{Énoncé :} Vous trouverez ci-dessous deux extraits. Pour chacun, dites s'il s'agit d'un texte fonctionnel relevant de la lettre officielle. Justifiez votre réponse en vous appuyant sur trois critères distincts.

    \textbf{Texte A :} « Cher ami, Je t'écris pour te dire que je suis très content de mon nouveau poste. L'ambiance est super et le patron est cool. Viens me voir quand tu veux. Bises, Paul. »

    \textbf{Texte B :} « Yaoundé, le 15 octobre 2025. Monsieur le Directeur des Ressources Humaines, Société CAMWATER, BP 1234 Yaoundé. Objet : Candidature au poste de Technicien en maintenance. Monsieur, Par la présente, je me permets de vous adresser ma candidature... Je vous prie d'agréer, Monsieur, l'expression de mes salutations distinguées. Signature. »

    \tcblower
    \textbf{Solution détaillée :}

    \textbf{1. Analyse du Texte A :}
    \begin{itemize}
        \item \textbf{Finalité :} Information amicale, partage d'émotion. Pas de visée administrative ou professionnelle formelle.
        \item \textbf{Destinataire :} « Cher ami », tutoiement (« te dire », « tu veux »). Relation privée, non hiérarchique.
        \item \textbf{Structure :} Absence d'en-tête complet (pas d'adresse destinataire, pas d'objet, pas de lieu/date formels), formules de politesse familières (« Bises »).
        \item \textbf{Ton :} Familier (« super », « cool »).
    \end{itemize}
    \textbf{Conclusion :} \textbf{Non}, ce n'est pas une lettre officielle. C'est une \cle{lettre personnelle (familière)}.

    \textbf{2. Analyse du Texte B :}
    \begin{itemize}
        \item \textbf{Finalité :} Obtenir un entretien d'embauche (action concrète, candidature).
        \item \textbf{Destinataire :} « Monsieur le Directeur des Ressources Humaines », Société CAMWATER. Relation professionnelle, hiérarchique (vouvoiement implicite).
        \item \textbf{Structure :} En-tête complet (Lieu/Date, Expéditeur implicite, Destinataire complet, Objet), Formule d'appel (« Monsieur »), Corps structuré, Formule de politesse longue codifiée (« Je vous prie d'agréer... »), Signature.
        \item \textbf{Ton :} Formel, respectueux (« je me permets », « salutations distinguées »).
    \end{itemize}
    \textbf{Conclusion :} \textbf{Oui}, c'est une \cle{lettre officielle de motivation}. Elle respecte tous les codes du genre épistolaire administratif/professionnel.
\end{exoresolu}

\begin{methode}[Maîtriser la structure externe (mise en page) de la lettre officielle]
    \textbf{Objectif :} Respecter scrupuleusement les normes de présentation (normes AFNOR / usage camerounais) pour une lisibilité immédiate.
    \begin{enumerate}
        \item \textbf{En-tête expéditeur (haut gauche) :} Prénom NOM, Adresse complète (Quartier, Ville, BP), Téléphone, Email. \emph{Astuce : Aligné à gauche.}
        \item \textbf{En-tête destinataire (sous expéditeur, à gauche) :} Qualité + Nom (ou Service), Nom de la structure, Adresse structure (BP, Ville).
        \item \textbf{Lieu et Date (à droite, sous destinataire) :} Ville, le [chiffres] [mois en toutes lettres] [année]. Ex: \emph{Yaoundé, le 15 octobre 2025}.
        \item \textbf{Objet (à gauche, 2 lignes sous la date, en gras/majuscules) :} Objet : [Nature de la demande] + [Poste/Concours visé] + [Référence annonce si existe].
        \item \textbf{Formule d'appel (à gauche, 1 ligne sous objet) :} Monsieur, / Madame, / Monsieur le Directeur, (Pas de nom propre, pas de « Cher »).
        \item \textbf{Corps de la lettre :} Aligné justifié, paragraphes séparés par une ligne blanche (Introduction, Développement, Conclusion).
        \item \textbf{Formule de politesse (fin, à gauche) :} Longue, adaptée au destinataire (Voir fiche «À retenir» ci-après).
        \item \textbf{Signature (sous formule, à droite) :} Prénom NOM manuscrit (si papier) ou tapé. Mention « Pièces jointes : » (aligné gauche, sous signature).
    \end{enumerate}
\end{methode}

\begin{aretenir}[Formules de politesse types (Bac)]
    \textbf{Supérieur hiérarchique / Administration :} « Je vous prie d'agréer, Monsieur le Directeur / Madame la Directrice, l'expression de mes salutations distinguées. »
    \textbf{DRH / Employeur privé :} « En vous remerciant de l'attention portée à ma candidature, je vous prie d'agréer, Madame, Monsieur, mes salutations distinguées. »
    \textbf{Interdit :} « Cordialement », « Amicalement », « Bien à vous » (trop courts/informels pour le Bac).
\end{aretenir}

\begin{exoresolu}[Mise en page complète d'une lettre]
    \textbf{Énoncé :} Marie NGOUMOU, habitante au quartier Mvog-Ada, BP 888 Yaoundé, Tél : 6XX XX XX XX, Email : m.ngoumou@email.cm, candidate au poste d'Assistante Comptable (Réf : AC-2025) à la Société CAMPOST (BP 234 Yaoundé). Nous sommes le 20 novembre 2025. Rédigez \textbf{uniquement la structure externe} (en-têtes, lieu/date, objet, formule d'appel, formule de politesse, signature, mentions PJ) sur une page blanche. Le corps de la lettre est remplacé par la mention \textit{[Corps de la lettre]}.

    \tcblower
    \textbf{Solution détaillée :}
    \begin{center}
    \begin{tabularx}{\linewidth}{|X|}\hline
    \textbf{MARIE NGOUMOU} \\
    Quartier Mvog-Ada, BP 888 Yaoundé \\
    Tél : 6XX XX XX XX \quad Email : m.ngoumou@email.cm \\ \hline
    \textbf{Monsieur le Directeur Général} \\
    \textbf{Société CAMPOST} \\
    BP 234 Yaoundé \\ \hline
    \hfill Yaoundé, le 20 novembre 2025 \\ \hline
    \textbf{Objet : Candidature au poste d'Assistante Comptable (Réf : AC-2025)} \\ \hline
    \textbf{Monsieur,} \\ \hline
    \textit{[Corps de la lettre : Introduction, Développement, Conclusion]} \\ \hline
    \textit{En vous remerciant de l'attention portée à ma candidature, je vous prie d'agréer, Monsieur, l'expression de mes salutations distinguées.} \\ \hline
    \hfill \textbf{Marie NGOUMOU} \\ \hline
    \textbf{Pièces jointes :} CV, Copie diplômes, Attestations de travail, Extrait d'acte de naissance. \\ \hline
    \end{tabularx}
    \end{center}
    \textbf{Justifications :}
    \begin{itemize}
        \item Alignement standard : Expéditeur/Destinataire/Objet/Appel/Corps/PJ à \textbf{gauche}. Lieu/Date/Signature à \textbf{droite}.
        \item Date en toutes lettres (mois), Objet en gras avec référence.
        \item Formule de politesse complète (type employeur privé).
        \item Signature Prénom NOM (majuscules au nom).
        \item Liste des PJ standard pour un dossier de candidature.
    \end{itemize}
\end{exoresolu}

\begin{methode}[Rédiger l'introduction et la conclusion (le « cadre » persuasif)]
    \textbf{Objectif :} Accrocher le lecteur (introduction) et provoquer l'action (entretien) en conclusion.
    \begin{enumerate}
        \item \textbf{Introduction (3 mouvements) :}
            \begin{enumerate}
                \item \textbf{La source (le « déclencheur ») :} « Suite à votre annonce parue dans \emph{Cameroon Tribune} du 10/11/2025... », « Sur les conseils de M. X... », « Actuellement à la recherche d'un stage... » (candidature spontanée).
                \item \textbf{L'objet précis :} « ...je me permets de vous adresser ma candidature au poste de... » / « ...je sollicite un stage de fin d'études... ».
                \item \textbf{L'accroche (la « valeur ajoutée » immédiate) :} Une phrase qui résume le profil : « Mon BTS en Génie Civil couplé à deux ans d'expérience sur chantier me permet d'être opérationnel immédiatement. »
            \end{enumerate}
        \item \textbf{Conclusion (2 mouvements) :}
            \begin{enumerate}
                \item \textbf{La disponibilité / L'appel à l'action :} « Disponible immédiatement pour un entretien à votre convenance, je me tiens à votre disposition pour tout complément d'information. »
                \item \textbf{La formule de politesse (voir fiche précédente).}
            \end{enumerate}
        \item \textbf{Reflexe « Je/Nous » :} Utiliser « Je » (candidature individuelle). Éviter « Nous » sauf si candidature groupée (rare).
        \item \textbf{Reflexe « Conditionnel de politesse » :} « Je \textbf{veuillez} » (impératif \textbf{interdit} dans le corps), « Je \textbf{souhaiterais} », « Je \textbf{serais} honoré(e) ».
    \end{enumerate}
\end{methode}

\begin{exoresolu}[Rédaction introduction et conclusion]
    \textbf{Énoncé :} Vous êtes Armand KAMGA, titulaire d'un Baccalauréat F4 (Électrotechnique) et d'un BTS Maintenance des Systèmes Automatisés. Vous répondez à une annonce de la SONEL (ENEO) pour un poste de « Technicien de Maintenance Réseau » (Réf : TMR-2025), parue sur le site web de l'entreprise le 5 mars 2025. Rédigez l'\textbf{introduction} et la \textbf{conclusion} de votre lettre.

    \tcblower
    \textbf{Solution détaillée :}

    \textbf{Introduction :}
    \begin{quote}
    \textit{Suite à votre avis de recrutement publié sur le site internet de la SONEL en date du 5 mars 2025, je me permets de vous adresser ma candidature au poste de Technicien de Maintenance Réseau (Réf : TMR-2025). Fort de mon BTS en Maintenance des Systèmes Automatisés et de mon stage de fin d'études effectué au sein de la direction régionale du Littoral où j'ai participé à la maintenance préventive des postes HTA/BT, je suis prêt à mettre mes compétences techniques et ma rigueur au service de la continuité de votre réseau de distribution.}
    \end{quote}
    \textbf{Analyse :} Source précise (site, date) $\rightarrow$ Objet précis (poste + réf) $\rightarrow$ Accroche (Diplôme + Expérience concrète + Qualité : rigueur/continuité).

    \textbf{Conclusion :}
    \begin{quote}
    \textit{Disponible dès à présent pour un entretien à votre convenance, je me tiens à votre entière disposition pour vous exposer de vive voix ma motivation et mes aptitudes. En vous remerciant de l'attention que vous voudrez bien porter à ma candidature, je vous prie d'agréer, Monsieur le Directeur, l'expression de mes salutations distinguées.}
    \end{quote}
    \textbf{Analyse :} Disponibilité immédiate + Proposition d'entretien (action) $\rightarrow$ Formule de politesse adaptée (Directeur).
\end{exoresolu}

\begin{methode}[Construire le corps de la lettre : Argumentation et outils langagiers]
    \textbf{Objectif :} Démontrer l'adéquation Profil/Poste (Compétences = Besoins) en utilisant le registre formel et des figures de style pour persuader.
    \begin{enumerate}
        \item \textbf{Structure en 2 ou 3 paragraphes distincts (idées forces) :}
            \begin{itemize}
                \item \textbf{Paragraphe 1 : Les Compétences Techniques (Savoir-faire).} Diplômes, stages, expériences, maîtrise d'outils/logiciels/machines. \emph{Mots-clés de l'annonce.}
                \item \textbf{Paragraphe 2 : Les Qualités Humaines / Transverses (Savoir-être).} Rigueur, esprit d'équipe, adaptation, sens du service, langues, permis. Preuves par l'exemple (ex: « Délégué de promo $\rightarrow$ leadership »).
                \item \textbf{Paragraphe 3 (optionnel) : La Connaissance de l'entreprise / Motivation spécifique.} Pourquoi EUX ? (Valeurs, projets, notoriété). Preuve qu'on ne fait pas du « publipostage ».
            \end{itemize}
        \item \textbf{Outils langagiers obligatoires (Programme) :}
            \begin{itemize}
                \item \textbf{Mode Impératif (Valeur incitative/Conseil) :} Uniquement dans la conclusion pour l'appel à l'action (\textit{« Veuillez agréer... », « Recevez... »}). \textbf{Interdit dans le corps} (trop directif/agressif).
                \item \textbf{Conditionnel (Politesse/Hypothèse) :} \textit{« Je souhaiterais », « Je pourrais », « Cela me permettrait »}.
                \item \textbf{Figures de style (Rhétorique) :}
                    \begin{itemize}
                        \item \cle{Énumération} : Lister compétences/qualités (\textit{« Rigueur, autonomie, sens du service »}). Effet : exhaustivité, densité.
                        \item \cle{Gradation} : Monter en puissance (\textit{« Apprendre, comprendre, maîtriser, innover »}). Effet : dynamique, progression.
                        \item \cle{Métonymie} : Remplacer le concret par l'abstrait ou l'institution (\textit{« Intégrer votre \textbf{maison} »} pour entreprise, \textit{« La \textbf{direction} »} pour les décideurs). Effet : appartenance, respect.
                        \item \cle{Hyperbole} (Mesurée) : \textit{« Une \textbf{forte} motivation »}, \textit{« Un \textbf{vif} intérêt »}. Effet : intensité. \textbf{Attention :} Pas de mensonge.
                    \end{itemize}
            \end{itemize}
        \item \textbf{Connecteurs logiques :} \textit{Par ailleurs, De plus, En outre, Ainsi, C'est pourquoi, Fort de...}
    \end{enumerate}
\end{methode}

\begin{exoresolu}[Rédaction du corps avec figures de style]
    \textbf{Énoncé :} Pour le même candidat Armand KAMGA (BTS Maintenance, stage SONEL Littoral), rédigez le \textbf{corps de la lettre} (2 paragraphes). Vous devez y intégrer \textbf{au moins une énumération, une gradation et une métonymie}. Mettez ces figures en \textbf{gras} dans votre copie.

    \tcblower
    \textbf{Solution détaillée :}

    \textbf{Paragraphe 1 (Compétences techniques) :}
    \begin{quote}
    \textit{Mon parcours technique, sanctionné par un BTS Maintenance des Systèmes Automatisés, m'a permis d'\textbf{acquérir, consolider, maîtriser} \textbf{(gradation)} les fondamentaux de l'électrotechnique et de l'automatisme industriel. Lors de mon stage de fin d'études à la Direction Régionale du Littoral, j'ai participé activement aux opérations de maintenance préventive et curative sur les postes HTA/BT, manipulé les relais de protection (Sepam, Micom) et assuré le suivi des tableaux de basse tension. Je maîtrise \textbf{la lecture de schémas électriques, la programmation d'automates (Siemens S7), l'utilisation d'oscilloscopes et de multimètres} \textbf{(énumération)}. Ces expériences m'ont forgé une rigueur indispensable à la sécurité des installations.}
    \end{quote}

    \textbf{Paragraphe 2 (Savoir-être + Motivation entreprise) :}
    \begin{quote}
    \textit{Au-delà des compétences techniques, mon engagement comme délégué de promotion m'a valu \textbf{écoute, diplomatie, sens des responsabilités} \textbf{(énumération/gradation)}. Intégrer \textbf{votre maison} \textbf{(métonymie : maison = entreprise SONEL)} serait pour moi l'aboutissement d'un parcours guidé par l'excellence du service public de l'énergie. Votre politique de modernisation du réseau, notamment le projet « Smart Grid », résonne particulièrement avec mon goût pour l'innovation technique.}
    \end{quote}

    \textbf{Vérification des figures :}
    \begin{itemize}
        \item \textbf{Gradation :} « acquérir, consolider, maîtriser » (progression dans la compétence).
        \item \textbf{Énumération :} Liste des outils/logiciels (densité technique) ; Liste qualités (densité humaine).
        \item \textbf{Métonymie :} « votre maison » pour « votre entreprise/votre structure ». Marque l'appartenance et le respect.
    \end{itemize}
    \textbf{Note :} L'impératif « Veuillez agréer » n'apparaîtra qu'en conclusion.
\end{exoresolu}

\begin{methode}[Intégrer la communication par l'image : Cohérence visuelle Lettre/CV]
    \textbf{Objectif :} Assurer une identité visuelle professionnelle (Personal Branding) entre la lettre et le CV.
    \begin{enumerate}
        \item \textbf{Police de caractères :} Une seule police (ou deux max : une pour titres, une pour corps). \cle{Police sans empattement (Sans Serif)} : Calibri, Arial, Helvetica, Roboto (lisibilité écran/impression). Taille 11 ou 12. Interligne 1,15 ou 1,5.
        \item \textbf{Mise en page (Marges, Alignement) :} Marges symétriques (2,5 cm). Justifié pour le corps. Alignement gauche pour en-têtes. \textbf{Cohérence :} Mêmes marges, même police, même en-tête (identité visuelle) sur Lettre et CV.
        \item \textbf{La Photo (sur le CV, pas sur la lettre) :} Professionnelle (fond neutre, tenue adaptée au poste, cadrage visage/épaules, sourire léger). \textbf{Interdit :} Photo de vacances, selfie, photo d'identité administrative stricte (sauf demande expresse).
        \item \textbf{Couleur (avec parcimonie) :} Un code couleur unique (ex: Bleu marine \#003366) pour : Nom (en-tête), Objet, Titres de parties du CV, Lignes de séparation. \textbf{Noir} pour le texte courant.
        \item \textbf{Format d'envoi :} \cle{PDF} (portrait, 1 page max pour la lettre). Nom du fichier : \textit{Nom\_Prenom\_LettreMotivation\_Poste.pdf}.
    \end{enumerate}
    \begin{attention}[Image et Lettre Officielle]
        La lettre officielle \textbf{ne contient pas d'image} (ni logo, ni photo, ni icône). L'image appartient au CV. La lettre joue sur la \cle{typographie} et la \cle{structure blanche} (aération) pour sa lisibilité.
    \end{attention}
\end{methode}

\begin{exoresolu}[Analyse de cohérence visuelle Lettre/CV]
    \textbf{Énoncé :} Un candidat envoie sa candidature par email.
    \textbf{Pièce 1 (Lettre) :} Police \textit{Times New Roman 12}, marges 3cm, en-tête bleu marine, objet en gras rouge, texte justifié, signature scannée en bleu. Fichier : \texttt{Lettre\_Armand.pdf}.
    \textbf{Pièce 2 (CV) :} Police \textit{Arial 10}, marges 1,5cm, en-tête vert fluo, photo en noir/blanc (fond blanc), titres en italique violet, puces en losanges roses. Fichier : \texttt{CV\_Final\_V2\_corrige.docx}.

    \textbf{Question :} Relevez 4 incohérences majeures nuisant à l'image professionnelle et proposez la correction pour chacune.

    \tcblower
    \textbf{Solution détaillée :}

    \begin{center}
    \begin{tabularx}{\linewidth}{|l|X|X|}\hline
    \textbf{Élément} & \textbf{Incohérence / Erreur} & \textbf{Correction recommandée} \\ \hline
    \textbf{Police} & Deux polices différentes (Times / Arial). Aspect brouillon. & \textbf{Unifier :} Arial 11 ou Calibri 11 pour les DEUX documents. \\ \hline
    \textbf{En-tête / Identité visuelle} & Couleurs différentes (Bleu marine / Vert fluo). Pas de logo commun. & \textbf{Harmoniser :} Même bandeau haut (Nom, Coordonnées, Couleur unique ex: Bleu \#003366) sur Lettre ET CV. \\ \hline
    \textbf{Format fichier} & CV en \texttt{.docx} (modifiable, mise en page instable selon OS). & \textbf{Exporter les DEUX en \texttt{.pdf}} (figé, professionnel, infalsifiable). \\ \hline
    \textbf{Nommage fichiers} & Noms vagues, versioning visible (\texttt{V2\_corrige}). & \textbf{Standard :} \texttt{Nom\_Prenom\_LettreMotivation\_IntitulePoste.pdf} / \texttt{Nom\_Prenom\_CV\_IntitulePoste.pdf}. \\ \hline
    \textbf{Mise en page CV} & Marges trop fines (1,5cm), couleurs multiples (violet, rose, vert), italique pour titres. & \textbf{Aérer :} Marges 2cm. \textbf{Sobriété :} 1 couleur d'accent max. Titres en \textbf{Gras}, pas italique. Puces standards (ronds/carrés). \\ \hline
    \end{tabularx}
    \end{center}
    \textbf{Conclusion :} La communication par l'image dans la candidature repose sur l'\textbf{unité graphique} (charte graphique personnelle) et la \textbf{stabilité technique} (PDF).
\end{exoresolu}

\begin{methode}[Production complète guidée : De l'analyse d'offre à la lettre finale]
    \textbf{Objectif :} Produire une lettre aboutie, sans erreur, en situation d'examen (Probatoire/Bac).
    \begin{enumerate}
        \item \textbf{Étape 1 : Analyse de l'offre (Brouillon - 10 min).} Surligner : \textbf{Mots-clés techniques} (logiciels, machines, diplômes), \textbf{Soft skills} (autonomie, équipe), \textbf{Verbes d'action} (gérer, concevoir, maintenir). Noter : Référence, Destinataire exact, Date limite.
        \item \textbf{Étape 2 : Matching Profil/Offre (Tableau - 10 min).}
            \begin{center}
            \begin{tabular}{|l|l|l|}\hline
            \textbf{Besoin Offre} & \textbf{Mon Preuve (Diplôme/Stage/Projet)} & \textbf{Mot-clé à réutiliser} \\ \hline
            Maintenance préventive & Stage SONEL : ronde quot. postes HTA & Maintenance préventive, HTA \\ \hline
            Automatismes Siemens & Projet fin d'études : Automate S7-1200 & Siemens, TIA Portal \\ \hline
            Esprit équipe & Délégué promo / Sport collectif & Coordination, communication \\ \hline
            \end{tabular}
            \end{center}
        \item \textbf{Étape 3 : Squelette rédactionnel (Brouillon - 5 min).} Rédiger les \textbf{mots-clés de transition} par paragraphe : Intro (Source/Poste/Atout majeur) -- P1 (Tech : Preuve 1, Preuve 2) -- P2 (Humain : Preuve 3 + Entreprise) -- Conclu (Dispo/Entretien/Politesse).
        \item \textbf{Étape 4 : Rédaction soignée (Propre - 20 min).} Appliquer structure externe, connecteurs, conditionnel, figures de style (1 ou 2 max), vocabulaire précis. \textbf{Une page max.}
        \item \textbf{Étape 5 : Rellecture systématique (Check-list - 5 min).} Voir la fiche «À retenir» ci-après.
    \end{enumerate}
\end{methode}

\begin{aretenir}[Check-list Rellecture « Zéro Faute » (Bac)]
    \begin{itemize}
        \item [$\checkmark$] \textbf{Orthographe/Grammaire :} Accords (participes passés, adjectifs), conjugaison (conditionnel vs futur), homophones (à/a, et/est, ou/où, son/sont, ces/ses/c'est).
        \item [$\checkmark$] \textbf{Structure externe :} Lieu/Date (mois en lettres), Objet (Gras, Réf), Appel (Monsieur,), Formule longue, Signature (Prénom NOM), PJ.
        \item [$\checkmark$] \textbf{Cohérence :} Nom destinataire = En-tête = Formule politesse. Poste/Ref = Objet = Intro.
        \item [$\checkmark$] \textbf{Ton/Style :} Pas d'impératif dans le corps. Pas de familier. Pas de répétitions lourdes (varier : « Mon stage / Cette expérience / Celle-ci »).
        \item [$\checkmark$] \textbf{Présentation :} Aération (paragraphes sautés), alignements, police unique, pas de ratures (brouillon séparé).
    \end{itemize}
\end{aretenir}

\begin{exoresolu}[Production intégrée : Lettre complète]
    \textbf{Énoncé :} \textit{Sujet type Probatoire.} Vous êtes \textbf{Chantal BILOA}, titulaire d'un \textbf{Baccalauréat G2 (Secrétariat Bureautique)} et d'un \textbf{BTS Assistant de Direction}. Vous avez effectué un stage de 3 mois à la \textbf{Chambre de Commerce de Douala} (gestion agenda, organisation réunions, rédaction PV). Vous répondez à l'annonce ci-dessous :

    \begin{quote}
    \textit{« ONG \textbf{PLAN INTERNATIONAL CAMEROUN} recrute un(e) Assistant(e) Administratif(ve) et Financier(ère). Réf : AAF-2025. Profil : Bac+2 Secrétariat/Assistanat. Maîtrise Pack Office (Word, Excel, PowerPoint), rédaction administrative, gestion agenda, organisation événements. Qualités : Rigueur, discrétion, sens de l'initiative. Envoyer CV + Lettre de motivation à : Directeur Pays, BP 1000 Douala, avant le 30/11/2025. »}
    \end{quote}
    \textbf{Rédigez la lettre de motivation complète.}

    \tcblower
    \textbf{Solution détaillée :}

    \begin{center}
    \begin{tabularx}{\linewidth}{|X|}\hline
    \textbf{CHANTAL BILOA} \\
    Quartier Bonapriso, BP 555 Douala \quad Tél : 6XX XX XX XX \quad Email : c.biloa@email.cm \\ \hline
    \textbf{Monsieur le Directeur Pays} \\
    \textbf{ONG PLAN INTERNATIONAL CAMEROUN} \\
    BP 1000 Douala \\ \hline
    \hfill Douala, le 18 novembre 2025 \\ \hline
    \textbf{Objet : Candidature au poste d'Assistant(e) Administratif(ve) et Financier(ère) (Réf : AAF-2025)} \\ \hline
    \textbf{Monsieur,} \\ \hline
    \textit{Suite à votre avis de recrutement publié sur votre site institutionnel, je vous adresse ma candidature pour le poste d'Assistant Administratif et Financier (Réf : AAF-2025). Titulaire d'un BTS Assistant de Direction et forte d'une expérience probante en gestion administrative au sein de la Chambre de Commerce de Douala, je souhaite mettre mon sens de l'organisation et ma rigueur au service de vos programmes humanitaires.} \\ \hline
    \textit{Sur le plan technique, mon stage de trois mois à la Chambre de Commerce m'a permis de \textbf{maîtriser la suite Office (Word, Excel, PowerPoint)} \textbf{(énumération)}, de gérer des agendas complexes de direction, d'organiser des sessions de formation et de rédiger des procès-verbaux de réunions statutaires. J'ai ainsi développé une \textbf{rapidité d'exécution, une discrétion absolue et une anticipation des besoins} \textbf{(gradation/énumération)} indispensables à ce poste.} \\ \hline
    \textit{Votre engagement pour la protection de l'enfance et l'égalité des genres, valeurs portées par \textbf{votre maison} \textbf{(métonymie)}, résonne profondément avec mes convictions personnelles. Intégrer vos équipes serait l'opportunité de conjuguer expertise administrative et utilité sociale.} \\ \hline
    \textit{Disponible immédiatement pour un entretien à votre convenance, je me tiens à votre disposition pour tout complément d'information. En vous remerciant de l'attention portée à ma candidature, je vous prie d'agréer, Monsieur le Directeur Pays, l'expression de mes salutations distinguées.} \\ \hline
    \hfill \textbf{Chantal BILOA} \\ \hline
    \textbf{Pièces jointes :} Curriculum Vitæ, Copie BTS Assistant de Direction, Copie Baccalauréat G2, Attestation de stage Chambre de Commerce, Extrait d'acte de naissance. \\ \hline
    \end{tabularx}
    \end{center}
    \textbf{Commentaire correcteur :} Structure externe parfaite. Introduction (Source/Poste/Accroche). Corps : P1 Technique (Mots-clés annonce repris + Énumération/Gradation). P2 Motivation/Valeurs (Métonymie). Conclusion (Action + Politesse adaptée). Vocabulaire précis. Une page. Zéro faute.
\end{exoresolu}

\begin{attention}[Les 5 erreurs qui coûtent des points au Bac (Probatoire)]
    \begin{enumerate}
        \item \textbf{Confondre « Lettre de motivation » et « Lettre manuscrite / Lettre à un ami » :} Oubli de l'en-tête destinataire, de l'objet, de la formule de politesse longue, utilisation du tutoiement, ton familier, absence de signature. $\rightarrow$ \textbf{Sanction :} Hors sujet / Note éliminatoire sur la forme.
        \item \textbf{Le « Copier-Coller » du CV en phrases :} Répéter la liste chronologique des diplômes/stages sans \textbf{lier} compétence $\rightarrow$ besoin du poste. Pas d'argumentation, pas de « valeur ajoutée ». $\rightarrow$ \textbf{Sanction :} Pénalité lourde sur « Contenu/Argumentation » (souvent 4 à 6 pts sur 20).
        \item \textbf{Erreurs de « Code » (Formules figées) :} « Je vous prie d'agréer, Monsieur, \textbf{mes salutations distinguées} » (oubli de « l'expression de » $\rightarrow$ formule tronquée). « Monsieur le Directeur » dans l'en-tête mais « Madame » dans la formule d'appel. Date en chiffres (15/11/2025) au lieu de « 15 novembre 2025 ». $\rightarrow$ \textbf{Sanction :} Perte de points « Forme/Conventions » (souvent 2 à 3 pts faciles à prendre).
        \item \textbf{Mauvaise gestion de l'Impératif / Conditionnel :} Écrire dans le corps : « \textbf{Regardez} mon CV », « \textbf{Appelez-moi} », « \textbf{Donnez-moi} une chance ». L'impératif est \textbf{interdit} hors formules de politesse finales. $\rightarrow$ \textbf{Sanction :} Faute de registre / Inconvenance (pénalité langue + communication).
        \item \textbf{Négliger la cohérence visuelle et le format :} Lettre sur 2 pages (interdit), police fantaisie, ratures/brouillon rendu comme copie finale, fichier .docx au lieu de .pdf (si dépôt numérique), nom de fichier « Lettre.doc ». $\rightarrow$ \textbf{Sanction :} Pénalité « Communication par l'image / Professionnalisme » (critère d'évaluation technique).
    \end{enumerate}
\end{attention}

\newpage

\coursec{Exercices}

\courssub{Je vérifie mes connaissances}

\begin{exercice}[QCM : Structure et vocabulaire de la lettre officielle]
Pour chaque question, une seule réponse est exacte. Recopiez le numéro et la lettre choisie.
\begin{enumerate}
    \item Dans la structure externe d'une lettre officielle, l'objet se place :
    \begin{itemize}
        \item[A.] Avant les coordonnées de l'expéditeur.
        \item[B.] Après la formule d'appel et avant le corps de la lettre.
        \item[C.] Après les coordonnées du destinataire et avant la formule d'appel.
        \item[D.] Après la signature.
    \end{itemize}
    \item La formule d'appel « Monsieur le Directeur Général, » correspond à :
    \begin{itemize}
        \item[A.] La prise de congé.
        \item[B.] La formule de politesse finale.
        \item[C.] L'incipit (formule d'appel).
        \item[D.] L'objet de la lettre.
    \end{itemize}
    \item Parmi les figures de style suivantes, laquelle repose sur un excès volontaire pour marquer l'intensité ?
    \begin{itemize}
        \item[A.] La métonymie.
        \item[B.] L'hyperbole.
        \item[C.] La gradation.
        \item[D.] L'énumération.
    \end{itemize}
    \item À l'impératif présent, la forme correcte pour le verbe \emph{recevoir} à la 2\textsuperscript{e} personne du pluriel (vous) est :
    \begin{itemize}
        \item[A.] Recevez.
        \item[B.] Recevoirez.
        \item[C.] Recevez.
        \item[D.] Recevrez.
    \end{itemize}
    \item Dans une lettre de motivation, le corps de la lettre a pour fonction principale de :
    \begin{itemize}
        \item[A.] Remercier le destinataire.
        \item[B.] Présenter le candidat, ses compétences et sa motivation (argumentation).
        \item[C.] Indiquer la date et le lieu.
        \item[D.] Répéter l'objet.
    \end{itemize}
\end{enumerate}
\end{exercice}

\begin{exercice}[Vrai ou Faux justifié]
Indiquez si les affirmations suivantes sont vraies ou fausses. Justifiez brièvement chaque réponse.
\begin{enumerate}
    \item La signature se place en haut à gauche de la lettre, au-dessus des coordonnées de l'expéditeur.
    \item La gradation ascendante (ex. : « apprendre, comprendre, maîtriser ») est pertinente dans une lettre de motivation pour montrer une progression professionnelle.
    \item L'impératif présent dans « Veuillez agréer, Monsieur, mes salutations distinguées » exprime un ordre strict et autoritaire.
    \item La métonymie « La direction a refusé ma candidature » (pour désigner les personnes qui dirigent) est un procédé courant dans le langage administratif.
    \item L'introduction d'une lettre de motivation doit obligatoirement mentionner la source de l'annonce (journal, site web, affichage) et le poste visé.
\end{enumerate}
\end{exercice}

\begin{exercice}[Définitions et mise en ordre]
\begin{enumerate}
    \item Définissez les termes suivants : \cle{structure externe}, \cle{structure interne}, \cle{formule de politesse}, \cle{objet de la lettre}.
    \item Remettez dans l'ordre logique les éléments suivants de la structure externe d'une lettre officielle :
    \begin{itemize}
        \item[1.] Formule d'appel (ex. : Monsieur le Directeur,)
        \item[2.] Coordonnées de l'expéditeur (en-tête)
        \item[3.] Signature
        \item[4.] Date et lieu
        \item[5.] Coordonnées du destinataire
        \item[6.] Objet
        \item[7.] Formule de politesse (prise de congé)
    \end{itemize}
    Justifiez la place de l'objet par rapport à la formule d'appel.
\end{enumerate}
\end{exercice}

\begin{exercice}[L'impératif : forme et valeur]
Conjuguez les verbes entre parenthèses à l'impératif présent (2\textsuperscript{e} personne du pluriel \emph{vous}) et précisez la \cle{valeur} (ordre, conseil, prière, invitation, souhait) de l'impératif dans chaque phrase extraite de lettres officielles.
\begin{enumerate}
    \item (\emph{veuillez}) \dots bien trouver ci-joint mon curriculum vitæ.
    \item (\emph{agréer}) \dots, Madame, l'expression de mes salutations distinguées.
    \item (\emph{recevoir}) \dots l'assurance de ma considération respectueuse.
    \item (\emph{croire}) \dots à ma détermination à servir votre entreprise.
    \item (\emph{permettre}) \dots que je vous rappelle ma disponibilité pour un entretien.
\end{enumerate}
\end{exercice}

\courssub{Je m'entraîne}

\begin{exercice}[Réécrire une introduction conforme aux usages]
Voici une introduction maladroite rédigée par un candidat :
\begin{quote}
« Salut ! Moi c'est Jean, je veux le job parce que j'ai besoin d'argent. J'ai vu l'annonce sur Facebook. »
\end{quote}
\begin{enumerate}
    \item Identifiez au moins quatre manquements aux conventions de la lettre officielle (registre, structure, politesse, précision).
    \item Réécrivez cette introduction en respectant les codes de la lettre de motivation (formule d'appel, objet, source de l'annonce, poste visé, accroche professionnelle). Le candidat s'appelle Jean Mbarga, il répond à une annonce pour un poste de \emph{Technicien de maintenance} parue dans \emph{Cameroon Tribune} du 10 mars 2026.
\end{enumerate}
\end{exercice}

\begin{exercice}[Figures de style dans le corps de la lettre]
Lisez l'extrait suivant du corps d'une lettre de motivation :
\begin{quote}
« Fort de cinq années d'expérience au sein de la CAMRAIL, j'ai acquis une \textbf{maîtrise totale} des protocoles de sécurité ferroviaire. Mon parcours, jalonné de \textbf{succès, de défis relevés et de responsabilités croissantes}, témoigne de ma rigueur. Je suis prêt à \textbf{porter la maison sur mon dos} pour garantir la satisfaction de vos clients. \textbf{La direction} de votre groupe, que je suis avec attention, incarne l'excellence. »
\end{quote}
\begin{enumerate}
    \item Identifiez et nommez les quatre figures de style présentes dans les passages en gras (dans l'ordre) : \textbf{maîtrise totale}, \textbf{succès, de défis relevés et de responsabilités croissantes}, \textbf{porter la maison sur mon dos}, \textbf{La direction}.
    \item Pour chaque figure, expliquez son effet de sens et dites si son emploi est \emph{pertinent} ou \emph{maladroit} dans le registre de la lettre officielle. Justifiez.
\end{enumerate}
\end{exercice}

\begin{exercice}[Communication par l'image : la photo d'identité]
Un candidat joint à son dossier une photographie d'identité aux caractéristiques suivantes : fond coloré (bleu clair), cadrage serré sur le visage, tenue décontractée (t-shirt à motifs), lunettes de soleil sur la tête, sourire large montrant les dents, éclairage latéral créant une ombre marquée sur la moitié du visage.
\begin{enumerate}
    \item Faites la \cle{dénotation} (description objective) de cette photographie.
    \item Analysez la \cle{connotation} (message implicite, impression sur le recruteur) de chacun de ces éléments : tenue, lunettes, sourire, éclairage, fond.
    \item Proposez \textbf{deux améliorations précises} pour rendre cette photo conforme aux attentes d'un dossier de candidature administratif ou technique.
\end{enumerate}
\end{exercice}

\begin{exercice}[Rédaction du corps de la lettre : argumentation]
Vous répondez à l'offre ci-dessous parue dans \emph{Cameroon Tribune} :
\begin{quote}
\textbf{RECRUTEMENT -- CAMRAIL} \\
Poste : \emph{Agent d'escale commercial} (Réf. AEC-2026/04) \\
Profil : BAC+2/3 en Commerce/Transport/Logistique. Maîtrise des outils bureautiques (Excel, Word). Anglais technique lu/parlé. Sens du service client, rigueur, disponibilité (horaires décalés). Expérience en gare ou aéroport appréciée. \\
Dossier : Lettre de motivation + CV + Copies diplômes. À déposer à la DRH Yaoundé avant le 30/04/2026.
\end{quote}
Le candidat : \textbf{Assana Bello}, titulaire d'un \textbf{BTS Transport et Logistique} (2023), stage de 6 mois à la \textbf{GA Camair-Co} (service escale bagages), pratique l'anglais (niveau B2), disponible immédiatement.
\begin{enumerate}
    \item Rédigez le \textbf{corps de la lettre} (2 à 3 paragraphes) en structurant l'argumentation : \textbf{Paragraphe 1} : Formation et adéquation poste/profil. \textbf{Paragraphe 2} : Expérience professionnelle (stage) et compétences techniques/comportementales. \textbf{Paragraphe 3} : Motivation spécifique pour la CAMRAIL et qualités personnelles (disponibilité, rigueur).
    \item Intégrez naturellement \textbf{une énumération} et \textbf{une gradation} dans votre rédaction.
\end{enumerate}
\end{exercice}

\courssub{Je me prépare au Bac}

\begin{exercice}[Sujet type Probatoire -- Lettre de motivation CAMRAIL]
\textbf{Durée : 45 minutes -- Barème indicatif sur 20 points}
\begin{quote}
\textbf{ANNONCE :} La Cameroon Railways (CAMRAIL) lance un recrutement pour le poste de \textbf{Technicien de Maintenance Voie} (Réf. TMV-2026/11).
\textbf{Profil recherché :} BTS/DUT Génie Civil ou Maintenance Industrielle. Connaissance des normes de sécurité ferroviaire. Capacité à travailler en équipe et en horaires décalés (nuit, week-end). Permis B exigé. Première expérience (stage ou emploi) souhaitée.
\textbf{Composition du dossier :} Lettre de motivation manuscrite + CV + Photocopies certifiées des diplômes.
\textbf{Dépôt :} Direction des Ressources Humaines, CAMRAIL, B.P. 1234 Douala, au plus tard le 15 mai 2026.
\end{quote}
\textbf{Consigne :} Vous êtes \textbf{Paul Ekedi}, titulaire d'un \textbf{BTS Génie Civil} (option Travaux Publics) obtenu en 2024 à l'ISTP de Douala. Vous avez effectué un stage de fin d'études de 3 mois chez \textbf{RAZEL} (chantier de réhabilitation RN10) où vous avez participé au suivi des ouvrages d'art. Vous êtes titulaire du Permis B. Vous maîtrisez les logiciels AutoCAD et MS Project.
Rédigez la \textbf{lettre de motivation complète} (structure externe + structure interne : introduction, corps, conclusion) en 25 à 30 lignes.
\begin{center}
\textbf{Barème de correction (indicatif) :} \\
Structure externe (en-têtes, date, objet, formules d'appel/congé, signature) : 4 pts \\
Structure interne (Introduction : source, poste, accroche / Corps : argumentation cohérente, lien profil-poste / Conclusion : disponibilité, demande entretien, formule de politesse) : 6 pts \\
Qualité de la langue (syntaxe, vocabulaire technique, registres, figures de style pertinentes, impératifs corrects) : 4 pts \\
Argumentation et pertinence du contenu (adéquation profil/poste, motivation réelle) : 4 pts \\
Présentation (propreté, paragraphage, lisibilité, respect du nombre de lignes) : 2 pts
\end{center}
\end{exercice}

\begin{exercice}[Sujet type Bac Technique -- Lettre de motivation SONARA]
\textbf{Durée : 1 heure -- Production écrite}
\begin{quote}
\textbf{SONARA -- SOCIÉTÉ NATIONALE DE RAFFINAGE} \\
\textbf{Avis de recrutement n\textsuperscript{o} 2026/003} \\
Poste : \textbf{Technicien Procédés (H/F)} -- Site de Limbé \\
\textbf{Missions :} Surveillance des unités de distillation, relevés de paramètres (température, pression, débit), participation aux arrêts programmés, respect strict des consignes HSE (Hygiène, Sécurité, Environnement).
\textbf{Profil :} BTS/DUT Génie des Procédés, Chimie Industrielle ou Maintenance des Systèmes Automatisés. Notions de raffinage/pétrochimie. Lecture de P\&ID. Rigueur, réactivité, esprit d'équipe. Anglais technique : atout majeur.
\textbf{Dossier :} Lettre de motivation + CV détaillé + Copies diplômes + Extrait casier judiciaire n\textsuperscript{o}3.
\textbf{Adresse :} DRH SONARA, B.P. 456 Limbé. Date limite : 20 juin 2026.
\end{quote}
\textbf{Consigne :} Vous êtes \textbf{Marie Ngono}, titulaire d'un \textbf{DUT Génie des Procédés} (ENSPY Yaoundé, 2023). Vous avez validé un stage de 4 mois à la \textbf{SONARA} même (unité Distillation Atmosphérique) où vous avez assisté l'opérateur de salle de commande. Vous lisez couramment les schémas P\&ID. Vous pratiquez l'anglais technique (niveau TOEIC 750). Vous êtes sensibilisée aux normes HSE (certificat « Sécurité Industrielle Niveau 1 » obtenu lors du stage).
Rédigez la \textbf{lettre de motivation} (25 à 30 lignes) en soignant particulièrement :
\begin{enumerate}
    \item L'\textbf{introduction} (rappel précis de l'annonce, poste, accroche valorisant le stage \emph{in situ}).
    \item Le \textbf{corps} (démonstration par les compétences techniques acquises : P\&ID, HSE, anglais, travail en équipe sur site Seveso).
    \item La \textbf{conclusion} (disponibilité immédiate, demande d'entretien, formule de politesse complète et correcte).
\end{enumerate}
\end{exercice}

\begin{exercice}[Sujet complet : Lettre à trous et analyse langagière]
\textbf{Partie A -- Complétion d'une lettre de motivation (10 pts)}
Complétez la lettre ci-dessous en rédigeant les parties manquantes (entre crochets) pour répondre à l'offre d'\textbf{ENEO} pour un poste de \textbf{Technicien Réseaux Aériens} (BTS Électrotechnique, travail en hauteur, permis B, sens du service public). Le candidat est \textbf{Kevin Abanda}, BTS Électrotechnique 2024, stage ENEO Douala-Bassa (maintenance postes HTA/BT), habilitation H0B0, permis B, sportif (endurance).

\begin{quote}
[Coordonnées expéditeur] \\[0pt]
[Coordonnées destinataire : Directeur Régional ENEO Littoral] \\[0pt]
[Date et Lieu] \\[0pt]
\textbf{Objet :} [Objet complet avec référence] \\[0pt]
[Formule d'appel] \\[0pt]
\textbf{Introduction :} [À rédiger : source annonce, poste visé, accroche stage ENEO] \\[0pt]
\textbf{Corps -- §1 :} [À rédiger : Formation BTS + modules pertinents (machines tournantes, protection réseaux) + Habilitation] \\[0pt]
\textbf{Corps -- §2 :} [À rédiger : Expérience stage (diagnostic pannes, remplacement isolateurs, travail en nacelle, respect consignes sécurité) + Permis B + Condition physique] \\[0pt]
\textbf{Corps -- §3 :} [À rédiger : Motivation service public, continuité électricité, valeurs ENEO, anglais technique lu] \\[0pt]
\textbf{Conclusion :} [À rédiger : Disponibilité, demande entretien, formule de politesse complète] \\[0pt]
[Signature]
\end{quote}

\textbf{Partie B -- Analyse langagière (10 pts)}
Reprenons la phrase suivante, extraite du paragraphe 2 rédigé par Kevin :
\begin{quote}
« Lors de mon stage, j'ai \textbf{diagnostiqué, isolé, réparé} et \textbf{validé} le rétablissement du courant sur une ligne 30 kV, \textbf{gravissant} les pylônes avec assurance malgré le vent. »
\end{quote}
\begin{enumerate}
    \item Identifiez la figure de style réalisée par la série « diagnostiqué, isolé, réparé, validé ». Justifiez.
    \item Quelle est la nature et la fonction grammaticale de « gravissant » ? Quelle valeur aspectuelle apporte-t-il à l'action ?
    \item Transformez la phrase à la voix passive (sujet : \emph{le rétablissement du courant}). Observez le changement de perspective.
    \item Dans la formule de politesse finale « Je vous prie d'agréer, Monsieur le Directeur, l'assurance de ma considération distinguée », conjuguez le verbe \emph{agréer} à la forme requise et nommez le mode et le temps. Quelle est la valeur de ce mode ici ?
\end{enumerate}
\end{exercice}

\newpage

\coursec{Corrigés des exercices}

\coursec{Corrigés des exercices — Leçon 21 : Produire une lettre officielle (lettre de motivation)}

\begin{corrige}[Exercice 1 — QCM : Structure et vocabulaire de la lettre officielle]
\begin{enumerate}
    \item \textbf{C.} L'objet se place après les coordonnées du destinataire et avant la formule d'appel (incipit). Il permet au destinataire d'identifier immédiatement le sujet du courrier.
    \item \textbf{C.} « Monsieur le Directeur Général, » est la \cle{formule d'appel} (ou incipit), qui ouvre la lettre en désignant le destinataire avec le titre requis.
    \item \textbf{B.} L'\cle{hyperbole} est une figure d'\textbf{amplification} qui repose sur un excès volontaire (exagération) pour marquer l'intensité d'un sentiment ou d'une qualité.
    \item \textbf{A.} (et C, identiques). À l'impératif présent, 2\textsuperscript{e} personne du pluriel, le verbe \emph{recevoir} (groupe 3, \textbf{verbes en -cevoir}) se conjugue : \emph{nous recevons $\rightarrow$ recevez}. La forme \emph{recevrez} est au futur simple.
    \item \textbf{B.} Le corps de la lettre de motivation est la partie argumentative où le candidat expose son parcours, ses compétences (\cle{savoir-faire}, \cle{savoir-être}) et sa motivation pour le poste visé.
\end{enumerate}
\end{corrige}

\begin{corrige}[Exercice 2 — Vrai ou Faux justifié]
\begin{enumerate}
    \item \textbf{Faux.} La signature se place en \textbf{bas à droite} de la lettre, après la formule de politesse (prise de congé). Les coordonnées de l'expéditeur sont en \textbf{haut à gauche} (en-tête).
    \item \textbf{Vrai.} La \cle{gradation ascendante} (ex. : « apprendre, comprendre, maîtriser ») organise les termes par intensité croissante. Elle est très pertinente pour illustrer une progression professionnelle, une montée en compétences ou en responsabilités.
    \item \textbf{Faux.} Dans « Veuillez agréer... », l'impératif présent (\emph{veuillez}, \emph{agréez}) a une valeur de \cle{prière} ou de \cle{souhait} conventionnel, figé dans le code de la politesse administrative. Il n'exprime pas un ordre autoritaire.
    \item \textbf{Vrai.} « La direction » (l'institution, le lieu) est utilisé pour désigner « les personnes qui dirigent » (les décideurs). C'est une \cle{métonymie} (contenant pour contenu / institution pour personnes), procédure standard dans le langage administratif et professionnel.
    \item \textbf{Vrai.} L'introduction doit être \emph{contextualisée} : mentionner la source (journal, site, affichage), la date de l'annonce, la référence de l'offre et le poste exact visé. Cela prouve la traçabilité de la candidature et l'attention du candidat.
\end{enumerate}
\end{corrige}

\begin{corrige}[Exercice 3 — Définitions et mise en ordre]
\begin{enumerate}
    \item \textbf{Définitions :}
    \begin{itemize}
        \item \cle{Structure externe} : Ensemble des éléments formels et matériels qui encadrent le texte de la lettre (en-têtes expéditeur/destinataire, date/lieu, objet, formules d'appel et de politesse, signature). Elle assure l'identification des parties et la validité administrative.
        \item \cle{Structure interne} : Organisation logique du contenu rédactionnel en trois parties : \emph{introduction} (accroche, origine de la candidature), \emph{corps} (argumentation, adéquation profil/poste), \emph{conclusion} (disponibilité, demande d'entretien, formule de politesse).
        \item \cle{Formule de politesse} : Expression codifiée, obligatoire, marquant le respect dû au destinataire. Elle existe en ouverture (formule d'appel : « Monsieur le Directeur, ») et en clôture (prise de congé : « Je vous prie d'agréer... »).
        \item \cle{Objet de la lettre} : Ligne concise, placée après les coordonnées du destinataire, résumant le but du courrier (ex. : « Candidature au poste de Technicien de maintenance - Réf. TM-2026 »). Elle facilite le tri et le traitement du courrier.
    \end{itemize}
    \item \textbf{Ordre logique de la structure externe :}
    \begin{enumerate}
        \item Coordonnées de l'expéditeur (en-tête)
        \item Date et lieu
        \item Coordonnées du destinataire
        \item \textbf{Objet}
        \item Formule d'appel (ex. : Monsieur le Directeur,)
        \item [Corps de la lettre : Introduction, Développement, Conclusion]
        \item Formule de politesse (prise de congé)
        \item Signature
    \end{enumerate}
    \textbf{Justification de la place de l'objet :} L'objet précède la formule d'appel car il constitue la \emph{référence administrative} du courrier. Il informe le destinataire (ou le service de courrier) du sujet \emph{avant} même la salutation, permettant un acheminement rapide vers le service concerné. C'est une norme de la correspondance professionnelle (norme NF Z 10-001).
\end{enumerate}
\end{corrige}

\begin{corrige}[Exercice 4 — L'impératif : forme et valeur]
\begin{center}
\begin{tabularx}{\linewidth}{|c|X|c|X|}\hline
\textbf{N°} & \textbf{Verbe (Infinitif)} & \textbf{Forme impératif présent (2\textsuperscript{e} pl.)} & \textbf{Valeur} \\ \hline
1 & \emph{vouloir} & \textbf{Veuillez} & \cle{Prière} / Demande polie (formule figée « Veuillez trouver... ») \\ \hline
2 & \emph{agréer} & \textbf{Agréez} & \cle{Souhait} / Formule de politesse finale figée (« Agréez, Madame... ») \\ \hline
3 & \emph{recevoir} & \textbf{Recevez} & \cle{Souhait} / Formule de politesse finale figée (« Recevez l'assurance... ») \\ \hline
4 & \emph{croire} & \textbf{Croyez} & \cle{Invitation} / \cle{Souhait} (inviter le destinataire à croire en la sincérité) \\ \hline
5 & \emph{permettre} & \textbf{Permettez} & \cle{Demande d'autorisation} / Politesse (demander la permission de rappeler un point) \\ \hline
\end{tabularx}
\end{center}
\textbf{Rappel morphologique :} Pour les verbes en \emph{-er} (agréer, permettre) et la plupart des verbes en \emph{-oir} (recevoir), l'impératif 2\textsuperscript{e} pl. se forme sur le radical de l'indicatif présent 2\textsuperscript{e} pl. (\emph{vous agréez, vous permettez, vous recevez}) en conservant la terminaison \emph{-ez}. \emph{Croire} est irrégulier : \emph{vous croyez $\rightarrow$ croyez}. \textbf{Attention :} \emph{Vouloir} est irrégulier : \emph{vous voulez $\rightarrow$ \textbf{veuillez}} (changement de radical).
\end{corrige}

\begin{corrige}[Exercice 5 — Réécrire une introduction conforme aux usages]
\begin{enumerate}
    \item \textbf{Quatre manquements majeurs :}
    \begin{enumerate}
        \item \cle{Registre de langue} : Familier/Inapproprié (« Salut ! », « Moi c'est Jean », « le job », « j'ai besoin d'argent »).
        \item \cle{Structure externe absente} : Pas de coordonnées expéditeur/destinataire, pas de date/lieu, pas d'objet, pas de formule d'appel ni de prise de congé.
        \item \cle{Politesse / Convention} : Tutoiement implicite, absence de formules de respect (« Monsieur, », « Je vous prie d'agréer... »).
        \item \cle{Précision / Argumentation} : Source vague (« Facebook » sans date ni référence), motivation purement pécuniaire non valorisante, pas d'accroche professionnelle.
    \end{enumerate}
    \item \textbf{Introduction réécrite (modèle) :}
    \begin{quote}
    \textbf{Jean Mbarga} \\
    12, Rue de la Paix — B.P. 456 Yaoundé \\
    Tel : 6XX XX XX XX — Email : jean.mbarga@email.cm \\[0.5em]
    \textbf{À l'attention de Monsieur le Directeur des Ressources Humaines} \\
    \textbf{Entreprise [Nom de l'entreprise]} \\
    B.P. 789 Douala \\[0.5em]
    Fait à Yaoundé, le 15 mars 2026 \\[0.5em]
    \textbf{Objet : Candidature au poste de Technicien de maintenance (Réf. Annonce Cameroon Tribune du 10/03/2026)} \\[0.5em]
    \textbf{Monsieur le Directeur,} \\[0.5em]
    Suite à votre annonce parue dans le \emph{Cameroon Tribune} du 10 mars 2026 relative au recrutement d'un Technicien de maintenance, je me permets de vous adresser ma candidature. Fort de mon diplôme de BTS Maintenance Industrielle et de trois années d'expérience en milieu industriel, je souhaite mettre mes compétences techniques et ma rigueur au service de votre structure.
    \end{quote}
\end{enumerate}
\end{corrige}

\begin{corrige}[Exercice 6 — Figures de style dans le corps de la lettre]
\begin{enumerate}
    \item \textbf{Identification des figures (dans l'ordre) :}
    \begin{enumerate}
        \item « \textbf{maîtrise totale} » : \cle{Hyperbole} (excès : « totale » est une exagération, une maîtrise est rarement absolue).
        \item « \textbf{succès, de défis relevés et de responsabilités croissantes} » : \cle{Énumération} (suite de trois noms) \textbf{et} \cle{Gradation ascendante} (progression logique : succès $\rightarrow$ défis $\rightarrow$ responsabilités croissantes).
        \item « \textbf{porter la maison sur mon dos} » : \cle{Hyperbole} (image forte, exagération pour dire « assumer une charge immense ») — \emph{Note : c'est aussi une image/métaphore figée, mais l'hyperbole est la figure rhétorique principale au programme}.
        \item « \textbf{La direction} » : \cle{Métonymie} (l'institution/le service « direction » désigne les personnes qui la composent : les décideurs, le DRH, le DG).
    \end{enumerate}
    \item \textbf{Analyse de la pertinence :}
    \begin{itemize}
        \item \textbf{Maîtrise totale (Hyperbole) :} \emph{Maladroit / Risqué.} Dans une lettre officielle, l'hyperbole sur une compétence technique (« totale ») peut passer pour de l'arrogance ou du mensonge. Mieux vaut : « une solide maîtrise », « une expertise confirmée ».
        \item \textbf{Succès, défis, responsabilités croissantes (Énumération + Gradation) :} \emph{Pertinent.} L'énumération structure le parcours ; la gradation ascendante démontre dynamisme et évolution professionnelle. C'est un argument fort.
        \item \textbf{Porter la maison sur mon dos (Hyperbole/Image) :} \emph{Maladroit / Inadapté.} Registre trop familier, imagé, peu professionnel. Une lettre de motivation exige le registre standard/soutenu. Remplacer par : « m'investir pleinement », « assumer l'entière responsabilité de... », « relever ce défi avec détermination ».
        \item \textbf{La direction (Métonymie) :} \emph{Pertinent.} Usage consacré, précis, respectueux du code administratif. Elle désigne l'autorité hiérarchique sans nommer les individus.
    \end{itemize}
\end{enumerate}
\end{corrige}

\begin{corrige}[Exercice 7 — Communication par l'image : la photo d'identité]
\begin{enumerate}
    \item \textbf{Dénotation (description objective) :} La photographie présente un visage de face, cadré serré (épaules coupées), sur un fond uni bleu clair. Le sujet porte un t-shirt à motifs colorés, des lunettes de soleil posées sur le front/cheveux. Il affiche un large sourire découvrant les dents. L'éclairage provient de la gauche (latéral), projetant une ombre marquée sur la moitié droite du visage.
    \item \textbf{Connotation (analyse de l'impression sur le recruteur) :}
    \begin{center}
    \begin{tabularx}{\linewidth}{|l|X|}\hline
    \textbf{Élément} & \textbf{Connotation (Message implicite)} \\ \hline
    Tenue (t-shirt motifs) & Décontraction excessive, manque de sérieux, inadaptation aux codes de l'entreprise, amateurisme. \\ \hline
    Lunettes de soleil & Dissimulation du regard (confiance, sincérité), attitude « vacances/loisirs », non-respect des normes (visage découvert obligatoire). \\ \hline
    Sourire large (dents) & Familiarité forcée, manque de gravité professionnelle, peut sembler peu naturel ou insincère. \\ \hline
    Éclairage latéral (ombre) & Amateurisme technique, photo non réalisée par un professionnel, qualité médiocre, cache une partie du visage. \\ \hline
    Fond bleu clair & Non-conformité aux standards administratifs (fond blanc/gris clair/neutre exigé), distraction visuelle. \\ \hline
    \end{tabularx}
    \end{center}
    \item \textbf{Deux améliorations précises :}
    \begin{enumerate}
        \item \textbf{Tenue et accessoires :} Porter une tenue professionnelle sobre (chemise/blouse unie, veste si possible), retirer les lunettes (ou porter lunettes de vue sans reflets si nécessaire), visage entièrement découvert, chevelure soignée.
        \item \textbf{Technique :} Fond uni blanc ou gris très clair (norme ISO/IEC 19794-5), éclairage frontal diffus (deux sources ou flash annulaire) pour supprimer les ombres, cadrage standard (visage 70-80\% de la hauteur, épaules visibles), expression neutre ou sourire léger bouche fermée, regard droit dans l'objectif.
    \end{enumerate}
\end{enumerate}
\end{corrige}

\begin{corrige}[Exercice 8 — Rédaction du corps de la lettre : argumentation CAMRAIL]
\textbf{Corps de la lettre (Modèle pour Assana Bello) :}
\begin{quote}
\textbf{Paragraphe 1 — Formation et adéquation :} \\
Titulaire d'un BTS Transport et Logistique obtenu en 2023, j'ai acquis une solide maîtrise des fondamentaux du transport ferroviaire et de la gestion des flux de voyageurs. Les modules « Exploitation ferroviaire », « Réglementation transports » et « Gestion commerciale » correspondent directement aux exigences de votre poste d'Agent d'escale commercial. Ma formation m'a permis de développer une rigueur analytique et une parfaite connaissance des procédures de sûreté, atouts indispensables pour garantir la régularité et la qualité du service à la clientèle.

\textbf{Paragraphe 2 — Expérience et compétences (avec énumération et gradation) :} \\
Mon stage de six mois au service escale bagages de \textbf{Camair-Co} a été le creuset de mon professionnalisme. J'ai participé activement à l'enregistrement des passagers, au suivi des bagages en correspondance et à la gestion des réclamations, dans le respect strict des normes IATA. Cette immersion m'a permis de développer des compétences techniques et comportementales complémentaires : \textbf{maîtrise des outils bureautiques (Excel, Word, Amadeus), pratique courante de l'anglais technique (niveau B2), sens aigu du service client et réactivité face aux imprévus}. J'ai ainsi appris à \textbf{écouter, comprendre, anticiper et résoudre} les situations complexes, une gradation qui illustre ma montée en autonomie.

\textbf{Paragraphe 3 — Motivation CAMRAIL et qualités :} \\
Intégrer la CAMRAIL représente pour moi une évidence : acteur majeur du corridor Douala-Ngaoundéré, votre entreprise incarne l'excellence technique et le service public que je souhaite servir. Votre engagement pour la modernisation du réseau et la satisfaction client résonne avec mes valeurs de \textbf{rigueur, de disponibilité (horaires décalés, week-ends, nuits) et d'esprit d'équipe}. Titulaire du permis B et immédiatement disponible, je suis prêt à m'investir pleinement dans la réussite de vos opérations au sol.
\end{quote}
\textbf{Figures intégrées :} \cle{Énumération} (compétences : « maîtrise..., pratique..., sens..., réactivité... ») et \cle{Gradation ascendante} (verbes : « écouter, comprendre, anticiper, résoudre »).
\end{corrige}

\begin{corrige}[Exercice 9 — Sujet type Probatoire : Lettre de motivation CAMRAIL (Technicien Maintenance Voie)]
\textbf{Proposition de lettre complète (environ 28 lignes) :}
\begin{quote}
\textbf{Paul Ekedi} \\
Quartier Bonabéri, Rue 12 — B.P. 555 Douala \\
Tel : 6XX XX XX XX — Email : paul.ekedi@email.cm \\[0.5em]
\textbf{À l'attention de Monsieur le Directeur des Ressources Humaines} \\
\textbf{CAMRAIL — Direction Régionale Littoral} \\
B.P. 1234 Douala \\[0.5em]
Fait à Douala, le 25 avril 2026 \\[0.5em]
\textbf{Objet : Candidature au poste de Technicien de Maintenance Voie (Réf. TMV-2026/11)} \\[0.5em]
\textbf{Monsieur le Directeur,} \\[0.5em]
Suite à votre avis de recrutement publié dans le \emph{Cameroon Tribune} en date du 15 avril 2026, je vous adresse ma candidature pour le poste de Technicien de Maintenance Voie. Titulaire d'un BTS Génie Civil option Travaux Publics (ISTP Douala, promotion 2024), je suis vivement intéressé par les défis techniques de la maintenance d'infrastructure ferroviaire. \\[0.5em]
Mon stage de fin d'études de trois mois chez \textbf{RAZEL} (chantier réhabilitation RN10, section ouvrages d'art) m'a confronté aux réalités du terrain : suivi géométrique des voies d'accès, contrôle des armatures et du béton, lecture de plans d'exécution (AutoCAD) et planification des tâches (MS Project). J'ai ainsi développé une rigueur de contrôle et une culture de la sécurité (port EPI, respect consignes) directement transposables à la maintenance voie (géométrie, rail, ballast, ouvrages d'art). \\[0.5em]
Titulaire du Permis B, habitué aux horaires décalés et au travail en équipe sur site linéaire, je suis prêt à m'engager durablement au sein de vos équipes de maintenance. Ma motivation pour le rail, ma capacité d'adaptation et ma maîtrise des outils numériques (CAO, GPAO) seront, je l'espère, des atouts pour la CAMRAIL. \\[0.5em]
Disponible immédiatement pour un entretien à votre convenance, je vous prie d'agréer, Monsieur le Directeur, l'assurance de ma considération distinguée. \\[1em]
\textbf{Paul Ekedi} \\
\textit{Signature}
\end{quote}

\textbf{Barème indicatif détaillé (sur 20) :}
\begin{itemize}
    \item \textbf{Structure externe (4 pts) :} En-têtes complets (1), Date/Lieu (0,5), Objet précis + Réf (1), Formule d'appel (0,5), Formule de politesse complète + Signature (1).
    \item \textbf{Structure interne (6 pts) :} Introduction : source + poste + accroche stage (2). Corps : Argumentation cohérente (formation + stage RAZEL + compétences techniques AutoCAD/MS Project + Permis B + dispo horaires) lien profil/poste explicite (3). Conclusion : Dispo + Demande entretien + Politesse (1).
    \item \textbf{Qualité langue (4 pts) :} Syntaxe correcte, vocabulaire technique (géométrique, ballast, ouvrages d'art, CAO, GPAO), registre soutenu, impératifs corrects (agréer, recevoir), figures pertinentes (ex: énumération compétences).
    \item \textbf{Argumentation/Pertinence (4 pts) :} Adéquation BTS TP / Maintenance Voie (transférabilité), Stage RAZEL pertinent (ouvrages d'art = ponts/dalais), Permis B cité, Dispo immédiate.
    \item \textbf{Présentation (2 pts) :} Propreté, paragraphage visible (sauts de ligne), lisibilité, respect 25-30 lignes.
\end{itemize}
\textbf{Points de vigilance (Probatoire) :} L'objet \textbf{doit} reprendre la référence exacte de l'annonce (TMV-2026/11). La formule de politesse finale \textbf{doit} contenir « agréer » à l'impératif (\emph{agréez}) et « considération distinguée » (pas « salutations » seul). Le stage RAZEL doit être explicitement lié aux compétences attendues (voie, ouvrages d'art, sécurité).
\end{corrige}

\begin{corrige}[Exercice 10 — Sujet type Bac Technique : Lettre de motivation SONARA (Technicien Procédés)]
\textbf{Proposition de lettre complète (environ 28 lignes) :}
\begin{quote}
\textbf{Marie Ngono} \\
Quartier Mvog-Ada, Yaoundé — B.P. 888 Yaoundé \\
Tel : 6XX XX XX XX — Email : marie.ngono@email.cm \\[0.5em]
\textbf{À l'attention de Monsieur le Directeur des Ressources Humaines} \\
\textbf{SONARA — Société Nationale de Raffinage} \\
B.P. 456 Limbé \\[0.5em]
Fait à Yaoundé, le 10 juin 2026 \\[0.5em]
\textbf{Objet : Candidature au poste de Technicien Procédés (Avis n° 2026/003 - Site Limbé)} \\[0.5em]
\textbf{Monsieur le Directeur,} \\[0.5em]
C'est avec un vif intérêt que j'ai pris connaissance de votre avis de recrutement n° 2026/003 pour le poste de Technicien Procédés sur le site de Limbé. Titulaire d'un DUT Génie des Procédés (ENSPY Yaoundé, 2023), je souhaite vivement réintégrer la SONARA où j'ai effectué un stage déterminant de 4 mois au sein de l'unité de Distillation Atmosphérique. \\[0.5em]
Cette expérience \emph{in situ} m'a permis d'assister l'opérateur de salle de commande dans la surveillance continue des paramètres critiques (température, pression, débit, niveau) et la lecture courante des schémas \textbf{P\&ID} (Piping and Instrumentation Diagram). J'ai participé aux rondes terrain, aux relevés d'échantillonnage et à la rédaction des fiches de quart, dans le respect strict des consignes \textbf{HSE} (Hygiène, Sécurité, Environnement) inhérentes à un site Seveso seuil haut. Mon certificat « Sécurité Industrielle Niveau 1 », obtenu lors de ce stage, atteste de ma culture préventionniste. \\[0.5em]
Au-delà de mes acquis techniques (thermodynamique, opérations unitaires, instrumentation), je dispose d'atouts majeurs pour ce poste : une pratique confirmée de l'\textbf{anglais technique} (TOEIC 750) pour exploiter la documentation constructeur, une rigueur éprouvée en environnement contraignant, et un esprit d'équipe forgé par le travail en 3x8. Intégrer vos équipes serait la suite logique de mon parcours. \\[0.5em]
Disponible immédiatement, je serais honorée de vous exposer ma motivation et mes compétences lors d'un entretien. Je vous prie d'agréer, Monsieur le Directeur, l'expression de mes salutations distinguées. \\[1em]
\textbf{Marie Ngono} \\
\textit{Signature}
\end{quote}

\textbf{Points de vigilance (Bac Technique) :}
\begin{itemize}
    \item \textbf{Introduction :} Citation précise de l'avis (n° 2026/003), site (Limbé), poste exact. Accroche forte : stage \emph{in situ} (unité Distillation Atmosphérique) = connaissance du site, des process, de la culture HSE SONARA.
    \item \textbf{Corps :} Preuves techniques concrètes : P\&ID (lecture), paramètres (T, P, D, N), HSE (Seveso, certif N1), Anglais (TOEIC 750), Travail en équipe (3x8). Vocabulaire métier exact.
    \item \textbf{Conclusion :} Dispo immédiate, demande entretien explicite, formule de politesse \emph{complète et correcte} (« agréer... expression de mes salutations distinguées »).
    \item \textbf{Langue :} Impératifs corrects (\emph{agréez}), registre soutenu, termes techniques (Seveso, P\&ID, quart, rondes, échantillonnage).
\end{itemize}
\end{corrige}

\begin{corrige}[Exercice 11 — Sujet complet : Lettre à trous ENEO + Analyse langagière]
\textbf{Partie A — Complétion de la lettre (Kevin Abanda — Technicien Réseaux Aériens)}
\begin{quote}
\textbf{Kevin Abanda} \\
Quartier Makepe, Douala — B.P. 1010 Douala \\
Tel : 6XX XX XX XX — Email : kevin.abanda@email.cm \\[0.5em]
\textbf{À l'attention de Monsieur le Directeur Régional} \\
\textbf{ENEO Cameroun — Direction Régionale Littoral} \\
B.P. 2020 Douala \\[0.5em]
Fait à Douala, le 20 mai 2026 \\[0.5em]
\textbf{Objet : Candidature au poste de Technicien Réseaux Aériens (Réf. TRA-2026/ENEO-LIT)} \\[0.5em]
\textbf{Monsieur le Directeur,} \\[0.5em]
\textbf{Introduction :} Suite à votre offre d'emploi parue sur le site corporate d'ENEO en date du 15 mai 2026, je me permets de vous proposer ma candidature pour le poste de Technicien Réseaux Aériens. Mon stage de fin d'études au sein de votre direction régionale (agence Douala-Bassa) m'a permis de découvrir l'exigence de votre métier et de confirmer mon attrait pour la maintenance des réseaux de distribution. \\[0.5em]
\textbf{Corps — §1 :} Titulaire d'un BTS Électrotechnique (2024), j'ai approfondi les modules de \textbf{machines tournantes}, de \textbf{protection des réseaux} et d'automatismes industriels. Cette formation théorique solide est complétée par l'obtention de mon \textbf{habilitation H0B0} (travaux au voisinage de pièces nues sous tension BT/HTA), prérequis indispensable pour intervenir en sécurité sur vos ouvrages. \\[0.5em]
\textbf{Corps — §2 :} Lors de mon stage à l'agence ENEO Douala-Bassa (service maintenance postes HTA/BT), j'ai participé activement au \textbf{diagnostic de pannes} sur cellules HTA, au \textbf{remplacement d'isolateurs} sur lignes aériennes 30 kV et aux manœuvres d'exploitation en nacelle. \textbf{J'ai ainsi diagnostiqué, isolé, réparé et validé le rétablissement du courant sur une ligne 30 kV, gravissant les pylônes avec assurance malgré le vent.} Ces interventions m'ont confronté au \textbf{travail en hauteur} et au respect scrupuleux des \textbf{consignes de sécurité} (cadenassage, balisage, PPE). Titulaire du \textbf{Permis B} et sportif (course à pied, endurance), je dispose de la condition physique requise pour les astreintes et les interventions en milieu difficile. \\[0.5em]
\textbf{Corps — §3 :} Rejoindre ENEO, c'est pour moi contribuer à une mission de \textbf{service public} essentielle : la continuité de la fourniture d'électricité. Je partage vos valeurs de performance, de sécurité et de proximité client. Ma lecture courante de l'\textbf{anglais technique} me permet d'exploiter la documentation constructeur internationale. \\[0.5em]
\textbf{Conclusion :} Disponible immédiatement pour un entretien et pour toute mission terrain, je vous prie d'agréer, Monsieur le Directeur, l'assurance de ma considération distinguée. \\[1em]
\textbf{Kevin Abanda} \\
\textit{Signature}
\end{quote}

\textbf{Partie B — Analyse langagière}
\begin{enumerate}
    \item \textbf{Figure de style :} \cle{Énumération} (suite de quatre verbes d'action au passé composé : « diagnostiqué, isolé, réparé, validé »). \textbf{Justification :} Cette série de verbes sans conjonction (asyndète) ou avec conjonction (polysyndète « et » avant le dernier) dresse la liste exhaustive et chronologique des étapes de l'intervention technique. Elle structure le récit professionnel et démontre la maîtrise du processus complet. On peut aussi y voir une \cle{gradation} logique (ordre opératoire : diagnostic $\rightarrow$ isolement $\rightarrow$ réparation $\rightarrow$ validation).
    \item \textbf{Nature et fonction de « gravissant » :} \textbf{Nature :} Participe présent (forme en \emph{-ant}) du verbe \emph{gravir}. \textbf{Fonction :} Complément circonstanciel de \textbf{moyen} / \textbf{manière} / \textbf{simultanéité} (réponse à la question « comment ? » ou « en faisant quoi ? »). \textbf{Valeur aspectuelle :} Il exprime le \textbf{processus en cours} (aspect imperfectif), la simultanéité de l'action de gravir par rapport à l'action principale (validé le rétablissement). Il met en scène l'effort physique concomitant à la réussite technique.
    \item \textbf{Transformation à la voix passive (sujet : \emph{le rétablissement du courant}) :} \\
    « Le rétablissement du courant sur une ligne 30 kV \textbf{a été diagnostiqué, isolé, réparé et validé} par mes soins, \textbf{tandis que je gravissais} les pylônes avec assurance malgré le vent. » \\
    \emph{Observation :} Le passage à la voix passive déplace le focus de l'agent (Kevin) vers le résultat (le rétablissement). Le participe présent « gravissant » (lié au sujet « je ») ne peut plus se rattacher au nouveau sujet inanimé (« le rétablissement »). Il faut donc le transformer en proposition subordonnée (« tandis que je gravissais ») ou en gérondif avec sujet explicite (« en gravissant... » - mais sujet différent). La perspective change : on met l'accent sur l'aboutissement technique plutôt que sur l'acteur.
    \item \textbf{Formule de politesse :} « Je vous prie d'\textbf{agréez}, Monsieur le Directeur, l'assurance de ma considération distinguée. » \\
    \textbf{Mode :} \cle{Impératif}. \textbf{Temps :} \cle{Présent}. \textbf{Forme :} 2\textsuperscript{e} personne du pluriel (\emph{vous agréez $\rightarrow$ agréez}). \textbf{Valeur :} \cle{Souhait} / \cle{Prière} conventionnelle (formule figée de politesse administrative). Ce n'est pas un ordre, mais une marque de déférence codifiée.
\end{enumerate}
\end{corrige}

\newpage

\coursec{Fiche bilan}

\begin{popfigure}
\begin{tikzpicture}[
    mindmap,
    grow cyclic,
    every node/.style={concept, execute at begin node=\hskip0pt},
    concept color=popblue,
    level 1/.append style={level distance=4.5cm, sibling angle=360/6},
    level 2/.append style={level distance=3cm, sibling angle=360/6},
    text=white,
    font=\small\bfseries
]
\node[concept, scale=1.1, align=center]{Lettre de\\ motivation}
    child[concept color=poporange] { node[concept, align=center]{Textes\\ fonctionnels} }
    child[concept color=popgreen] { node[concept, align=center]{Structure\\ externe} }
    child[concept color=poppurple] { node[concept, align=center]{Structure\\ interne\\ (Intro / Conclusion)} }
    child[concept color=poppink] { node[concept, align=center]{Corps de lettre\\ \& Outils langagiers} }
    child[concept color=popgold] { node[concept, align=center]{Communication\\ par l'image} }
    child[concept color=popdark] { node[concept, align=center]{Figures de style\\ \& Impératif} };
\end{tikzpicture}
\legende{Carte mentale de la leçon 21 : Produire une lettre de motivation}
\end{popfigure}

\begin{bilan}
\courssub{Définitions clés}
\begin{itemize}
    \item \textbf{Texte fonctionnel} : texte à visée utilitaire (administratif, professionnel) codifié par des normes strictes.
    \item \textbf{Lettre officielle} : correspondance formelle respectant une structure externe (mise en page) et interne (contenu) normalisées.
    \item \textbf{Lettre de motivation} : lettre d'accompagnement (candidature) qui persuade le destinataire de l'adéquation profil/poste.
\end{itemize}

\courssub{Structure externe (mise en page)}
\begin{center}
\begin{tabularx}{\linewidth}{|l|X|}\hline
\rowcolor{popblueL} \textbf{Zone} & \textbf{Contenu / Règles} \\ \hline
Expéditeur & Haut gauche : Nom, Prénom, Adresse, Tél., Email \\ \hline
Destinataire & Haut droite (sous expéditeur) : Nom/Qualité, Structure, Adresse \\ \hline
Lieu \& Date & Ligne séparée, alignée à droite (ex. : Yaoundé, le 15 octobre 2025) \\ \hline
Objet & Gras, concis, mentionne le poste/concours + référence \\ \hline
Formule d'appel & \emph{Monsieur le Directeur,} / \emph{Madame la Directrice,} (sans nom propre) \\ \hline
Corps & Introduction -- Développement (2--3 §) -- Conclusion \\ \hline
Formule de politesse & Adaptée au destinataire (ex. : \emph{Je vous prie d'agréer, Monsieur le Directeur, l'expression de ma considération distinguée.}) \\ \hline
Signature & Manuscritte (entre formule et nom tapé) \\ \hline
Pièces jointes & Mention \emph{Pièces jointes :} + liste (CV, diplômes, etc.) \\ \hline
\end{tabularx}
\end{center}

\courssub{Structure interne : Introduction \& Conclusion}
\begin{methode}[Introduction (3 mouvements)]
\begin{enumerate}
    \item \textbf{Accroche / Origine} : Source de l'offre (journal, site, affichage, relation).
    \item \textbf{Objet précis} : Poste visé, référence de l'annonce.
    \item \textbf{Annonce du plan} : \emph{« Mon parcours et mes motivations me permettent de... »}
\end{enumerate}
\end{methode}
\begin{methode}[Conclusion (2 mouvements)]
\begin{enumerate}
    \item \textbf{Disponibilité} : Entretien, informations complémentaires.
    \item \textbf{Formule de politesse complète} (voir tableau externe).
\end{enumerate}
\end{methode}

\courssub{Corps de la lettre : Stratégie \cle{Vous -- Moi -- Nous}}
\begin{itemize}
    \item \textbf{§ Vous (L'entreprise)} : Intérêt pour la structure, valeurs, projets connus (preuve de recherche).
    \item \textbf{§ Moi (Le candidat)} : Parcours (diplômes techniques, stages, projets), compétences \cle{dures} (techniques) et \cle{douces} (rigueur, travail en équipe), résultats chiffrés si possible.
    \item \textbf{§ Nous (La rencontre)} : Adéquation profil/poste, valeur ajoutée, motivation profonde.
\end{itemize}

\courssub{Outils langagiers obligatoires}
\begin{center}
\begin{tabularx}{\linewidth}{|l|X|X|}\hline
\rowcolor{popblueL} \textbf{Outil} & \textbf{Définition / Valeur} & \textbf{Exemple en candidature} \\ \hline
\textbf{Hyperbole} & Exagération pour insister & \emph{« Une expérience d'une richesse inestimable »} \\ \hline
\textbf{Métonymie} & Substitution (contenant/contenu, auteur/œuvre, lieu/institution) & \emph{« Intégrer vos ateliers » (pour l'entreprise) } \\ \hline
\textbf{Gradation} & Enchaînement termes croissants/décroissants & \emph{« Apprendre, maîtriser, innover »} \\ \hline
\textbf{Énumération} & Liste d'éléments (effet d'abondance ou exhaustivité) & \emph{« Rigueur, autonomie, esprit d'équipe »} \\ \hline
\textbf{Mode impératif} & Ordre, conseil, invitation (souvent à la 2e personne) & \emph{« Considérez ma candidature »} (formule de politesse), \emph{« Veuillez agréer... »} \\ \hline
\end{tabularx}
\end{center}

\courssub{Communication par l'image (CV / Portfolio joints)}
\begin{itemize}
    \item \textbf{Cohérence visuelle} : Même police, en-tête, code couleur que la lettre.
    \item \textbf{Lisibilité} : Espaces, puces, icônes sobres (technique).
    \item \textbf{Photo} : Professionnelle, fond neutre, cadre buste (norme camerounaise).
    \item \textbf{Portfolio technique} : Photos de réalisations (soudure, menuiserie, code, plans), légendées.
\end{itemize}
\end{bilan}

\begin{aretenir}[Checklist avant le Bac -- Lettre de motivation]
$\square$ \textbf{Format A4}, marges normales, police lisible (Times New Roman / Arial 12), interligne 1,15.
$\square$ \textbf{Structure externe complète} : Expéditeur, Destinataire, Lieu/Date, Objet (gras), Appel, Corps, Formule, Signature, PJ.
$\square$ \textbf{Introduction en 3 temps} : Source $\rightarrow$ Poste visé $\rightarrow$ Annonce du plan.
$\square$ \textbf{Corps structuré « Vous -- Moi -- Nous »} (3 paragraphes distincts).
$\square$ \textbf{Arguments « Moi »} : Diplômes techniques (Probatoire/Bac Tech), stages, compétences \textbf{chiffrées} ou \textbf{certifiées}.
$\square$ \textbf{Zéro faute d'orthographe/grammaire} (relecture à voix haute + correcteur).
$\square$ \textbf{Au moins 2 figures de style} (hyperbole, métonymie, gradation, énumération) utilisées à bon escient.
$\square$ \textbf{Impératif} correct dans la formule de politesse (\emph{Veuillez agréer...}) et éventuellement en appel à l'action.
$\square$ \textbf{Ton professionnel} : Vouvoiement, pas de familier, pas d'argot, vocabulaire technique précis.
$\square$ \textbf{Cohérence avec le CV} (dates, intitulés de poste, compétences identiques).
$\square$ \textbf{Pièces jointes listées} et présentes dans l'enveloppe / email.
$\square$ \textbf{Envoi respecté} : Délai, mode (dépôt physique, email PDF unique nommé \emph{Nom\_Prenom\_LettreMotivation.pdf}).
\end{aretenir}

\findecours

\end{document}
